\chapter{Numerical methods}
\label{vol02:ch16}

\epigraph{The purpose of computing is insight, not numbers.}{Richard
W.\ Hamming, \emph{Numerical Methods for Scientists and Engineers},
2nd ed., 1973, preface; the orientation paraphrased throughout this
chapter follows Higham
\cite[ch.~1]{text:higham-numerical}.}

\chapmeta{Half-life: foundational. Archetypes: failure (numerical),
uncertainty (conditioning and round-off as the formal handling of
input-error and finite-precision amplification). Exercise target: 35.
Project track: simulation with analysis (reproduce one published
numerical result, account for any drift, plus a 1500-word
reflection). Hazard class: none. Reader path default: standard, with
sections explicitly tagged. Named cases: Patriot missile system at
Dhahran (1991); Ariane 5 Flight 501 (1996); Pentium FDIV bug (1994);
Vancouver Stock Exchange index drift (1982-1983).}

\noindent
A computation that returns a number is not the same thing as a
computation that returns the right number. Volume II Chapters 5
through 15 set up the analytic apparatus on which engineers fit
models, integrate dynamics, solve linear systems, and find optima;
each of those chapters worked in exact arithmetic on the page. A
working engineer does not have exact arithmetic. The arithmetic the
engineer has is finite-precision, often IEEE 754 double, and the
gap between the page and the silicon is the subject of this chapter.

The discipline is to make the gap visible. Three quantities sit
behind every numerical answer the engineer reports: the truncation
error of the algorithm, the round-off error of the arithmetic, and
the conditioning of the problem. When the engineer can name all
three, the answer travels. When even one is hidden, the answer is
liable to be replaced silently by a different one. The chapter
develops the apparatus that names them, the algorithms whose
behaviour the engineer should know without lookup, and the named
cases in which an undocumented numerical decision destroyed an
engineering system.

\section{Floating-point arithmetic}\pathtag{core}

The arithmetic that runs in modern engineering software is, in the
overwhelming majority of cases, the IEEE 754 binary floating-point
arithmetic ratified in 1985 and revised in 2008
\cite[ch.~2]{text:higham-numerical}. The standard fixes the
representation, the rounding rules, the special values, and the
required operations. Engineers who do not know what the standard
guarantees, and what it explicitly does not, will sooner or later
report a number whose provenance they cannot defend.

\subsection{The representation}

A binary floating-point number $x$ is represented as

\[
x = (-1)^{s} \cdot (1.f) \cdot 2^{e},
\]

where $s \in \{0, 1\}$ is the sign bit, $f$ is a fraction in $[0, 1)$
stored as the bits of the \engterm{mantissa} (or significand), and
$e$ is a binary exponent stored with a bias. The leading $1.$ is
implicit for normal numbers; the bits actually stored are $s$, the
biased $e$, and the bits of $f$. The standard formats relevant to
engineering are summarised in Table~\ref{tab:ieee754}.

\begin{center}
\begin{tabular}{lcccc}
\hline
Format & Bits & Mantissa bits & Exponent bits & Approx.~range \\
\hline
binary16 (half) & 16 & 10 + 1 implicit & 5 & $6\times 10^{-5}$ to $6.5\times 10^{4}$ \\
binary32 (single) & 32 & 23 + 1 implicit & 8 & $1.2\times 10^{-38}$ to $3.4\times 10^{38}$ \\
binary64 (double) & 64 & 52 + 1 implicit & 11 & $2.2\times 10^{-308}$ to $1.8\times 10^{308}$ \\
binary128 (quad) & 128 & 112 + 1 implicit & 15 & $3.4\times 10^{-4932}$ to $1.2\times 10^{4932}$ \\
\hline
\end{tabular}
\end{center}
\label{tab:ieee754}

The mantissa-bit count fixes the precision; the exponent-bit count
fixes the range. Double precision is the working default in
scientific computing because its precision (about 15 to 17 decimal
digits) and its range cover most engineering quantities without
intermediate scaling. Single precision (about 7 decimal digits) is
the default in graphics and in some machine-learning workloads
because its memory and bandwidth costs are half of double's. Half
precision (about 3 to 4 decimal digits) appears in deep-learning
training and in some embedded systems, and is too narrow for most
classical numerical work.

\begin{figure}[ht]
\centering
\begin{tikzpicture}[x=0.16cm, y=0.55cm, font=\footnotesize]

% Top: bit-field layout of IEEE 754 binary64
\draw[thick] (0,0) rectangle (64,1);
% sign bit (1)
\draw[thick] (1,0) -- (1,1);
% exponent (11)
\draw[thick] (12,0) -- (12,1);
% Fill sign
\fill[engblue!20] (0,0) rectangle (1,1);
\node[anchor=center] at (0.5,0.5) {s};
% Fill exponent
\fill[engred!20] (1,0) rectangle (12,1);
\node[anchor=center] at (6.5,0.5) {biased exponent $e$ (11 bits)};
% Fill mantissa
\fill[engamber!25] (12,0) rectangle (64,1);
\node[anchor=center] at (38,0.5) {fraction $f$ (52 bits)};

% Bit indices
\node[anchor=south, font=\scriptsize] at (0.5, 1.05) {63};
\node[anchor=south, font=\scriptsize] at (6.5, 1.05) {62..52};
\node[anchor=south, font=\scriptsize] at (38, 1.05) {51..0};

% Labels for value formula below
\node[anchor=north] at (32, -0.3)
  {$x = (-1)^{s} \cdot (1.f)_{2} \cdot 2^{\,e - 1023}$ for normal numbers};

% Decoded example for value 0.1
\node[anchor=north, align=center] at (32, -1.6)
  {Example: $0.1_{10} = (-1)^{0} \cdot 1.\overline{1001100110011001100110011001100110011001100110011010}_{2} \cdot 2^{-4}$,\\
  the closing $\ldots10$ is the round-to-nearest tie-break; absolute error $\approx 5.55 \times 10^{-18}$.};

\end{tikzpicture}
\caption[IEEE 754 binary64 bit-field layout]{IEEE 754 binary64 bit-field
layout. The sign bit, the 11-bit biased exponent, and the 52-bit fraction
together encode a normal number with an implicit leading 1. The decimal
constant $0.1$ has no exact binary representation; the round-to-nearest
representation differs from the true value by approximately
$5.55 \times 10^{-18}$, the source of the per-conversion error that
accumulated to lethal scale in the Patriot Dhahran case (section 16.8.1).
Source: authors' diagram from IEEE Std 754-2019
\cite{std:ieee754-2019}.}
\label{fig:vol02:ch16:binary64-layout}
\end{figure}


Figure~\ref{fig:vol02:ch16:binary64-layout} shows the bit-field
layout of binary64 and the decoded representation of the decimal
$0.1$. The infinite repeating binary fraction terminates at the
mantissa's 52nd bit by the round-to-nearest-ties-to-even rule. The
chapter's Failure section (section 16.8.1) returns to this constant
in deployed software.

\begin{estimation}
\textbf{Question.} A measurement chain produces a sequence of
$N = 10^{6}$ readings, each rounded to double precision and
summed in the order produced. Estimate the round-off error in the
final sum, before computing it.

\smallskip
\textit{Estimate, before reading on.}

\smallskip
Each addition is rounded with relative error at most the unit
round-off $u \approx 1.1 \times 10^{-16}$ for double precision. The
naive sequential sum accumulates error in roughly the worst-case
form $N u \cdot \abs*{S}$, where $S$ is the sum's magnitude, giving
a relative error bound near $10^{-10}$. The bound is loose; in
practice, with random-sign rounding errors, the typical accumulated
error grows as $\sqrt{N} u$ and the relative error is closer to
$10^{-13}$. Either bound is far below the measurement noise of any
realistic engineering chain. The lesson is the order of magnitude:
naive summation of a million double-precision numbers is safe; a
billion is borderline; a trillion calls for compensated summation
(Kahan, pairwise) or for an extended precision.
\end{estimation}

\subsection{Rounding}

The result of a real-arithmetic operation between two
floating-point numbers is, in general, not representable. The
standard requires the operation to round to the nearest
representable number, with ties broken toward the value whose last
mantissa bit is zero (\engterm{round-to-nearest, ties to even}).
This is the default rounding mode and the only mode the engineer
should normally rely on. Three further modes (round toward zero,
toward $+\infty$, toward $-\infty$) exist for interval arithmetic
and for specific algorithms; their use is documented per call site.

The bound on the rounding error of a single basic operation $\circ
\in \{+, -, \times, /\}$ is

\[
\operatorname{fl}(a \circ b) = (a \circ b)(1 + \delta), \qquad \abs{\delta} \leq u,
\]

where $\operatorname{fl}(\cdot)$ denotes the floating-point result and $u$ is the
\engterm{unit round-off} of the format \cite[\S2.2]{text:higham-numerical}.
For binary64, $u = 2^{-53} \approx 1.11 \times 10^{-16}$; for
binary32, $u = 2^{-24} \approx 5.96 \times 10^{-8}$. The square root
operation is also correctly rounded under IEEE 754; transcendental
functions (sin, cos, exp, log) are not, in general, required to be
correctly rounded, and an engineer who depends on the last bit of a
transcendental's value is depending on the specific math library
shipped with the platform.

\subsection{Machine epsilon}

The \engterm{machine epsilon} $\epsilon_{m}$ is the gap between
$1$ and the next representable floating-point number. For binary64,
$\epsilon_{m} = 2^{-52} \approx 2.22 \times 10^{-16}$; for binary32,
$\epsilon_{m} = 2^{-23} \approx 1.19 \times 10^{-7}$. The machine
epsilon is twice the unit round-off in the round-to-nearest scheme,
$\epsilon_{m} = 2 u$, because $u$ is the worst-case relative
rounding error and $\epsilon_{m}$ is the spacing between adjacent
representables near $1$. Both quantities appear in algorithm
analysis; the engineer should know both, and should know which one
a given source uses without ambiguity.

A short computation fixes the two constants in the reader's hands.
The next representable double above $1.0$ is $1 + 2^{-52}$, because
the mantissa of $1.0 = 1.\underbrace{00\cdots0}_{52}\times 2^{0}$
increments in its last bit. The gap is $2^{-52} = \epsilon_{m}$. The
worst-case relative error of rounding a real number in $[1, 2)$ to
the nearest double is half that gap, $2^{-53} = u$, attained when the
real number sits exactly midway between two representables. The
relationship $\epsilon_{m} = 2u$ is then mechanical, and the reader
who is handed a formula in $u$ and a colleague's formula in
$\epsilon_{m}$ can reconcile the two without guessing which the
author meant.

\begin{example}
\textbf{Spacing at a large exponent.} The spacing between adjacent
doubles is not constant; it scales with the binade. Near $x = 2^{e}$
the spacing is $\epsilon_{m} \cdot 2^{e} = 2^{e - 52}$. At
$x = 2^{52} \approx 4.5 \times 10^{15}$ the spacing is exactly $1$:
every integer up to $2^{53}$ is representable, but no half-integer
above $2^{52}$ is. At $x = 2^{53}$ the spacing is $2$, and odd
integers above $2^{53}$ are not representable. An accumulator that
counts events past four quadrillion in a double silently stops
incrementing when the increment falls below the local spacing. The
working lesson: an integer counter that may exceed $2^{53}$ belongs
in a 64-bit integer, not a double. The same arithmetic explains why
$10^{16} + 1 = 10^{16}$ in double precision, a result that surprises
the reader once and never again.
\end{example}

\begin{mastery}
\textbf{Why subnormals exist.} Without subnormal numbers the smallest
positive normal double is $2^{e_{\min}} = 2^{-1022}$, and the next
representable below it is $0$. The gap from the smallest normal to
zero would then be $2^{-1022}$, while the gap between the smallest
normal and its successor is only $2^{-1022} \cdot \epsilon_{m} =
2^{-1074}$, a jump of $2^{52}$ in the spacing as the value crosses
into zero. Worse, the identity $x - y = 0 \iff x = y$ would fail:
two distinct normals closer together than $2^{-1022}$ would subtract
to a nonzero real that rounds to zero, so $x - y = 0$ with $x \neq
y$. Subnormal numbers fill the gap $[0, 2^{-1022})$ with values of
the form $0.f \times 2^{-1022}$, spaced uniformly at $2^{-1074}$,
restoring the guarantee that $\operatorname{fl}(x - y) = 0$ implies
$x = y$ (\engterm{gradual underflow}). The cost is that subnormals
carry fewer significant bits as they approach zero, and on some
hardware a subnormal operand triggers a microcode trap that runs an
order of magnitude slower than a normal operation; performance-tuned
numerical kernels sometimes enable a flush-to-zero mode that trades
the gradual-underflow guarantee for uniform speed, a trade the
engineer documents at the call site. Kahan, the principal architect
of IEEE 754, argued the gradual-underflow guarantee on exactly the
$x - y = 0 \iff x = y$ grounds \cite[ch.~2]{text:overton-ieee754}.
\end{mastery}

\subsection{Special values}

IEEE 754 reserves several bit patterns for non-finite results.

\begin{itemize}[leftmargin=*]
  \item Signed zero: $+0$ and $-0$ are distinct representations and
    compare equal under $==$, but distinguish themselves under
    $1/x$, which yields $+\infty$ for $+0$ and $-\infty$ for $-0$.
  \item Signed infinity: $\pm\infty$ arise from overflow and from
    division of a finite nonzero by zero. Subsequent arithmetic
    propagates infinity deterministically: $\infty + 1 = \infty$,
    $\infty / \infty = \mathrm{NaN}$, $\infty \times 0 = \mathrm{NaN}$.
  \item Not-a-Number (NaN): the result of $0/0$, $\infty - \infty$,
    $\sqrt{-1}$, and any operation with at least one NaN operand.
    NaN compares unequal to every value including itself; the
    \texttt{NaN != NaN} test is the canonical NaN check.
  \item Subnormal (denormal) numbers: representations with the
    minimum biased exponent and a zero leading bit, used to fill the
    gap between zero and the smallest normal number. Subnormals
    extend the underflow region with gradual loss of precision; on
    some hardware they are slow, and a flush-to-zero mode replaces
    them with zero.
\end{itemize}

The engineer's discipline is to know that these values are
contagious, that they are sometimes legitimate (an unattainable
optimum returns $+\infty$ as a flag), and that they are sometimes
the silent shadow of an unhandled corner. A computation that
returns NaN is not a failed computation in the sense of a crash; it
is a successful computation of a non-result, and the engineer who
treats it as a number is in trouble. Section 16.8 makes this point
concrete.

\subsection{What the standard does not guarantee}

IEEE 754 fixes the basic operations and the formats. It does not
fix the order in which a compiler evaluates a sum, the precision
of intermediate values held in extended-precision registers, the
semantics of fused multiply-add (FMA) versus separate multiply and
add, or the bit-for-bit reproducibility of a parallel reduction
across thread counts. Section 16.7 returns to these. A program
that produces the same floating-point output on the same input on
the same hardware on the same day is a program. A program that
produces the same output on the same input on different hardware,
different optimisation levels, or different thread counts is a
program plus a discipline.

\begin{mastery}
\textbf{Error-free transformations and the fused multiply-add.} The
single-rounding bound $\operatorname{fl}(a \circ b) = (a \circ b)(1 +
\delta)$ hides a sharper fact the engineer can exploit. For addition and
subtraction the rounding error is itself a floating-point number, and it
can be computed exactly with no extra precision. The \engterm{TwoSum}
algorithm, given $a$ and $b$, returns $s = \operatorname{fl}(a + b)$ and
$e$ such that $a + b = s + e$ exactly, in six floating-point operations:
$s = \operatorname{fl}(a+b)$; $b' = \operatorname{fl}(s - a)$;
$e = \operatorname{fl}((a - (s - b')) + (b - b'))$. The recovered $e$ is
the rounding error the addition discarded, and Kahan summation is the
special case that folds $e$ back into the running sum. For multiplication
the analogous \engterm{TwoProduct} algorithm recovers the exact product
error, and it is here that the fused multiply-add earns its place: a
hardware FMA computes $\operatorname{fl}(a \cdot b + c)$ with a single
rounding, so $e = \operatorname{fl}(\text{fma}(a, b, -\operatorname{fl}(a
b)))$ delivers the multiplication error in one instruction rather than the
seventeen operations the FMA-free Dekker split requires. These
\engterm{error-free transformations} are the foundation of
compensated and doubly-compensated algorithms, which deliver an accuracy
as if the working precision were doubled while running entirely in the
native format \cite[ch.~4]{text:higham-numerical}. The same FMA that
improves accuracy here is a reproducibility hazard elsewhere: whether the
compiler contracts \texttt{a*b + c} into an FMA is a target and
optimisation-level decision, and the contracted and uncontracted forms
differ at the last bit. The engineer who wants the accuracy benefit
requests the FMA explicitly (the C \texttt{fma} function, the
\texttt{std::fma} call); the engineer who wants bit-reproducibility
across targets disables contraction (\texttt{-ffp-contract=off}). Section
16.7 returns to the second face of this trade.
\end{mastery}

\begin{mastery}
\textbf{Double rounding.} A value computed in a wider format and then
stored to a narrower one is rounded twice, and the two-step rounding can
land on a different representable than a single direct rounding to the
narrow format would. The historical instance is the x87 floating-point
unit, which evaluated expressions in an 80-bit extended format and rounded
to 64-bit double on store. Consider a real number that sits just above the
midpoint between two doubles but, in the 80-bit intermediate, rounds
down to exactly the midpoint; the subsequent round-to-even of that
midpoint then goes the other way from a direct round-to-nearest of the
original. The discrepancy is at most one unit in the last place, and it
is rare, but it breaks the clean single-rounding model that every error
bound in this chapter assumes. Double rounding is why bit-reproducible
numerics on the x86 architecture moved from the x87 unit to the SSE2
registers, which compute directly in the target precision, and why a
build that still routes through x87 (a 32-bit legacy target, certain
embedded toolchains) can disagree with an SSE2 build at the last bit on
the same source \cite[ch.~3]{text:overton-ieee754}. The durable lesson
outlives the specific hardware: when an intermediate is held in a wider
format than the destination, the engineer states the intermediate
precision as part of the computation's specification, because the result
depends on it.
\end{mastery}

\section{Error: truncation, round-off, conditioning}\pathtag{core}

Three error sources sit behind every numerical answer. The
discipline is to name each one, to bound it, and to choose the
algorithm whose dominant error matches the engineer's tolerance.

\subsection{Truncation error}

\engterm{Truncation error} (also called \engterm{discretisation
error}, or \engterm{algorithm error}) is the error introduced by
replacing the exact mathematical operation with a finite
approximation. The forward-difference approximation to the
derivative,

\[
f'(x) \approx \frac{f(x+h) - f(x)}{h},
\]

has truncation error $\frac{h}{2} f''(\xi)$ for some $\xi$ near
$x$, by Taylor's theorem. The error vanishes as $h \to 0$. A
five-term truncation of a Taylor series has truncation error of
order $h^{5}$ in the discarded terms. A trapezoidal-rule
quadrature on $n$ intervals has truncation error of order $h^{2}$.

Truncation error is independent of the arithmetic. It is a
property of the algorithm relative to the underlying mathematical
operation. It scales predictably with the discretisation parameter
($h$, $\Delta t$, the number of terms retained), and the engineer
controls it by choosing the parameter.

\subsection{Round-off error}

\engterm{Round-off error} is the error introduced by the
finite-precision arithmetic of the previous section. Each basic
operation contributes a relative error of size at most $u$, and
the accumulated round-off depends on the algorithm and on the data.
For a sum of $n$ floating-point numbers, the worst-case
accumulated round-off is bounded by $n u$, while the typical
behaviour for random-sign errors is closer to $\sqrt{n} u$
\cite[\S4.2]{text:higham-numerical}. For matrix-vector multiplication
in $\R^{n}$, the worst-case relative round-off is bounded by $n u$
times a small constant.

Round-off error is independent of the algorithm's truncation. It
is a property of the arithmetic relative to the algorithm as
written. It scales predictably with the format (single, double,
quad), and the engineer controls it by choosing the format and by
preferring numerically stable variants of the algorithm (Kahan
summation \cite{text:v2ch16-kahan-1965}, Householder QR over
Gram-Schmidt, partial pivoting in LU). The cancellation pattern in
which two nearly equal floating-point numbers are subtracted, and
the leading identical digits annihilate, is the canonical worst
case of round-off amplification; it is sometimes called
\engterm{catastrophic cancellation} because the algorithm choice
turns an otherwise well-behaved problem into an unrecoverable
loss of digits.

\begin{figure}[ht]
\centering
\begin{tikzpicture}[x=0.32cm, y=0.32cm, font=\footnotesize]

% Axis
\draw[->, thick] (0,0) -- (38,0) node[right] {$\log_{10}(1/x)$};
\draw[->, thick] (0,0) -- (0,18) node[above, align=center] {decimal digits of\\correct result};

% Tick marks for x = 10^{-k}, k=0..16
\foreach \k in {0,2,4,6,8,10,12,14,16} {
  \pgfmathsetmacro{\xp}{2*\k}
  \draw (\xp,0) -- (\xp,-0.3);
  \node[anchor=north, font=\scriptsize] at (\xp,-0.3) {$10^{-\k}$};
}
% y ticks
\foreach \d in {0,4,8,12,16} {
  \draw (0,\d) -- (-0.3,\d);
  \node[anchor=east, font=\scriptsize] at (-0.3,\d) {\d};
}

% Reference: ideal 16 digits constant line
\draw[thick, engblue, dashed] (0,16) -- (32,16) node[right, font=\scriptsize, engblue] {Taylor $x + x^{3}/6$};

% Catastrophic cancellation: naive sin(x) - x + x^3/6 in double
% Digits lost ~ 2k for x = 10^{-k}; clamp at 0
% Sample points (k, digits) approximately: (0,16) (2,12) (4,8) (6,4) (7,2) (8,0)
\draw[thick, engred] plot[smooth] coordinates {
  (0,16) (2,14) (4,10) (6,6) (8,2) (10,0) (16,0) (32,0)
};
\node[engred, font=\scriptsize, anchor=west] at (10.5,0.8)
  {naive double-precision evaluation};

% Annotated region of cancellation
\draw[thick, engamber, fill=engamber!20, opacity=0.4]
  (10,0) rectangle (32,16);
\node[anchor=center, font=\scriptsize] at (21,10)
  {cancellation region: $\sin x \approx x$ to many digits;};
\node[anchor=center, font=\scriptsize] at (21,8.5)
  {the difference $\sin x - x$ in double};
\node[anchor=center, font=\scriptsize] at (21,7)
  {underflows the significand of $\sin x$.};

\end{tikzpicture}
\caption[Catastrophic cancellation in $\sin(x) - x + x^{3}/6$]{Catastrophic
cancellation in the naive evaluation of $\sin(x) - x + x^{3}/6$ at small
$x$ in IEEE 754 binary64. The Taylor-rewritten form preserves all 16
digits of precision; the naive form loses approximately two digits per
decade of $x$ because the subtraction $\sin(x) - x$ cancels the leading
identical digits of two nearly equal numbers. By $x \approx 10^{-10}$
the naive result has no correct digits. The remedy is rewriting before
evaluating: the cancellation lives in the algorithm's choice, not in the
arithmetic. Source: authors' computation; analytic prediction
$\sim 2 \log_{10}(1/x)$ digits lost.}
\label{fig:vol02:ch16:catastrophic-cancellation}
\end{figure}


Figure~\ref{fig:vol02:ch16:catastrophic-cancellation} reads the
mechanism on the example $\sin(x) - x + x^{3}/6$. The naive
evaluation loses approximately two decimal digits per decade of
$x$; the Taylor-rewritten form $x^{5}/120 - x^{7}/5040$ preserves
all 16. The estimation block above runs the same diagnostic on
$1 - \cos(x)$. Across this chapter the same recipe applies. The
algorithm's correctness up to algebraic identity does not imply
its numerical safety; the engineer rewrites before evaluating
whenever the input drives two summands into near-equality. Goldberg
\cite{text:goldberg1991-floating-point} catalogues additional
patterns (Heron's formula at near-collinear triangles, the
quadratic formula at small discriminant, the recursive variance
computation) that share the structure and the remedy.

\begin{example}
\textbf{Kahan compensated summation.} Summing $n$ numbers in the
naive sequential order accumulates a round-off error bounded by
$(n - 1) u$ times the sum of magnitudes. \engterm{Kahan summation}
carries a running compensation term that recaptures the low-order
bits lost at each addition \cite{text:v2ch16-kahan-1965}: maintain a
running sum $s$ and a compensation $c$, and at each step compute
$y = x_{i} - c$, $t = s + y$, $c = (t - s) - y$, $s = t$. The
quantity $(t - s) - y$ is, in exact arithmetic, zero; in
floating-point it captures the part of $y$ that did not fit into
$s$, and subtracting it from the next addend restores it. The
worst-case error of Kahan summation is bounded by $2u$ plus a term
of order $n u^{2}$, independent of $n$ to first order, against the
naive $(n - 1) u$. The contrast is sharp on a stress test. Summing
$1$ followed by $10^{8}$ copies of $10^{-8}$ in binary64: the exact
sum is $2$, the naive left-to-right sum returns about $1.9999998$
(the small addends fall below the spacing of $s$ once $s \approx 1$
and are partly swallowed), and Kahan summation returns $2.0000000$
to all visible digits. The
\texttt{code/kahan\_summation.py} script reproduces both sums and
prints the per-step compensation. Pairwise (cascade) summation,
which adds in a balanced binary tree, achieves an $O(u \log n)$
bound at no compensation-arithmetic cost and is what NumPy's
\texttt{sum} uses by default; Kahan is the choice when the engineer
needs the tighter $O(u)$ bound and can pay the four operations per
element.
\end{example}

\subsection{Conditioning}

\engterm{Conditioning} is a property of the problem, not of any
algorithm. The condition number measures how much a small relative
perturbation of the input amplifies into a relative perturbation of
the output. For a differentiable scalar function $f: \R \to \R$, the
relative condition number is

\[
\kappa_{f}(x) = \abs*{\frac{x f'(x)}{f(x)}},
\]

defined where $f(x) \neq 0$. For a linear system $\mat{A}\vect{x} =
\vect{b}$, the condition number is the matrix-norm quantity
$\kappa(\mat{A}) = \norm{\mat{A}} \norm{\mat{A}^{-1}}$ introduced
in section 9.7. For a root-finding problem $f(x) = 0$ at a simple
root $x^{*}$, the condition number with respect to perturbations of
$f$ is $1/\abs*{f'(x^{*})}$ in absolute terms.

The dimensional reading is essential. A condition number of
$10^{6}$ does not mean the algorithm is bad; it means the problem
amplifies input perturbations by a factor of $10^{6}$. Even an
algorithm that introduces no round-off cannot recover digits the
problem has lost. The engineer who reports a result on an
ill-conditioned problem reports the conditioning, the precision,
and the rule-of-thumb digit count, and refuses to over-claim.

\begin{figure}[ht]
\centering
\begin{tikzpicture}[x=0.6cm, y=0.6cm, font=\footnotesize]

% Two side-by-side panels: well-conditioned vs ill-conditioned 2x2 system
% Each panel shows two lines (rows of A x = b) and their intersection;
% small perturbation of b becomes small/large displacement of solution.

% --- Left panel: well-conditioned
\begin{scope}[shift={(0,0)}]
  \node[anchor=south] at (3,4.4) {\textbf{Well-conditioned} $\kappa \approx 1$};
  \draw[->, thick] (-0.2,0) -- (6,0) node[right] {$x_{1}$};
  \draw[->, thick] (0,-0.2) -- (0,4) node[above] {$x_{2}$};

  % Line 1: x1 + x2 = 4
  \draw[thick, engblue] (0,4) -- (4,0);
  \node[engblue, anchor=south west, font=\scriptsize] at (4.0,0.0) {$L_{1}$};
  % Line 2: x1 - x2 = 0, slope 1
  \draw[thick, engred] (0,0) -- (3.5,3.5);
  \node[engred, anchor=south, font=\scriptsize] at (3.6,3.5) {$L_{2}$};
  % Perturbed line 2: shifted by 0.2 down
  \draw[thick, engred, dashed] (0.2,0) -- (3.7,3.5);

  % Solutions: intersection at (2,2) and ~ (2.1, 1.9)
  \fill[black] (2,2) circle (2pt);
  \fill[black] (2.1,1.9) circle (2pt);
  \draw[thick, ->, gray] (2,2) -- (2.1,1.9);
  \node[anchor=west, font=\scriptsize] at (2.2,1.85)
    {small $\Delta x$};
\end{scope}

% --- Right panel: ill-conditioned (near-parallel lines)
\begin{scope}[shift={(8,0)}]
  \node[anchor=south] at (3,4.4) {\textbf{Ill-conditioned} $\kappa \gg 1$};
  \draw[->, thick] (-0.2,0) -- (6,0) node[right] {$x_{1}$};
  \draw[->, thick] (0,-0.2) -- (0,4) node[above] {$x_{2}$};

  % Line 1: gentle slope
  \draw[thick, engblue] (0,0.5) -- (5.5,3.25);
  \node[engblue, anchor=south, font=\scriptsize] at (5.5,3.25) {$L_{1}$};
  % Line 2: slightly different slope (near-parallel)
  \draw[thick, engred] (0,0.2) -- (5.5,3.5);
  \node[engred, anchor=west, font=\scriptsize] at (5.5,3.5) {$L_{2}$};
  % Perturbed line 2: shift by 0.1
  \draw[thick, engred, dashed] (0,0.4) -- (5.5,3.7);

  % Intersections: very far apart
  \fill[black] (2.4,1.7) circle (2pt);
  \fill[black] (4.6,2.8) circle (2pt);
  \draw[thick, ->, gray] (2.4,1.7) -- (4.55,2.78);
  \node[anchor=south west, font=\scriptsize] at (3.2,2.2)
    {large $\Delta x$};
\end{scope}

\end{tikzpicture}
\caption[Geometric reading of the condition number]{Geometric reading of
the condition number for a $2 \times 2$ linear system. The two rows of
$\mat{A}$ are two lines in the $(x_{1}, x_{2})$ plane; the solution
$\vect{x}$ is their intersection. When the lines are nearly orthogonal
the system is well-conditioned: a small perturbation of $\vect{b}$
(shifting one of the lines by a small amount) moves the intersection by
a comparable small amount. When the lines are nearly parallel the
system is ill-conditioned: the same shift moves the intersection by a
large amount, amplified by $\kappa(\mat{A})$. The Hilbert matrix
exercise of section 16.4 makes this concrete in $n$ dimensions.
Source: authors' diagram; geometric framing follows
\cite[ch.~12]{text:trefethen-bau1997}.}
\label{fig:vol02:ch16:condition-number}
\end{figure}


Figure~\ref{fig:vol02:ch16:condition-number} reads the condition
number geometrically. The rows of $\mat{A}$ are lines (in two
dimensions); the solution is their intersection. When the lines are
near-orthogonal the intersection is stable under small perturbations
of $\vect{b}$; when the lines are near-parallel a small perturbation
moves the intersection by a large amount, amplified by
$\kappa(\mat{A})$. The Hilbert-matrix simulation exercise of section
16.4 reproduces this picture at higher dimension and at numerically
catastrophic conditioning ($\kappa \sim 10^{12}$ at $n = 10$).

\begin{archetype}[Uncertainty]
The condition number is the bridge between this chapter and the
uncertainty apparatus of Volume I Chapter 4. There, the propagation
of measurement uncertainty through a function $y = f(x)$ multiplied
the input standard uncertainty by the sensitivity $\abs{f'(x)}$.
Here, the relative condition number $\kappa_{f}(x) = \abs{x f'(x) /
f(x)}$ multiplies the input relative perturbation by the same
sensitivity, recast in relative terms. Finite-precision round-off is
one source of input perturbation; measurement error is another;
model error is a third. The condition number does not care which
source supplied the perturbation. It states how much the problem
amplifies whatever uncertainty the input carries, and it sets a
floor on the output uncertainty that no algorithm, however stable,
can beat. The engineer who has internalised the uncertainty budget
of Volume I reads conditioning as the same accounting applied to the
arithmetic: the input has a relative uncertainty (from measurement,
from rounding, from the model), the problem has a condition number,
and the product is the irreducible relative uncertainty of the
output. Reporting more digits than that product allows is the
numerical form of claiming an error bar the data do not support.
\end{archetype}

\begin{estimation}
\textbf{Question.} An engineer evaluates
$g(x) = 1 - \cos(x)$ at $x = 10^{-7}$ in double precision. How many
correct decimal digits does the engineer expect, and why? Before
calculating, estimate the worst error source.

\smallskip
\textit{Estimate, before reading on.}

\smallskip
For $x = 10^{-7}$, the exact value is $g(x) \approx x^{2}/2 =
5 \times 10^{-15}$. The naive evaluation computes $\cos(x)$ to a
relative error $\sim u$, returning a value close to $1$ that
differs from the exact $\cos(x)$ in roughly the 15th decimal
digit. The subtraction $1 - \cos(x)$ then cancels the 14 identical
leading digits of $1$ and $\cos(x)$, leaving an absolute error
near $u$ in a result whose true magnitude is $5 \times 10^{-15}$.
The relative error of the naive evaluation is approximately
$u / (x^{2}/2) \approx 4$; the answer has \emph{no} correct
digits. The condition number of $g$ at this $x$ is
$\kappa_{g}(x) = \abs*{x g'(x)/g(x)} = \abs{x \sin x / (1 - \cos
x)} \approx 2$, so the problem itself is well conditioned. The
algorithm is the buyer: the cancellation error is unforced. The
remedy is to rewrite, using either the trig identity $1 - \cos x
= 2 \sin^{2}(x/2)$ (which evaluates as $2 (\sin(x/2))^{2}$ with no
cancellation) or the Taylor truncation $x^{2}/2 - x^{4}/24$. Either
form recovers all 16 digits. The estimation drills the section's
core point: the three error sources are independent. The arithmetic
is correctly rounded; the problem is well conditioned; the algorithm
loses everything. Section 16.8 returns to this pattern at deployed
scale.
\end{estimation}

\subsection{Forward and backward error}

Two error notions are formal in numerical analysis. The
\engterm{forward error} is the difference between the computed
output $\hat{y}$ and the exact output $y$ for the given input.
The \engterm{backward error} is the smallest perturbation
$\Delta x$ of the input such that the computed output equals the
exact output for the perturbed input,

\[
\hat{y} = f(x + \Delta x).
\]

An algorithm is \engterm{backward stable} if its backward error
is small relative to $u$ for every input. Backward stability is
the strongest stability property an engineer can usually expect
from a general-purpose numerical algorithm; LU with partial
pivoting, QR by Householder, the symmetric-tridiagonal eigenvalue
QR algorithm, and the SVD by the Golub-Reinsch routine are all
backward stable
\cite[ch.~7]{text:trefethen-bau1997}.

The relationship between forward and backward error is

\[
\text{forward error} \lesssim \kappa(\text{problem}) \cdot
\text{backward error}.
\]

A backward-stable algorithm on an ill-conditioned problem still
loses digits; the digits the engineer loses are bought by the
problem, not by the algorithm. The discipline is to identify
which buyer is responsible.

\begin{example}
\textbf{The same digits, two different buyers.} Consider evaluating
$y = f(x) = \sqrt{x + 1} - \sqrt{x}$ at $x = 10^{12}$ in double
precision. The naive evaluation subtracts two numbers that agree to
about twelve digits, $\sqrt{x + 1} \approx \sqrt{x} \approx 10^{6}$,
and the cancellation leaves an absolute error near $u \cdot 10^{6}
\approx 10^{-10}$ in a result whose true magnitude is $1/(\sqrt{x+1}
+ \sqrt{x}) \approx 5 \times 10^{-7}$. The relative error of the
naive computation is near $2 \times 10^{-4}$, about four correct
digits. Is the loss the problem's fault or the algorithm's? The
condition number is $\kappa_{f}(x) = \abs{x f'(x)/f(x)}$. With
$f'(x) = \tfrac{1}{2}(x+1)^{-1/2} - \tfrac{1}{2}x^{-1/2}$ the
algebra gives $\kappa_{f}(x) = \tfrac{1}{2} \cdot
\frac{\sqrt{x+1} + \sqrt{x}}{\sqrt{x+1} - \sqrt{x}} \cdot
\frac{x}{\sqrt{x(x+1)}} \approx \tfrac{1}{2}$ for large $x$. The
problem is well conditioned. The algorithm is the buyer. Rewriting
$f(x) = 1/(\sqrt{x + 1} + \sqrt{x})$ adds two like-signed numbers
and divides, with no cancellation, and recovers all sixteen digits.
The diagnostic is the same one the estimation block ran on $1 -
\cos x$: compute $\kappa$, and if $\kappa$ is small while the answer
is wrong, rewrite the algorithm rather than reaching for higher
precision.
\end{example}

\begin{example}
\textbf{When the problem is the buyer.} Evaluate the polynomial
$p(x) = (x - 1)^{6}$ near $x = 1$ in its expanded monomial form
$x^{6} - 6x^{5} + 15x^{4} - 20x^{3} + 15x^{2} - 6x + 1$. At
$x = 1.01$ the exact value is $10^{-12}$, but the expanded
evaluation sums seven terms each of magnitude near $1$ to $20$ whose
true sum is $10^{-12}$; the result loses about twelve digits to
cancellation among the partial sums. Here higher precision genuinely
helps, because the conditioning of polynomial evaluation in the
monomial basis near a high-multiplicity root is itself large: the
condition number of evaluating $\sum a_k x^k$ is $\sum \abs{a_k}
\abs{x}^k / \abs{p(x)}$, which at $x = 1.01$ is roughly $77 /
10^{-12} \approx 8 \times 10^{13}$. The problem, as posed in the
monomial basis, is ill conditioned; the algorithm cannot recover
what the representation has lost. The remedy is to change the
problem, not the precision: evaluate $(x - 1)^{6}$ in factored form,
which is perfectly conditioned. The Wilkinson polynomial of the
following box is the extreme case of this pathology.
\end{example}

\begin{mastery}
\textbf{Wilkinson's polynomial.} Wilkinson's perturbation study of
the polynomial $w(x) = \prod_{k=1}^{20}(x - k)$, with roots at the
integers $1$ through $20$, is the historical demonstration that root
conditioning can be catastrophic even when every root is real,
simple, and well separated \cite[ch.~2]{text:v2ch16-wilkinson-1965}.
The coefficient of $x^{19}$ in the expanded form is $-210$.
Perturbing that single coefficient by $2^{-23} \approx 10^{-7}$ (one
part in $10^{9}$ of its magnitude) splits several of the larger
roots into complex-conjugate pairs displaced by order unity: the
root near $x = 16$ moves by about $2.8$. The condition number of the
root $x_j$ with respect to the coefficient $a_k$ is
$\abs{a_k x_j^{k}} / \abs{x_j \, w'(x_j)}$, and for the clustered
larger roots of $w$ this reaches $10^{10}$ and beyond. Wilkinson
recorded that the discovery, that a polynomial with integer roots
this benign-looking could be so violently ill conditioned in its
coefficients, was the most traumatic experience of his career as a
numerical analyst. The engineering lesson is durable past any
language or library: the map from polynomial coefficients to roots
is, for clustered or high-degree polynomials, an ill-conditioned
problem, and a root-finder reported without its coefficient-to-root
condition number is a result whose error the engineer cannot bound.
The practical defence is to avoid forming the monomial coefficients
at all when the application allows it (work with the factored form,
or with an eigenvalue formulation of the companion matrix, whose
conditioning is governed by the matrix rather than the
coefficients).
\end{mastery}

\subsection{A backward-error analysis carried to the bound}

The forward-error inequality $\text{forward error} \lesssim \kappa \cdot
\text{backward error}$ is a slogan until the engineer has run a backward
analysis end to end on a concrete computation. Horner's rule for
evaluating a polynomial is the cleanest vehicle, because its backward
error has a closed form and the analysis exposes every step where round-off
enters.

Evaluate $p(x) = a_{n} x^{n} + \cdots + a_{1} x + a_{0}$ by Horner's
nested form,

\[
p(x) = (\cdots((a_{n} x + a_{n-1}) x + a_{n-2}) x \cdots) x + a_{0},
\]

which costs $n$ multiplications and $n$ additions. In floating-point each
operation carries a relative error bounded by $u$. Writing the computed
partial result after absorbing the multiplication by $x$ at stage $k$ as
$\hat{b}_{k}$, the recurrence is $\hat{b}_{k} = \operatorname{fl}
(\hat{b}_{k+1} x + a_{k}) = (\hat{b}_{k+1} x (1 + \mu_{k}) + a_{k})(1 +
\alpha_{k})$ with $\abs{\mu_{k}}, \abs{\alpha_{k}} \leq u$, where $\mu_{k}$
is the rounding of the multiply and $\alpha_{k}$ the rounding of the add.
Carrying the products of $(1 + \mu_{k})$ and $(1 + \alpha_{k})$ factors
through the recurrence and collecting them with the standard
$\prod(1 + \delta_{i}) = 1 + \theta$ device, where $\abs{\theta_{m}} \leq
m u / (1 - m u) \equiv \gamma_{m}$ for $m u < 1$, the computed value
satisfies

\[
\hat{p}(x) = \sum_{k=0}^{n} a_{k} (1 + \eta_{k}) x^{k},
\qquad \abs{\eta_{k}} \leq \gamma_{2n} \approx 2 n u.
\]

This is a backward-error statement of the strongest kind. The computed
$\hat{p}(x)$ is the \emph{exact} value at $x$ of a polynomial whose
coefficients are each perturbed by no more than a relative $2 n u$. Horner's
rule is backward stable: it returns the exact answer to a nearby problem,
and the nearby problem differs from the posed one only by coefficient
perturbations at the level of the arithmetic.

The forward error follows by multiplying the backward error by the
conditioning. The perturbation $\delta a_{k} = a_{k} \eta_{k}$ changes the
value by $\sum_{k} \delta a_{k} x^{k}$, bounded in magnitude by $\gamma_{2n}
\sum_{k} \abs{a_{k}} \abs{x}^{k}$. Dividing by $\abs{p(x)}$ gives the
relative forward error bound

\[
\frac{\abs{\hat{p}(x) - p(x)}}{\abs{p(x)}} \leq \gamma_{2n}
\frac{\sum_{k=0}^{n} \abs{a_{k}} \abs{x}^{k}}{\abs{p(x)}}
= \gamma_{2n} \cdot \kappa_{\mathrm{Horner}}(x),
\]

and the ratio $\kappa_{\mathrm{Horner}}(x) = \sum_{k} \abs{a_{k}}
\abs{x}^{k} / \abs{p(x)}$ is exactly the condition number of polynomial
evaluation that the second worked example of section 16.2 quoted. The
slogan is now an inequality with named parts: the backward error is
$\gamma_{2n} \approx 2 n u$, set by the algorithm and the arithmetic; the
condition number is $\kappa_{\mathrm{Horner}}$, set by the problem; the
forward error is bounded by their product.

\begin{example}
\textbf{Horner near a root, by the numbers.} Evaluate $p(x) = (x-2)^{3}
= x^{3} - 6 x^{2} + 12 x - 8$ at $x = 2.001$ in binary64. The exact value
is $(0.001)^{3} = 10^{-9}$. The condition number is
$\kappa_{\mathrm{Horner}} = (\abs{x}^{3} + 6\abs{x}^{2} + 12\abs{x} + 8)
/ \abs{p(x)}$, and at $x = 2.001$ the numerator is $8.012 + 24.02 + 24.01
+ 8 = 64.05$ while the denominator is $10^{-9}$, so
$\kappa_{\mathrm{Horner}} \approx 6.4 \times 10^{10}$. The forward-error
bound is $\gamma_{6} \cdot \kappa \approx 6 u \cdot 6.4 \times 10^{10}
\approx 6.7 \times 10^{-16} \cdot 6.4 \times 10^{10} \approx 4 \times
10^{-5}$. The engineer therefore expects about
$-\log_{10}(4 \times 10^{-5}) \approx 4$ correct digits in the
Horner-evaluated value, and a direct computation confirms roughly four:
the computed value at $x = 2.001$ lands near $1.00004 \times 10^{-9}$
against the true $10^{-9}$. The algorithm is blameless; Horner delivered
the exact value of a polynomial whose coefficients moved by a relative
$6 u$, and the eleven-decade conditioning of evaluation near the triple
root turned that microscopic coefficient perturbation into a four-digit
forward error. The remedy is not a better evaluator but a better
representation: the factored form $(x - 2)^{3}$ has condition number $3$
at this point (the condition number of the cubing operation), and recovers
all sixteen digits. The reader who runs the backward analysis once, and
watches the four-digit prediction match the four-digit observation,
internalises the decomposition that the rest of numerical analysis rests
on.
\end{example}

\begin{mastery}
\textbf{Conditioning is the problem, stability is the algorithm.} The two
words are routinely conflated, and the conflation is the source of the most
common misdiagnosis in numerical work: reaching for higher precision when
the algorithm is at fault, or rewriting the algorithm when the problem is
at fault. The clean separation is this. \engterm{Conditioning}
$\kappa$ is a property of the map $f$ and the input $x$ alone; it is
defined before any algorithm is chosen, it does not change when the
language or the library changes, and no algorithm can compute $f(x)$ to a
relative accuracy better than $\kappa \cdot u$ in working precision $u$.
\engterm{Stability} (backward stability) is a property of the algorithm
$\hat{f}$ that approximates $f$; a backward-stable algorithm returns the
exact value of $f$ at some input within a relative $O(u)$ of $x$,
regardless of $\kappa$. The two combine in the master inequality
$\text{relative forward error} \lesssim \kappa \cdot (\text{backward
error})$, and the diagnostic that follows is mechanical. Compute $\kappa$.
If $\kappa$ is small and the answer is wrong, the algorithm is unstable:
rewrite it, and higher precision will not help much. If $\kappa$ is large,
the problem has lost the digits no matter what: change the problem
(reformulate, refactor, rescale, recondition) or accept fewer digits and
report the loss. The four worked examples of this section are the four
quadrants of that diagnostic: $1 - \cos x$ and $\sqrt{x+1} - \sqrt{x}$ are
the small-$\kappa$, unstable-algorithm cases that rewriting fixes; the
expanded $(x-1)^{6}$ and the Wilkinson polynomial are the large-$\kappa$
cases that only reformulation fixes. An engineer who carries the two words
as distinct quantities never asks higher precision to solve an algorithm
problem, and never blames an algorithm for digits the problem never had.
\end{mastery}

\section{Root-finding: bisection, Newton, secant}\pathtag{core}

The root-finding problem is: given $f: \R \to \R$ continuous,
find $x^{*}$ with $f(x^{*}) = 0$. The problem appears in every
volume that follows: equilibrium states in thermodynamics, fixed
points in control, equation-of-state inversions in fluid mechanics,
optimality conditions in optimisation. The three algorithms
covered here span the working tradeoff between robustness and
speed.

\subsection{Bisection}

If $f$ is continuous on $[a, b]$ and $f(a) f(b) < 0$, the
intermediate-value theorem guarantees a root in $[a, b]$.
\engterm{Bisection} repeatedly halves the bracket: compute
$c = (a + b)/2$, evaluate $f(c)$, replace whichever endpoint
shares its sign. The interval length halves at each step, so
after $k$ iterations the bracket has length $(b - a)/2^{k}$, and
the error of the midpoint approximation is at most half that.
To reach a tolerance $\tau$, the iteration needs

\[
k \geq \log_{2}\!\left(\frac{b - a}{2 \tau}\right)
\]

steps. Bisection has \engterm{linear convergence} with
contraction factor $1/2$.

The strength of bisection is its robustness. It requires only
continuity and a sign-change bracket; it does not require
differentiability, smoothness, or a good initial guess; the
interval length is a guaranteed forward error. Its weakness is
its speed. Twenty bisection steps shrink the bracket by a factor
of $10^{6}$; reaching machine precision from a unit bracket
requires about 53 steps in double precision.

\subsection{Newton's method}

If $f$ is differentiable and $x_{0}$ is close enough to a simple
root $x^{*}$, the iteration

\[
x_{k+1} = x_{k} - \frac{f(x_{k})}{f'(x_{k})}
\]

converges to $x^{*}$ at \engterm{quadratic} rate: the error
$e_{k} = x_{k} - x^{*}$ satisfies, for $k$ large enough,

\[
\abs{e_{k+1}} \leq C \abs{e_{k}}^{2},
\]

where $C = \abs*{f''(x^{*})}/(2\abs*{f'(x^{*})})$
\cite[ch.~4]{text:moler-numerical}. The number of correct digits
roughly doubles per step. Five or six steps from a reasonable
initial guess is typical to reach double precision.

The strength of Newton is its speed. The weakness is its
fragility. The method requires $f'$ (analytic or numerical), can
diverge or cycle from a poor initial guess, slows to linear
convergence at multiple roots, and amplifies catastrophically when
$\abs{f'(x^{*})}$ is small. Engineers use Newton inside a
safeguarded outer routine: Newton step proposed; if the proposal
falls outside a bracket maintained by bisection or by trust-region
logic, fall back to bisection. The hybrid routine \texttt{brentq}
in scientific Python and the \texttt{fzero} routine in MATLAB
implement variants of this discipline.

\begin{example}
\textbf{Three ways Newton fails.} The fragility is concrete, not
abstract. First, divergence from a poor start: for $f(x) =
\arctan(x)$, the unique root is $x^{*} = 0$, but Newton from any
$\abs{x_{0}} \gtrsim 1.39$ overshoots so far that successive iterates
grow in magnitude and diverge, because the Newton step $x -
\arctan(x)(1 + x^{2})$ exceeds $\abs{x}$ once $\abs{x}$ is large.
Second, a two-cycle: for $f(x) = x^{3} - 2x + 2$, Newton from
$x_{0} = 0$ produces $x_{1} = 1$, and from $x_{1} = 1$ produces
$x_{2} = 0$, oscillating forever between $0$ and $1$ without
approaching the real root near $-1.77$. Third, the multiple-root
slowdown: for $f(x) = (x - 1)^{2}$, the root $x^{*} = 1$ has
multiplicity $2$, $f'(x^{*}) = 0$, and Newton converges only
linearly with ratio $1/2$ rather than quadratically, halving the
error per step like bisection while paying for two evaluations. Each
failure violates a precondition of the quadratic-convergence theorem
(good initial guess, no cycling, simple root), and each is invisible
to a Newton loop that lacks the bracket safeguard. The safeguarded
hybrid of the design exercises detects all three: a step that leaves
the bracket triggers a bisection fallback, and a stalled bracket
width flags the multiple-root case.
\end{example}

\subsection{Secant method}

The \engterm{secant method} replaces the analytic derivative
with a finite-difference approximation across two iterates,

\[
x_{k+1} = x_{k} - f(x_{k}) \frac{x_{k} - x_{k-1}}
{f(x_{k}) - f(x_{k-1})}.
\]

The method has \engterm{superlinear convergence} of order
$\varphi = (1 + \sqrt{5})/2 \approx 1.618$, the golden ratio. It
needs no derivative and one function evaluation per step instead
of two. The method's robustness sits between bisection and
Newton; like Newton, it can diverge from a poor initial pair.

\begin{figure}[ht]
\centering
\begin{tikzpicture}[x=0.32cm, y=0.32cm, font=\footnotesize]

% Axis: iteration count vs log10 |error|
\draw[->, thick] (0,0) -- (22,0) node[right] {iteration $k$};
\draw[->, thick] (0,0) -- (0,18) node[above, align=center] {$-\log_{10}\abs{e_{k}}$\\(correct decimal digits)};

\foreach \k in {0,5,10,15,20} {
  \draw (\k,0) -- (\k,-0.3);
  \node[anchor=north, font=\scriptsize] at (\k,-0.3) {\k};
}
\foreach \d in {0,4,8,12,16} {
  \draw (0,\d) -- (-0.3,\d);
  \node[anchor=east, font=\scriptsize] at (-0.3,\d) {\d};
}

% Bisection: linear, slope = log10(2) per step ~ 0.301 digits/step
\draw[thick, engblue] plot coordinates {
  (0,0.5) (5,2.0) (10,3.5) (15,5.0) (20,6.5)
};
\node[engblue, anchor=west, font=\scriptsize] at (20.2,6.5)
  {bisection};

% Secant: superlinear, order phi ~ 1.618
\draw[thick, engamber] plot[smooth] coordinates {
  (0,0.5) (1,0.7) (2,1.2) (3,2.1) (4,3.5) (5,5.5) (6,8.3) (7,12.5) (8,16) (15,16) (20,16)
};
\node[engamber!70!black, anchor=west, font=\scriptsize] at (8.5,14.5)
  {secant ($\varphi \approx 1.618$)};

% Newton: quadratic
\draw[thick, engred] plot[smooth] coordinates {
  (0,0.5) (1,1.0) (2,2.0) (3,4.0) (4,8.0) (5,16) (10,16) (20,16)
};
\node[engred, anchor=west, font=\scriptsize] at (5.3,16.5)
  {Newton (quadratic)};

% Double-precision ceiling
\draw[thick, dashed, gray] (0,16) -- (22,16);
\node[gray, anchor=west, font=\scriptsize] at (0.5,16.7)
  {binary64 precision ceiling};

\end{tikzpicture}
\caption[Convergence orders of bisection, secant, and Newton]{Convergence
orders of bisection, the secant method, and Newton's method on a smooth
problem with a simple root and a good initial guess. Bisection adds
$\log_{10}(2) \approx 0.30$ correct digits per iteration; the secant
method has order $\varphi = (1 + \sqrt{5})/2 \approx 1.618$; Newton has
order $2$ and roughly doubles its correct-digit count per step. All
three saturate at the binary64 precision ceiling near 16 decimal
digits. The plot illustrates why the robustness-vs-speed tradeoff of
section 16.3 favours Newton when its precondition is met and
bisection only as a fall-back. Source: authors' computation; the
slopes are theoretical predictions from the convergence-order theory
in \cite[ch.~4]{text:moler-numerical}.}
\label{fig:vol02:ch16:convergence-orders}
\end{figure}


Figure~\ref{fig:vol02:ch16:convergence-orders} reads the three rates
side-by-side on a problem all three solve (the smooth scalar root of
$\cos x - x = 0$ on $[0, 1]$). Bisection adds approximately $0.30$
correct decimal digits per step; the secant method's rate climbs
geometrically toward order $\varphi \approx 1.618$; Newton roughly
doubles the digit count per step. All three saturate near 16 digits
at the binary64 precision ceiling. The cost per step is one function
evaluation for bisection and secant; two (one $f$, one $f'$) for
Newton.

\begin{example}
\textbf{A secant trace, and the superlinear order read from it.} Solve
$f(x) = x^{2} - 2 = 0$ for the positive root $x^{*} = \sqrt{2} =
1.4142135623730951$, starting from $x_{0} = 1$ and $x_{1} = 2$. The secant
update is $x_{k+1} = x_{k} - f(x_{k})(x_{k} - x_{k-1})/(f(x_{k}) -
f(x_{k-1}))$. Step by step: $f(x_{0}) = -1$, $f(x_{1}) = 2$, so $x_{2} =
2 - 2(2 - 1)/(2 - (-1)) = 2 - 2/3 = 1.3333333$, $e_{2} = -8.09 \times
10^{-2}$. Then $f(x_{2}) = 1.3333^{2} - 2 = -0.2222$, $x_{3} = 1.3333 -
(-0.2222)(1.3333 - 2)/(-0.2222 - 2) = 1.4000000$, $e_{3} = -1.42 \times
10^{-2}$. Then $f(x_{3}) = -0.04$, $x_{4} = 1.4146341$, $e_{4} = +4.21
\times 10^{-4}$; $x_{5} = 1.4142114$, $e_{5} = -2.1 \times 10^{-6}$;
$x_{6} = 1.41421356205$, $e_{6} = -3.2 \times 10^{-11}$; $x_{7}$ at the
binary64 floor. The order shows itself in the ratio $\log\abs{e_{k+1}} /
\log\abs{e_{k}}$, which approaches $\varphi = 1.618$: the digit count is
multiplied by about $1.6$ per step, between bisection's additive $0.3$ and
Newton's doubling. Reading the per-step error reduction directly,
$-\log_{10}\abs{e_{k}}$ runs $1.09, 1.85, 3.38, 5.68, 10.5, \ldots$, and
the successive increments $0.76, 1.53, 2.30, 4.82$ themselves grow by the
factor $\varphi$, which is the signature of order-$\varphi$ convergence.
The secant method reached eleven digits in six steps at one function
evaluation per step, where Newton would reach the same in five steps at
two evaluations each; counting evaluations, the secant method's efficiency
index $\varphi^{1} = 1.618$ beats Newton's $2^{1/2} = 1.414$, which is why
a derivative-free secant step is the working default when $f'$ is expensive
to form.
\end{example}

\subsection{Fixed-point iteration}

Many root-finding problems arrive already in, or are easily
rewritten into, \engterm{fixed-point} form $x = g(x)$, where a root
of $f$ is a fixed point of $g$. The iteration $x_{k+1} = g(x_{k})$
converges to a fixed point $x^{*}$ when $g$ is a contraction near
$x^{*}$. The governing quantity is $\abs{g'(x^{*})}$: by the
mean-value theorem $e_{k+1} = g(x_{k}) - g(x^{*}) = g'(\xi_{k})
e_{k}$, so

\[
\abs{e_{k+1}} \approx \abs{g'(x^{*})} \, \abs{e_{k}}
\]

near the fixed point. The iteration converges linearly with
asymptotic rate $\abs{g'(x^{*})}$ when $\abs{g'(x^{*})} < 1$, and
diverges when $\abs{g'(x^{*})} > 1$. The case $g'(x^{*}) = 0$ is
the source of Newton's quadratic rate: Newton is the fixed-point
iteration with $g(x) = x - f(x)/f'(x)$, for which $g'(x^{*}) =
f(x^{*}) f''(x^{*}) / f'(x^{*})^{2} = 0$ at a simple root (since
$f(x^{*}) = 0$). The vanishing first derivative of the iteration map
is the structural reason Newton converges quadratically and a
generic fixed-point iteration converges only linearly.

\begin{example}
\textbf{Two rearrangements, two fates.} To solve $x^{3} - x - 2 =
0$ near $x^{*} \approx 1.5214$, two natural rearrangements are
$g_{1}(x) = x^{3} - 2$ (from $x = x^{3} - 2$) and $g_{2}(x) =
(x + 2)^{1/3}$ (from $x^{3} = x + 2$). At the root, $g_{1}'(x) =
3x^{2}$ gives $g_{1}'(x^{*}) \approx 6.94 > 1$: the iteration
$x_{k+1} = x_{k}^{3} - 2$ diverges from any start near the root.
The map $g_{2}'(x) = \tfrac{1}{3}(x + 2)^{-2/3}$ gives
$g_{2}'(x^{*}) \approx 0.135 < 1$: the iteration $x_{k+1} =
(x_{k} + 2)^{1/3}$ converges linearly, gaining about
$-\log_{10}(0.135) \approx 0.87$ correct digits per step, roughly
eighteen steps to machine precision. The same root, the same
problem, two algebraically equivalent rearrangements, and one
converges while the other flies apart. The discipline is to inspect
$\abs{g'(x^{*})}$ before iterating, or to estimate it from the
observed per-step error ratio in the first few iterations and to
abandon the rearrangement when the ratio exceeds one.
\end{example}

\begin{mastery}
\textbf{Aitken acceleration and Steffensen's method.} A linearly
convergent sequence $x_{k} \to x^{*}$ with asymptotic ratio $r =
g'(x^{*})$ satisfies $x_{k+1} - x^{*} \approx r(x_{k} - x^{*})$. Three
consecutive iterates then determine $r$ and, by extrapolation,
$x^{*}$ itself. \engterm{Aitken's $\Delta^{2}$ process} forms

\[
\hat{x}_{k} = x_{k} - \frac{(x_{k+1} - x_{k})^{2}}
{x_{k+2} - 2x_{k+1} + x_{k}},
\]

which converges to $x^{*}$ faster than the original sequence. Folding
Aitken extrapolation back into the iteration (apply $g$ twice, then
extrapolate, then restart) is \engterm{Steffensen's method}, which
recovers quadratic convergence from a linearly convergent fixed-point
map without requiring the derivative $g'$. Steffensen is the
working answer when the engineer has a contraction map but no cheap
derivative and wants Newton-like speed. The cost is two function
evaluations per accelerated step and a fragility near $x_{k+2} -
2x_{k+1} + x_{k} = 0$, where the denominator cancels and the
extrapolation must be guarded.
\end{mastery}

\subsection{Choosing among the three}

The working choice depends on three properties of the problem.
Is $f$ smooth, with a known derivative, near a simple root, with
a good initial guess available: use Newton, with a safeguarded
fallback. Is $f$ continuous, with a sign-change bracket
identifiable, with no derivative information: use bisection
followed by Newton (or by a hybrid like Brent's method). Is the
derivative expensive but the function smooth: use the secant
method, with the same safeguarding pattern.

The engineer should not write a custom Newton or secant routine
for production work; \texttt{scipy.optimize.brentq},
\texttt{scipy.optimize.newton}, MATLAB \texttt{fzero}, and the
GNU Scientific Library equivalents are tested, safeguarded, and
documented. The discipline this section teaches is the choice
among them, not the implementation.

\begin{archetype}[Failure (numerical)]
A numerical method has a precondition, a convergence rate, and a
failure mode. The convergence rate is the quantity by which the
method earns its place; the failure mode is the quantity that
costs the engineer the answer when the precondition is violated.
Newton at a simple root with a good initial guess earns quadratic
convergence; Newton at a multiple root degrades to linear; Newton
at a stationary point of $f$ diverges. The discipline is to know
the precondition before deploying the method, and to safeguard
the deployment for the cases where the precondition is silently
violated. The same archetype recurs across this chapter:
trapezoidal quadrature on a discontinuous integrand, RK4 on a
stiff system, Gaussian elimination without pivoting on a matrix
with a small pivot, and naive summation on a vector with
catastrophic cancellation. Each algorithm earns its rate from a
property the input may or may not satisfy; the engineer's job is
to verify, not to assume.
\end{archetype}

\section{Interpolation and approximation}\pathtag{standard}

A table of data points, or a function too expensive to evaluate
everywhere, leaves the engineer needing a value between the samples.
\engterm{Interpolation} builds a function that passes through the
given points; \engterm{approximation} builds a function that stays
close to the given points without necessarily passing through them.
The distinction matters when the data carry noise: interpolating
noisy data reproduces the noise, where a least-squares
approximation of lower degree filters it. This section develops the
interpolation apparatus, the way it fails (the Runge phenomenon),
and the piecewise remedy (splines).

\subsection{Polynomial interpolation exists and is unique}

Given $n + 1$ distinct nodes $x_{0}, \ldots, x_{n}$ and values
$y_{0}, \ldots, y_{n}$, there is exactly one polynomial of degree at
most $n$ that passes through all $n + 1$ points. Existence and
uniqueness follow from the Vandermonde system: writing the
interpolant as $p(x) = \sum_{j=0}^{n} c_{j} x^{j}$ and imposing
$p(x_{i}) = y_{i}$ gives a linear system whose matrix is the
Vandermonde matrix $V_{ij} = x_{i}^{j}$, nonsingular when the nodes
are distinct (its determinant is $\prod_{i < j}(x_{j} - x_{i})$).
Solving the Vandermonde system directly is the wrong way to compute
the interpolant: the Vandermonde matrix is notoriously ill
conditioned, with $\kappa$ growing exponentially in $n$ for
equispaced nodes. The Lagrange and Newton forms compute the same
unique polynomial without forming or inverting the Vandermonde
matrix.

\subsection{The Lagrange form}

The \engterm{Lagrange form} writes the interpolant as a weighted
sum of \engterm{cardinal polynomials},

\[
p(x) = \sum_{j=0}^{n} y_{j} \, \ell_{j}(x), \qquad
\ell_{j}(x) = \prod_{\substack{i=0 \\ i \neq j}}^{n}
\frac{x - x_{i}}{x_{j} - x_{i}}.
\]

Each $\ell_{j}$ is the degree-$n$ polynomial equal to $1$ at
$x_{j}$ and $0$ at every other node, so $p(x_{i}) = y_{i}$ by
construction. The Lagrange form is the cleanest expression for
theory and for proofs; it is poor for computation as written,
because adding a new data point requires recomputing every cardinal
polynomial. The \engterm{barycentric} rewrite of the Lagrange form
fixes the computational cost and is numerically stable; it is the
form a production interpolation routine actually uses.

\subsection{Newton's divided differences}

The \engterm{Newton divided-difference form} builds the interpolant
incrementally, so that adding a data point costs one more term
rather than a full recomputation. Define the divided differences
recursively:

\begin{align*}
f[x_{i}] &= y_{i}, \\
f[x_{i}, \ldots, x_{i+k}] &= \frac{f[x_{i+1}, \ldots, x_{i+k}]
- f[x_{i}, \ldots, x_{i+k-1}]}{x_{i+k} - x_{i}}.
\end{align*}

The interpolant is then

\[
p(x) = f[x_{0}] + f[x_{0}, x_{1}](x - x_{0}) + \cdots
+ f[x_{0}, \ldots, x_{n}] \prod_{i=0}^{n-1}(x - x_{i}).
\]

The leading divided difference $f[x_{0}, \ldots, x_{n}]$ equals
$f^{(n)}(\xi)/n!$ for some $\xi$ in the interval spanned by the
nodes when $f \in C^{n}$, which connects the divided difference to
the Taylor coefficient and underwrites the interpolation-error
formula below.

\begin{example}
\textbf{A three-point interpolant computed two ways.} Interpolate
$f(x) = \ln x$ at $x_{0} = 1, x_{1} = 2, x_{2} = 4$, where $y_{0} =
0$, $y_{1} = \ln 2 = 0.6931$, $y_{2} = \ln 4 = 1.3863$. Divided
differences: $f[x_{0}, x_{1}] = (0.6931 - 0)/(2 - 1) = 0.6931$;
$f[x_{1}, x_{2}] = (1.3863 - 0.6931)/(4 - 2) = 0.3466$;
$f[x_{0}, x_{1}, x_{2}] = (0.3466 - 0.6931)/(4 - 1) = -0.1155$. The
Newton interpolant is $p(x) = 0 + 0.6931(x - 1) - 0.1155(x - 1)(x -
2)$. Evaluating at $x = 3$: $p(3) = 0.6931 \cdot 2 - 0.1155 \cdot 2
\cdot 1 = 1.3863 - 0.2310 = 1.1553$, against the true $\ln 3 =
1.0986$, an error of $0.057$. The same value comes from the
Lagrange form $p(3) = 0 \cdot \ell_{0}(3) + 0.6931 \cdot \ell_{1}(3)
+ 1.3863 \cdot \ell_{2}(3)$ with $\ell_{1}(3) = \frac{(3-1)(3-4)}
{(2-1)(2-4)} = 1$ and $\ell_{2}(3) = \frac{(3-1)(3-2)}{(4-1)(4-2)} =
\tfrac{1}{3}$, giving $0.6931 + 0.4621 = 1.1552$, the same
polynomial to rounding. The $0.057$ error reflects the curvature of
$\ln$ that a quadratic cannot capture across so wide an interval;
the interpolation-error formula below predicts its size.
\end{example}

\subsection{The interpolation error formula}

For $f \in C^{n+1}$ on an interval containing the nodes, the error
of the degree-$n$ interpolant at any point $x$ is

\[
f(x) - p(x) = \frac{f^{(n+1)}(\xi)}{(n+1)!}
\prod_{i=0}^{n}(x - x_{i})
\]

for some $\xi$ in the interval. Two factors govern the error. The
derivative factor $f^{(n+1)}(\xi)/(n+1)!$ depends on the function:
for an entire function whose high derivatives stay bounded, it
shrinks with $n$. The \engterm{node polynomial} $\prod(x - x_{i})$
depends only on the node placement, and its size between the nodes
is the engineer's design variable. Equispaced nodes make the node
polynomial large near the interval endpoints, which is the seed of
the Runge phenomenon. Chebyshev nodes minimise the maximum of the
node polynomial and tame the error.

\subsection{The Runge phenomenon}

Raising the degree of an equispaced polynomial interpolant does not
in general improve the fit. For $f(x) = 1/(1 + 25 x^{2})$ on
$[-1, 1]$, the classic example of Runge (1901), the equispaced
interpolant develops oscillations near the endpoints whose amplitude
grows with the degree; the maximum-norm error diverges as $n \to
\infty$. The derivative factor $f^{(n+1)}$ grows fast enough, and
the equispaced node polynomial grows large enough near the
endpoints, that their product blows up.

\begin{figure}[ht]
\centering
\begin{tikzpicture}[x=1.05cm, y=2.6cm, font=\footnotesize]

% Runge function f(x) = 1/(1 + 25 x^2) on [-1, 1], with the degree-10
% equispaced interpolant overshooting near the endpoints.

% Axes
\draw[->, thick] (-1.25,0) -- (1.25,0) node[right] {$x$};
\draw[->, thick] (0,-0.85) -- (0,1.35) node[above] {$y$};

\foreach \x in {-1,-0.5,0.5,1} {
  \draw (\x,0) -- (\x,-0.04);
  \node[anchor=north, font=\scriptsize] at (\x,-0.04) {$\x$};
}
\foreach \y in {0.5,1} {
  \draw (0,\y) -- (-0.04,\y);
  \node[anchor=east, font=\scriptsize] at (-0.04,\y) {$\y$};
}

% Runge function (smooth), sampled densely
\draw[thick, engblue] plot[smooth, samples=80, domain=-1:1]
  ({\x}, {1/(1 + 25*\x*\x)});
\node[engblue, anchor=west, font=\scriptsize] at (0.12,0.95)
  {$f(x) = \dfrac{1}{1 + 25 x^{2}}$};

% Degree-10 equispaced polynomial interpolant: caricature with the
% characteristic endpoint oscillation (overshoot near +/- 1).
\draw[thick, engred] plot[smooth] coordinates {
  (-1.0,1.0) (-0.92,-0.35) (-0.8,0.45) (-0.6,0.06) (-0.4,0.10)
  (-0.2,0.50) (0,1.0) (0.2,0.50) (0.4,0.10) (0.6,0.06)
  (0.8,0.45) (0.92,-0.35) (1.0,1.0)
};
\node[engred, anchor=north, font=\scriptsize] at (-0.78,-0.42)
  {degree-10 equispaced interpolant};

% Interpolation nodes (11 equispaced)
\foreach \x in {-1,-0.8,-0.6,-0.4,-0.2,0,0.2,0.4,0.6,0.8,1} {
  \fill[black] (\x, {1/(1 + 25*\x*\x)}) circle (1.3pt);
}

\end{tikzpicture}
\caption[Runge phenomenon]{The Runge phenomenon. Interpolating the
smooth function $f(x) = 1/(1 + 25 x^{2})$ at $11$ equally spaced nodes
on $[-1, 1]$ with a single degree-$10$ polynomial (red) produces large
oscillations near the endpoints, with overshoot growing as the degree
rises. The interpolant matches $f$ exactly at the nodes (black dots) and
diverges from it between them. Raising the polynomial degree on an
equispaced grid makes the overshoot worse, not better: the interpolation
error in the maximum norm grows without bound as the degree increases.
The cure is to abandon the equispaced grid (Chebyshev nodes cluster at
the endpoints and tame the oscillation) or to abandon the single
high-degree polynomial (piecewise cubic splines stay bounded). Source:
authors' computation; the example is Runge's 1901 original, discussed in
\cite[ch.~5]{text:v2ch05-burden-faires-numerical-2015}.}
\label{fig:vol02:ch16:runge-phenomenon}
\end{figure}


Figure~\ref{fig:vol02:ch16:runge-phenomenon} shows the
degree-$10$ equispaced interpolant overshooting near $x = \pm 1$.
The cure is one of two changes. Cluster the nodes at the endpoints
using the \engterm{Chebyshev nodes} $x_{k} = \cos(k\pi/n)$, whose
node polynomial has the smallest possible maximum and whose
interpolant converges for any function analytic on the interval.
Abandon the single high-degree polynomial for a piecewise low-degree
construction, the spline.

\begin{estimation}
\textbf{Question.} An engineer has $21$ tabulated values of a smooth
calibration curve on a uniform grid and wants the value at a point
between two table entries. Before computing, estimate whether a
single degree-$20$ polynomial interpolant or a cubic spline through
the same $21$ points will give the smaller worst-case error, and by
roughly what factor.

\smallskip
\textit{Estimate, before reading on.}

\smallskip
The degree-$20$ equispaced interpolant sits squarely in the Runge
regime: the node polynomial's maximum near the endpoints grows like
$2^{n}$ with $n = 20$, so the endpoint overshoot is governed by a
factor near $2^{20} \approx 10^{6}$ relative to the interior, and
the maximum-norm error is dominated by that overshoot. The cubic
spline has error $O(h^{4})$ with $h = 1/20$, giving a factor near
$(1/20)^{4} \approx 6 \times 10^{-6}$ times a curvature constant of
order the fourth derivative. The spline wins, and the rough ratio of
the two worst-case errors is the product of the polynomial's
endpoint blow-up and the spline's $h^{4}$ decay, plausibly five to
ten orders of magnitude in the spline's favour for a curve with
modest higher derivatives. The order-of-magnitude lesson governs the
working choice: for interpolation on a fixed equispaced grid of more
than a handful of points, reach for the spline, not the high-degree
polynomial. The exact factor depends on the calibration curve's
derivatives, but no realistic smooth curve reverses the ranking.
\end{estimation}

\subsection{Cubic splines}

A \engterm{cubic spline} interpolant is a piecewise cubic
polynomial, one cubic per subinterval $[x_{k}, x_{k+1}]$, joined so
that the interpolant and its first two derivatives are continuous
across every interior node. The continuity conditions plus the
interpolation conditions leave two degrees of freedom, fixed by a
boundary condition: the \engterm{natural} spline sets the second
derivative to zero at both ends; the \engterm{clamped} spline sets
the first derivative to specified values; the \engterm{not-a-knot}
condition (the default in most libraries) requires the third
derivative to be continuous across the first and last interior
nodes. Assembling the continuity conditions produces a tridiagonal
linear system in the interior second derivatives, solvable in
$O(n)$ work.

\begin{figure}[ht]
\centering
\begin{tikzpicture}[x=1.05cm, y=2.6cm, font=\footnotesize]

% Same Runge target, but now contrasting the equispaced high-degree
% polynomial (oscillates) with a natural cubic spline (stays bounded
% and tracks the data).

\draw[->, thick] (-1.25,0) -- (1.25,0) node[right] {$x$};
\draw[->, thick] (0,-0.25) -- (0,1.35) node[above] {$y$};

\foreach \x in {-1,-0.5,0.5,1} {
  \draw (\x,0) -- (\x,-0.04);
  \node[anchor=north, font=\scriptsize] at (\x,-0.04) {$\x$};
}
\foreach \y in {0.5,1} {
  \draw (0,\y) -- (-0.04,\y);
  \node[anchor=east, font=\scriptsize] at (-0.04,\y) {$\y$};
}

% Underlying Runge function
\draw[thick, engblue] plot[smooth, samples=80, domain=-1:1]
  ({\x}, {1/(1 + 25*\x*\x)});
\node[engblue, anchor=west, font=\scriptsize] at (0.12,0.95)
  {$f(x)$ (target)};

% Natural cubic spline through the 11 equispaced nodes: bounded,
% tracks the data without endpoint overshoot.
\draw[thick, engamber!80!black] plot[smooth] coordinates {
  (-1.0,0.0385) (-0.8,0.0588) (-0.6,0.1) (-0.4,0.2) (-0.2,0.5)
  (0,1.0) (0.2,0.5) (0.4,0.2) (0.6,0.1) (0.8,0.0588) (1.0,0.0385)
};
\node[engamber!70!black, anchor=west, font=\scriptsize] at (0.25,0.62)
  {cubic spline};

% Interpolation nodes
\foreach \x in {-1,-0.8,-0.6,-0.4,-0.2,0,0.2,0.4,0.6,0.8,1} {
  \fill[black] (\x, {1/(1 + 25*\x*\x)}) circle (1.3pt);
}

\end{tikzpicture}
\caption[Cubic spline versus high-degree polynomial]{The natural cubic
spline (amber) through the same $11$ equispaced nodes that produced the
oscillating degree-$10$ polynomial of
Figure~\ref{fig:vol02:ch16:runge-phenomenon}. The spline is a
piecewise-cubic that matches the data at the nodes and enforces
continuity of the function and its first two derivatives across each
knot; it stays bounded between $0$ and $1$ and tracks the target. The
spline trades the polynomial's global smoothness (a single analytic
expression) for local control (each piece sees only its neighbours),
and the trade buys boundedness. The maximum-norm error of the spline
falls as the node spacing shrinks, where the equispaced polynomial's
error grows. Source: authors' computation; cubic-spline interpolation as
in \cite[ch.~3]{text:v2ch05-burden-faires-numerical-2015}.}
\label{fig:vol02:ch16:spline-vs-poly}
\end{figure}


Figure~\ref{fig:vol02:ch16:spline-vs-poly} contrasts the spline
with the oscillating polynomial on the Runge data. The spline trades
the polynomial's single analytic expression for local control: each
cubic piece sees only its immediate neighbours, so a perturbation in
one data value moves the interpolant only locally, and the
maximum-norm error falls as the node spacing shrinks (as $O(h^{4})$
for a clamped spline of a smooth function). The spline is the
working default for interpolating tabulated engineering data, in CAD
geometry, and in the smooth reconstruction of sampled signals.

\begin{example}
\textbf{A natural cubic spline built by hand.} Construct the natural cubic
spline through the four points $(0, 0)$, $(1, 1)$, $(2, 0)$, $(3, 1)$ on
the uniform grid $h = 1$. The interior second derivatives $M_{1}, M_{2}$
at the two interior nodes satisfy the tridiagonal system that the
continuity-of-first-derivative conditions produce; for a uniform grid the
standard form is, at interior node $k$,

\[
M_{k-1} + 4 M_{k} + M_{k+1} = \frac{6}{h^{2}}(y_{k-1} - 2 y_{k} + y_{k+1}),
\]

with the natural end conditions $M_{0} = M_{3} = 0$. The right-hand sides
are second differences: at node $1$, $(6/1)(0 - 2 \cdot 1 + 0) = -12$; at
node $2$, $(6/1)(1 - 2 \cdot 0 + 1) = 12$. With $M_{0} = M_{3} = 0$ the
system reduces to

\[
\begin{aligned}
4 M_{1} + M_{2} &= -12, \\
M_{1} + 4 M_{2} &= 12.
\end{aligned}
\]

Solve by elimination: from the first, $M_{1} = (-12 - M_{2})/4$; substitute
into the second, $(-12 - M_{2})/4 + 4 M_{2} = 12$, giving $-3 + (15/4) M_{2}
= 12$, so $M_{2} = (15)(4)/(15) = 4$ and $M_{1} = (-12 - 4)/4 = -4$. The
second derivatives at the four nodes are therefore $M = (0, -4, 4, 0)$. On
each subinterval $[x_{k}, x_{k+1}]$ the cubic piece is

\[
S_{k}(x) = M_{k}\frac{(x_{k+1} - x)^{3}}{6h} + M_{k+1}\frac{(x - x_{k})^{3}}
{6h} + \Big(\frac{y_{k}}{h} - \frac{M_{k} h}{6}\Big)(x_{k+1} - x)
+ \Big(\frac{y_{k+1}}{h} - \frac{M_{k+1} h}{6}\Big)(x - x_{k}).
\]

On the first piece $[0, 1]$, with $M_{0} = 0$, $M_{1} = -4$, $y_{0} = 0$,
$y_{1} = 1$: $S_{0}(x) = -4 \frac{x^{3}}{6} + (0)(1 - x) + (1 + \frac{4}{6})
x = -\frac{2}{3}x^{3} + \frac{5}{3}x$. A check at the endpoints gives
$S_{0}(0) = 0$ and $S_{0}(1) = -\frac{2}{3} + \frac{5}{3} = 1$, matching
the data, and $S_{0}''(0) = 0$, satisfying the natural left condition.
Evaluate the spline at the midpoint $x = 0.5$: $S_{0}(0.5) = -\frac{2}{3}
(0.125) + \frac{5}{3}(0.5) = -0.0833 + 0.8333 = 0.75$. The interpolant
rises to $0.75$ halfway to the first data point, where a straight-line
interpolant would give $0.5$; the spline's curvature, fixed by the
continuity conditions and the data on both sides, lifts the midpoint. The
$O(n)$ tridiagonal solve that the libraries run is exactly this elimination
at scale; carrying it once by hand on four points shows the engineer where
the smoothness comes from and why the spline's influence is local: each
$M_{k}$ couples only to its neighbours, so a single bad data value
perturbs the curve over a few intervals and not the whole range.
\end{example}

\begin{mastery}
\textbf{Why Chebyshev nodes are optimal.} Among all choices of $n +
1$ nodes on $[-1, 1]$, the Chebyshev nodes minimise the maximum
over $[-1, 1]$ of the node polynomial $\prod(x - x_{i})$. The
minimised maximum is $2^{-n}$, achieved when the node polynomial is
the monic Chebyshev polynomial $T_{n+1}/2^{n}$, which equioscillates
between $\pm 2^{-n}$ across the interval. The proof is a classic
equioscillation argument: any monic degree-$(n+1)$ polynomial with a
smaller maximum would have to cross the Chebyshev polynomial more
than $n + 1$ times, which is impossible for two distinct monic
polynomials of degree $n + 1$ whose difference has degree at most
$n$. Equispaced nodes, by contrast, give a node polynomial whose
maximum grows like $2^{n}/(e\, n \log n)$, exponentially worse. The
factor between the two, $2^{2n}$ at the endpoints, is the
quantitative seed of the Runge phenomenon, and the reason a
spectral method or a Chebyshev-based quadrature places its nodes by
the cosine rule rather than uniformly.
\end{mastery}

\section{Numerical linear algebra revisited}\pathtag{standard}

Volume II Chapter 9 introduced linear systems, LU factorisation,
QR factorisation, condition numbers, and the SVD at the level of
recognition. This section completes the apparatus at the level a
working engineer should be able to defend.

\subsection{LU with partial pivoting}

Gaussian elimination as written in Chapter 9 produces $\mat{A} =
\mat{L}\mat{U}$ when no zero pivot appears. In floating-point, a
\emph{small} pivot is almost as dangerous as a zero pivot: the
elimination scales subsequent rows by the inverse of the pivot,
and a small pivot amplifies the round-off error of those rows.
\engterm{Partial pivoting} fixes this by permuting rows at each
elimination step so that the pivot is the largest-magnitude
candidate in its column. The factorisation becomes

\[
\mat{P}\mat{A} = \mat{L}\mat{U},
\]

where $\mat{P}$ is a permutation matrix that records the row
swaps. The algorithm is backward stable in practice for almost all
matrices, with backward error of order $u$ times a quantity called
the \engterm{growth factor}, which is bounded by $2^{n-1}$ in
theory and is observed to be $O(n)$ in practice
\cite[ch.~22]{text:higham-numerical}.

Partial pivoting is not optional in production code. Every
library LU routine pivots; an engineer who writes Gaussian
elimination from scratch and omits pivoting has written an
algorithm whose stability the engineer cannot defend. The
discipline of section 9.8's refusal list applies: a reported
linear-system solve without partial pivoting (or full pivoting,
or a stability argument) is incomplete.

\begin{example}
\textbf{Why the growth factor matters.} Partial pivoting bounds each
multiplier by $1$ in magnitude, but it does not bound the
\emph{growth} of the matrix entries during elimination; the
worst-case growth factor $\rho = \max_{i,j,k}\abs{u_{ij}^{(k)}} /
\max_{i,j}\abs{a_{ij}}$ can reach $2^{n-1}$. The standard worst-case
matrix is

\[
\mat{A} = \begin{psmallmatrix}
1 & 0 & 0 & 1 \\
-1 & 1 & 0 & 1 \\
-1 & -1 & 1 & 1 \\
-1 & -1 & -1 & 1
\end{psmallmatrix},
\]

whose last column doubles at each elimination step: the bottom-right
entry grows from $1$ to $2$ to $4$ to $8 = 2^{3}$ for this $4 \times
4$ case, and to $2^{n-1}$ for the $n \times n$ analogue. Partial
pivoting does not swap any rows here, because every pivot is already
the column maximum, so the growth proceeds unchecked. The backward
error of the solve is bounded by $\rho \cdot n \cdot u$, so this
matrix loses about $n - 1$ bits of accuracy that a benign matrix
keeps. Such matrices are vanishingly rare in practice (the observed
growth factor on random and on application matrices is $O(\sqrt{n})$
to $O(n)$, never the exponential bound), which is why partial
pivoting is the production default despite its non-bound on growth.
The engineer who must guarantee the bound, rather than rely on the
empirical typical case, uses complete pivoting (which bounds growth
polynomially) at roughly double the cost, or monitors the computed
growth factor and re-solves with higher precision when it is large.
\end{example}

\subsection{QR by Householder reflections}

The Gram-Schmidt procedure of Chapter 9 produces $\mat{A} =
\mat{Q}\mat{R}$ in exact arithmetic, but in floating-point it
loses orthogonality: the computed $\mat{Q}$ has $\mat{Q}^{\transpose}
\mat{Q} - \mat{I}$ of size $\norm{\cdot}_{2} \approx \kappa(\mat{A})
u$, which can be much larger than $u$. \engterm{Modified
Gram-Schmidt} improves the situation but still loses orthogonality
on ill-conditioned matrices.

The \engterm{Householder QR} algorithm produces a $\mat{Q}$ that
is orthogonal to working precision regardless of the conditioning
of $\mat{A}$. The algorithm constructs a sequence of orthogonal
\engterm{Householder reflectors} $\mat{H}_{k}$, each of the form

\[
\mat{H} = \mat{I} - 2 \frac{\vect{v}\vect{v}^{\transpose}}
{\vect{v}^{\transpose}\vect{v}},
\]

chosen so that $\mat{H}_{k}$ zeros the subdiagonal of column $k$
of the partially triangularised matrix. The product $\mat{H}_{n}
\cdots \mat{H}_{1} \mat{A} = \mat{R}$ is upper triangular, and
$\mat{Q} = \mat{H}_{1}^{\transpose}\cdots\mat{H}_{n}^{\transpose}$
is orthogonal by construction. The factorisation is backward
stable: the computed $\hat{\mat{Q}}$ and $\hat{\mat{R}}$ satisfy
$\hat{\mat{Q}}\hat{\mat{R}} = \mat{A} + \mat{E}$ with
$\norm{\mat{E}}/\norm{\mat{A}} \leq c_{n} u$ for a small constant
$c_{n}$ \cite[ch.~10]{text:trefethen-bau1997}.

The Householder $\mat{Q}$ is rarely formed explicitly; the
algorithm stores the reflectors and applies them on demand to
$\vect{b}$ when solving $\mat{A}\vect{x} = \vect{b}$ via
$\mat{R}\vect{x} = \mat{Q}^{\transpose}\vect{b}$. The storage is
the same as the input matrix's, and the operation count is
$2 m n^{2} - \frac{2}{3} n^{3}$ flops for an $m \times n$ matrix
with $m \geq n$, twice that of LU but with the orthogonality
bonus.

\subsection{Singular value decomposition}

The SVD of Chapter 9 takes the form $\mat{A} = \mat{U}\mat{\Sigma}
\mat{V}^{\transpose}$. The standard algorithm for computing it is
the \engterm{Golub-Reinsch} two-phase routine: first reduce
$\mat{A}$ to bidiagonal form by left and right Householder
reflectors, then iterate an implicit QR algorithm on the
bidiagonal until the off-diagonal entries are below a tolerance.
The cost is $O(m n^{2})$ for an $m \times n$ matrix with $m \geq
n$, and the algorithm is backward stable
\cite[ch.~31]{text:trefethen-bau1997}.

The SVD's working role for the engineer is the rank-revealing
diagnosis. Three quantities follow from the singular values that
no other factorisation gives directly:

\begin{itemize}[leftmargin=*]
  \item The \engterm{numerical rank} of $\mat{A}$, defined as the
    number of singular values above a tolerance (typically
    $\sigma_{1} \cdot u \cdot \max(m, n)$).
  \item The \engterm{2-norm condition number} $\kappa_{2}(\mat{A})
    = \sigma_{1}/\sigma_{r}$, where $r$ is the numerical rank.
  \item The \engterm{best low-rank approximation} of $\mat{A}$ in
    the $\norm{\cdot}_{2}$ and $\norm{\cdot}_{F}$ norms, which is
    obtained by truncating the SVD to its $k$ leading singular
    values (the Eckart-Young theorem).
\end{itemize}

The truncated SVD is the engineer's tool for regularising
ill-conditioned least-squares problems and for compressing
high-rank approximate matrices. Section 16.7 returns to its role
in reproducibility.

\subsection{Iterative methods (recognition level)}

Direct factorisations cost $O(n^{3})$ flops and $O(n^{2})$
storage, prohibitive when $n$ exceeds about $10^{4}$ on a single
node, or when $\mat{A}$ is sparse and the engineer wants to
preserve the sparsity. \engterm{Iterative methods} produce a
sequence of approximate solutions $\vect{x}_{k}$ that converges to
$\vect{x}^{*}$, using $\mat{A}$ only through matrix-vector
products. Three families recur in engineering practice:

\begin{itemize}[leftmargin=*]
  \item \engterm{Stationary methods} (Jacobi, Gauss-Seidel, SOR)
    iterate $\vect{x}_{k+1} = \vect{x}_{k} + \mat{M}^{-1}(\vect{b}
    - \mat{A}\vect{x}_{k})$ for an easily invertible
    \emph{preconditioner} $\mat{M}$. They converge linearly and
    are simple to implement; they are increasingly displaced by
    Krylov methods.
  \item \engterm{Krylov methods} (conjugate gradients for
    symmetric positive definite systems, GMRES for general
    systems, BiCGSTAB and MINRES for special cases) build the
    iterate in the Krylov subspace $\spn\{\vect{r}_{0},
    \mat{A}\vect{r}_{0}, \mat{A}^{2}\vect{r}_{0}, \ldots\}$ and
    converge in at most $n$ steps in exact arithmetic. With a
    good preconditioner they converge in far fewer.
  \item \engterm{Multigrid} methods are tailored to systems
    arising from the discretisation of PDEs, exploiting the
    multi-scale structure of the residual. They achieve $O(n)$
    work per solution for problems where direct methods cost
    $O(n^{3})$.
\end{itemize}

The engineer's working role is recognition: when $n$ is large and
$\mat{A}$ is sparse, the choice between direct and iterative is a
project decision, with trade-offs in time, memory, robustness,
and the engineer's familiarity with the convergence theory.
Volume IV Chapter 8 (computational continuum mechanics) and
Volume IX develop the iterative apparatus at the level needed for
PDE discretisations.

\begin{mastery}
\textbf{Conjugate gradients and the role of the condition number.}
For a symmetric positive definite system $\mat{A}\vect{x} =
\vect{b}$, the conjugate-gradient method's error in the
$\mat{A}$-norm after $k$ iterations is bounded by

\[
\frac{\norm{\vect{x}_{k} - \vect{x}^{*}}_{\mat{A}}}
{\norm{\vect{x}_{0} - \vect{x}^{*}}_{\mat{A}}}
\leq 2 \left(\frac{\sqrt{\kappa} - 1}{\sqrt{\kappa} + 1}\right)^{k},
\]

where $\kappa = \kappa_{2}(\mat{A})$. The bound says the number of
iterations to a fixed relative accuracy scales as $\sqrt{\kappa}$,
where the stationary methods (Jacobi, Gauss-Seidel) scale as
$\kappa$ itself. A problem with $\kappa = 10^{6}$ needs of order
$10^{3}$ conjugate-gradient iterations and of order $10^{6}$
Gauss-Seidel iterations to the same accuracy, the difference between
a routine and an intractable solve. This is why \engterm{preconditioning},
replacing $\mat{A}\vect{x} = \vect{b}$ with $\mat{M}^{-1}\mat{A}
\vect{x} = \mat{M}^{-1}\vect{b}$ for an $\mat{M}$ that clusters the
spectrum, dominates large-scale iterative practice: a good
preconditioner brings the effective $\kappa$ down to a small
constant and the iteration count down with it. The art of
large-sparse-system solving is largely the art of choosing $\mat{M}$,
and it is problem-specific: incomplete-LU, algebraic multigrid, and
domain-decomposition preconditioners each suit a different class of
$\mat{A}$. The condition number, introduced in section 16.2 as the
amplifier of input error, reappears here as the governor of
iteration count, the same quantity wearing a second hat.
\end{mastery}

\begin{example}
\textbf{An iterative solve carried by hand.} Solve

\[
\mat{A}\vect{x} = \vect{b}, \qquad
\mat{A} = \begin{psmallmatrix} 4 & 1 \\ 1 & 3 \end{psmallmatrix},
\qquad \vect{b} = \begin{psmallmatrix} 1 \\ 2 \end{psmallmatrix},
\]

whose exact solution is $\vect{x}^{*} = (1/11, 7/11)^{\transpose} =
(0.0909, 0.6364)^{\transpose}$. The matrix is symmetric positive definite
with eigenvalues $\tfrac{1}{2}(7 \pm \sqrt{5}) = 4.618, 2.382$, so
$\kappa_{2} = 1.94$, a benign system that lets the iteration counts stay
readable. \emph{Jacobi} updates both components from the previous iterate:
$x_{1}^{(k+1)} = (1 - x_{2}^{(k)})/4$, $x_{2}^{(k+1)} = (2 -
x_{1}^{(k)})/3$. From $\vect{x}^{(0)} = (0, 0)$: $\vect{x}^{(1)} =
(0.25, 0.6667)$; $\vect{x}^{(2)} = (0.0833, 0.5833)$; $\vect{x}^{(3)} =
(0.1042, 0.6389)$; $\vect{x}^{(4)} = (0.0903, 0.6319)$; the error norm
falls by a factor near $0.29$ per step (the spectral radius of the Jacobi
iteration matrix is $1/\sqrt{12} = 0.289$), reaching $10^{-6}$ in about
eleven steps. \emph{Gauss-Seidel} uses each freshly computed component
immediately: $x_{1}^{(k+1)} = (1 - x_{2}^{(k)})/4$ then $x_{2}^{(k+1)} =
(2 - x_{1}^{(k+1)})/3$. From $(0, 0)$: $\vect{x}^{(1)} = (0.25, 0.5833)$;
$\vect{x}^{(2)} = (0.1042, 0.6319)$; $\vect{x}^{(3)} = (0.0920, 0.6360)$;
$\vect{x}^{(4)} = (0.0910, 0.6363)$; the per-step factor is the square of
Jacobi's, near $0.083$, and the iteration reaches $10^{-6}$ in about six
steps. The immediate reuse of the updated first component is the whole of
Gauss-Seidel's advantage, and it costs nothing extra. \emph{Conjugate
gradients} on the same system, starting from $\vect{x}^{(0)} = (0,0)$ with
residual $\vect{r}_{0} = \vect{b} = (1, 2)$, takes the steepest first step
along $\vect{r}_{0}$, then a single $\mat{A}$-conjugate correction, and
reaches the \emph{exact} solution in two steps, because CG on an $n \times
n$ SPD system terminates in at most $n$ steps in exact arithmetic. The
three methods rank by the structure they exploit: Jacobi uses only the
diagonal; Gauss-Seidel uses the lower triangle as it goes; CG uses the
full $\mat{A}$-geometry of the Krylov subspace. The ranking that this
$2 \times 2$ shows at a glance is the ranking that governs the
million-unknown sparse systems of Volume IV and Volume IX, where the
constant factors are the same and only $n$ and $\kappa$ change.
\end{example}

\begin{figure}[ht]
\centering
\begin{tikzpicture}[x=0.85cm, y=0.5cm, font=\footnotesize]

% Semilog: iteration count (x) vs log10 relative residual (y).
% Gauss-Seidel converges linearly with rate ~ (1 - 1/kappa); slow.
% CG converges with rate ~ (sqrt(kappa)-1)/(sqrt(kappa)+1); faster,
% and terminates by step n in exact arithmetic. Curves drawn for a
% symmetric positive definite system with kappa ~ 100.

% Axes
\draw[->, thick] (0,0) -- (13,0) node[right] {iteration $k$};
\draw[->, thick] (0,-12) -- (0,1) node[above, align=center]
  {$\log_{10}\dfrac{\norm{\vect{r}_{k}}}{\norm{\vect{r}_{0}}}$};

\foreach \k in {0,2,4,6,8,10,12} {
  \draw (\k,0) -- (\k,0.25);
  \node[anchor=south, font=\scriptsize] at (\k,0.25) {\k};
}
\foreach \d in {0,-3,-6,-9,-12} {
  \draw (0,\d) -- (-0.2,\d);
  \node[anchor=east, font=\scriptsize] at (-0.2,\d) {$10^{\d}$};
}

% Gauss-Seidel: linear, rate ~ 0.82 per step (slope ~ -0.086 dec/step)
% drawn over the full window, reaching only ~10^{-1.1} by step 12.
\draw[thick, engred] plot coordinates {
  (0,0) (2,-0.18) (4,-0.36) (6,-0.55) (8,-0.74) (10,-0.93) (12,-1.12)
};
\node[engred, font=\scriptsize, anchor=west] at (7.2,-0.55)
  {Gauss-Seidel};

% Conjugate gradients: faster geometric drop, hitting 10^{-11} by step 12
\draw[thick, engblue] plot coordinates {
  (0,0) (1,-0.6) (2,-1.5) (3,-2.6) (4,-3.8) (5,-5.0) (6,-6.1)
  (7,-7.2) (8,-8.2) (9,-9.1) (10,-9.9) (11,-10.6) (12,-11.2)
};
\node[engblue, font=\scriptsize, anchor=west] at (5.4,-4.4)
  {conjugate gradients};

\end{tikzpicture}
\caption[Gauss-Seidel versus conjugate gradients on an SPD system]{Relative
residual norm against iteration count for a symmetric positive definite
linear system of condition number $\kappa \approx 100$, comparing the
stationary Gauss-Seidel iteration with the Krylov conjugate-gradient (CG)
method. Gauss-Seidel contracts the residual by a near-constant factor per
step (linear convergence with asymptotic rate governed by $1 - O(1/\kappa)$),
losing roughly one decimal digit over twelve iterations. CG contracts at
the faster rate $(\sqrt{\kappa}-1)/(\sqrt{\kappa}+1)$, reaching eleven
digits in the same twelve iterations, and would terminate exactly by step
$n$ in exact arithmetic. The gap widens with $\kappa$: the iteration count
to a fixed accuracy scales as $\kappa$ for the stationary method and as
$\sqrt{\kappa}$ for CG. Source: authors' analysis; the convergence bounds
are standard, see \cite[ch.~38]{text:trefethen-bau1997}.}
\label{fig:vol02:ch16:cg-vs-gs}
\end{figure}


Figure~\ref{fig:vol02:ch16:cg-vs-gs} carries the same comparison to a
worse-conditioned system, where the gap between the stationary method and
the Krylov method is the difference between a tractable and an intractable
solve. The stationary iteration's residual falls by a near-constant factor
per step set by $1 - O(1/\kappa)$; the conjugate-gradient residual falls
at the rate $(\sqrt{\kappa} - 1)/(\sqrt{\kappa} + 1)$, so the iteration
count to a fixed accuracy scales as $\kappa$ for the stationary method and
as $\sqrt{\kappa}$ for CG. The square-root in the exponent is the entire
reason Krylov methods, with a preconditioner that lowers the effective
$\kappa$, dominate large-scale linear-system practice.

\subsection{GMRES for nonsymmetric systems}

Conjugate gradients exploits symmetric positive definiteness, and most
matrices an engineer meets in a discretised convection problem or a
circuit-network solve are neither symmetric nor definite. The
\engterm{generalised minimal residual} method (GMRES) is the Krylov
method for the general nonsingular $\mat{A}$. At step $k$ it builds an
orthonormal basis $\mat{Q}_{k} = [\vect{q}_{1}, \ldots, \vect{q}_{k}]$ of
the Krylov subspace $\mathcal{K}_{k} = \spn\{\vect{r}_{0},
\mat{A}\vect{r}_{0}, \ldots, \mat{A}^{k-1}\vect{r}_{0}\}$ by the Arnoldi
process, and chooses the iterate $\vect{x}_{k} \in \vect{x}_{0} +
\mathcal{K}_{k}$ that minimises the residual norm $\norm{\vect{b} -
\mat{A}\vect{x}_{k}}_{2}$ over that subspace
\cite[ch.~35]{text:trefethen-bau1997}. The minimisation reduces to a small
$(k+1) \times k$ least-squares problem on the upper-Hessenberg matrix
$\mat{H}_{k}$ that the Arnoldi process produces, solved cheaply by Givens
rotations that are updated one column per step.

GMRES carries two costs that CG does not. The Arnoldi orthogonalisation
stores all $k$ basis vectors and costs $O(k)$ inner products at step $k$,
so the work and the memory grow with the iteration count. Production GMRES
therefore \engterm{restarts}: after $m$ steps (a typical $m$ is $20$ to
$50$) it discards the basis, takes the current iterate as a new $\vect{x}_{0}$,
and begins a fresh cycle, written GMRES($m$). The restart caps the memory
at $m$ vectors and the per-step cost at $O(m)$, at the price of slower
convergence, since the restarted method forgets the subspace it had built.

\begin{example}
\textbf{GMRES on a $3 \times 3$ nonsymmetric system, by hand.} Solve
$\mat{A}\vect{x} = \vect{b}$ with

\[
\mat{A} = \begin{psmallmatrix} 2 & 1 & 0 \\ 0 & 2 & 1 \\ 1 & 0 & 2
\end{psmallmatrix}, \qquad
\vect{b} = \begin{psmallmatrix} 1 \\ 0 \\ 0 \end{psmallmatrix},
\]

a matrix that is nonsymmetric (so CG does not apply) and nonsingular (its
determinant is $9$). Start from $\vect{x}_{0} = \vect{0}$, so $\vect{r}_{0}
= \vect{b} = (1, 0, 0)^{\transpose}$, $\norm{\vect{r}_{0}} = 1$, and
$\vect{q}_{1} = \vect{r}_{0}$. The first Arnoldi step forms $\mat{A}
\vect{q}_{1} = (2, 0, 1)^{\transpose}$; the Hessenberg entry is $h_{11} =
\vect{q}_{1}^{\transpose}\mat{A}\vect{q}_{1} = 2$; the orthogonal residual
is $\vect{w} = \mat{A}\vect{q}_{1} - h_{11}\vect{q}_{1} = (0, 0,
1)^{\transpose}$, so $h_{21} = \norm{\vect{w}} = 1$ and $\vect{q}_{2} =
(0, 0, 1)^{\transpose}$. The least-squares problem at $k = 1$ minimises
$\norm{\,\norm{\vect{r}_{0}}\vect{e}_{1} - \bar{\mat{H}}_{1} y\,}$ with
$\bar{\mat{H}}_{1} = (2, 1)^{\transpose}$ and target $(1, 0)^{\transpose}$,
giving $y_{1} = (2)(1)/(2^{2} + 1^{2}) = 0.4$ and residual norm
$\abs{1 - 2 \cdot 0.4} / \sqrt{1} \cdot \ldots$, concretely
$\norm{(1, 0)^{\transpose} - (2, 1)^{\transpose}(0.4)} = \norm{(0.2,
-0.4)^{\transpose}} = 0.447$. The residual has fallen from $1$ to $0.447$
in one step, exactly the monotone decrease GMRES guarantees. The second
Arnoldi step forms $\mat{A}\vect{q}_{2} = (0, 1, 2)^{\transpose}$,
orthogonalises against $\vect{q}_{1}$ and $\vect{q}_{2}$ to extend
$\bar{\mat{H}}$, and the $k = 2$ least-squares solve drops the residual
again; the third step closes the Krylov subspace (it now spans $\R^{3}$)
and GMRES returns the exact $\vect{x} = \mat{A}^{-1}\vect{b} =
(0.444, 0.111, -0.222)^{\transpose}$ at step $3$, the finite-termination
property every Krylov method has on an $n$-dimensional space. The example
exposes the machinery the library hides: GMRES is least-squares over a
growing orthonormal Krylov basis, the residual decreases monotonically by
construction, and the work per step is the Arnoldi orthogonalisation whose
cost is the reason the method restarts on large problems.
\end{example}

\begin{example}
\textbf{Preconditioning changes the iteration count, not the answer.} The
convergence of GMRES is governed not by $\kappa$ alone but by how the
spectrum of $\mat{A}$ clusters. A matrix whose eigenvalues lie in a tight
cluster away from the origin converges in a handful of steps; a matrix with
eigenvalues scattered across a wide region, or straddling the origin,
converges slowly. \engterm{Preconditioning} replaces $\mat{A}\vect{x} =
\vect{b}$ with $\mat{M}^{-1}\mat{A}\vect{x} = \mat{M}^{-1}\vect{b}$ for an
$\mat{M}$ cheap to invert and close to $\mat{A}$, so that $\mat{M}^{-1}
\mat{A} \approx \mat{I}$ has a clustered spectrum. Take the
diagonally dominant $\mat{A} = \mat{D} + \mat{E}$ with $\mat{D} = 100\,
\mat{I}$ and $\mat{E}$ a random matrix of entries near $1$. Unpreconditioned,
the eigenvalues cluster near $100$ already and GMRES converges in a few
steps. Take instead a matrix with $\mat{D}$ spanning $\{1, 10, 100, 1000,
\ldots\}$ on the diagonal: now the spectrum spans four decades, GMRES
crawls, and the Jacobi (diagonal) preconditioner $\mat{M} = \mat{D}$ maps
every diagonal entry to $1$, clustering the preconditioned spectrum near
unity and cutting the iteration count by an order of magnitude. The
preconditioned and unpreconditioned solves return the same
$\vect{x}$ to working precision; the preconditioner buys iteration count,
not accuracy. The art of large-scale solving, as the conjugate-gradient
mastery box stated, is the art of choosing $\mat{M}$: incomplete-LU for
general sparse matrices, algebraic multigrid for elliptic-PDE
discretisations, and block-Jacobi or additive-Schwarz for domain-decomposed
problems, each matched to the structure of $\mat{A}$.
\end{example}

\subsection{Iterative refinement and mixed precision}

A backward-stable direct solve returns $\hat{\vect{x}}$ whose forward error
is bounded by $\kappa(\mat{A}) u$. When that bound costs the engineer more
digits than the application allows, \engterm{iterative refinement} recovers
them without re-factorising. Given the computed $\hat{\vect{x}}$ from one
LU solve, the refinement loop is: compute the residual $\vect{r} =
\vect{b} - \mat{A}\hat{\vect{x}}$; solve $\mat{A}\vect{d} = \vect{r}$ reusing
the existing $\mat{L}$ and $\mat{U}$; update $\hat{\vect{x}} \leftarrow
\hat{\vect{x}} + \vect{d}$; repeat. Each pass costs an $O(n^{2})$ residual
and triangular solve against the $O(n^{3})$ of the original factorisation,
so refinement is cheap. The crucial detail is the precision of the
residual: when $\vect{r}$ is computed in higher precision than the working
factorisation, the loop drives the forward error down to working-precision
accuracy even for condition numbers up to $1/u$
\cite[ch.~12]{text:higham-numerical}.

The modern form of this idea drives the fastest dense solvers in production.
Graphics and machine-learning hardware computes low-precision matrix
products (binary16 or bfloat16) at several times the throughput of binary64.
\engterm{Mixed-precision iterative refinement} factorises $\mat{A}$ in
binary16 or binary32 at the high low-precision throughput, then refines in
a higher precision (sometimes three precisions: low for the factorisation,
working for the update, high for the residual), recovering binary64 accuracy
on well-conditioned problems while paying the $O(n^{3})$ cost at
low-precision speed. The technique reached the top of the high-performance
computing benchmarks in the late 2010s and is the reason a tensor-core GPU
solves a dense linear system far faster than its binary64 peak flop rate
would predict \cite[ch.~12]{text:higham-numerical}.

\begin{example}
\textbf{One refinement pass recovers digits.} Solve $\mat{A}\vect{x} =
\vect{b}$ for a matrix with $\kappa(\mat{A}) = 10^{8}$ in binary64. The
plain LU solve returns $\hat{\vect{x}}$ with relative forward error near
$\kappa u \approx 10^{8} \cdot 10^{-16} = 10^{-8}$, eight correct digits.
Iterative refinement with a residual computed in the same binary64 working
precision typically halves the exponent of the error per pass when
$\kappa u < 1$, so two or three passes bring the error toward $10^{-15}$,
full working accuracy, at the cost of a few $O(n^{2})$ residual-and-solve
steps. The same loop on a matrix with $\kappa = 10^{17} > 1/u$ stalls: the
residual itself cannot be computed to enough relative accuracy in working
precision, and refinement needs an extended-precision residual to make
progress. The diagnostic the engineer reads from this is the same separation
the chapter rests on. Refinement attacks the algorithm's contribution to
the error (it cleans up round-off in the solve), and it works while the
problem still carries recoverable digits ($\kappa u < 1$); once the problem
itself has destroyed the digits ($\kappa u \geq 1$), no working-precision
loop recovers them and the engineer must raise the residual precision or
reformulate.
\end{example}

\section{Numerical integration revisited}\pathtag{standard}

Volume II Chapter 6 introduced the trapezoidal rule and Simpson's
rule. This section completes the apparatus.

\subsection{Trapezoidal and Simpson on a uniform grid}

For a smooth function $f$ on $[a, b]$, divided into $n$ equal
subintervals of width $h = (b - a)/n$, the \engterm{composite
trapezoidal rule} approximates

\[
\int_{a}^{b} f(x)\, \dd x \approx \frac{h}{2}\!\left[f(x_{0}) +
2 \sum_{k=1}^{n-1} f(x_{k}) + f(x_{n})\right],
\]

with truncation error $-\frac{(b-a)h^{2}}{12} f''(\xi)$ for some
$\xi \in (a, b)$. The error is $O(h^{2})$. The \engterm{composite
Simpson's rule} (requiring $n$ even) approximates

\[
\int_{a}^{b} f(x)\, \dd x \approx \frac{h}{3}\!\left[f(x_{0}) +
4 \sum_{k\, \mathrm{odd}} f(x_{k}) + 2 \sum_{k\, \mathrm{even}}
f(x_{k}) + f(x_{n})\right],
\]

with truncation error $-\frac{(b-a)h^{4}}{180} f^{(4)}(\xi)$, of
order $O(h^{4})$. Each doubling of $n$ reduces the error by 4 for
trapezoidal and by 16 for Simpson, when the integrand has the
required smoothness.

\subsection{Gaussian quadrature}

\engterm{Gaussian quadrature} chooses both the nodes $x_{i}$ and
the weights $w_{i}$ so that the rule

\[
\int_{-1}^{1} f(x)\, \dd x \approx \sum_{i=1}^{n} w_{i} f(x_{i})
\]

is exact for polynomials of degree up to $2n - 1$. The optimal
nodes are the roots of the degree-$n$ Legendre polynomial; the
weights are determined by the integration of the Lagrange
interpolating polynomials at those nodes. The truncation error
for an integrand $f$ with $2n$ continuous derivatives is bounded
by a quantity proportional to $\max\abs*{f^{(2n)}}$ over $[-1,
1]$, with a small constant
\cite[ch.~10]{text:moler-numerical}.

The intervals $[a, b]$ other than $[-1, 1]$ are handled by the
linear change of variables $x = \frac{a+b}{2} + \frac{b-a}{2} t$,
which scales the weights by $(b-a)/2$. Variants exist for
weighted integrals over $[0, \infty)$ (Gauss-Laguerre), over
$(-\infty, \infty)$ (Gauss-Hermite), and over $[-1, 1]$ with
endpoint singularities (Gauss-Chebyshev). The engineer chooses
the variant whose weight matches the integrand's singular
structure.

For smooth integrands, Gauss with $n$ nodes is dramatically more
accurate than trapezoidal or Simpson with comparable cost. For
integrands with discontinuities, weak singularities, or rapid
oscillation, fixed-order Gaussian quadrature is no better than
the lower-order rules and is often worse.

\begin{example}
\textbf{Three nodes against eight.} Evaluate $\int_{0}^{1} e^{x}\,
\dd x = e - 1 = 1.718282$ by three-point Gauss-Legendre and compare
to the composite Simpson rule. Mapping $[0, 1]$ to $[-1, 1]$ by
$x = (t + 1)/2$, the three Gauss nodes $t_{i} = -\sqrt{3/5}, 0,
+\sqrt{3/5}$ map to $x_{i} = 0.1127, 0.5, 0.8873$ with weights
$w_{i} = \tfrac{1}{2}(5/9, 8/9, 5/9) = 0.2778, 0.4444, 0.2778$ after
the $(b-a)/2 = \tfrac{1}{2}$ scaling. The Gauss estimate is
$0.2778\, e^{0.1127} + 0.4444\, e^{0.5} + 0.2778\, e^{0.8873} =
0.3109 + 0.7328 + 0.6746 = 1.718281$, in error by about $7 \times
10^{-7}$ with only three function evaluations. Composite Simpson
matches that accuracy only near eight subintervals (nine
evaluations). The three-point Gauss rule is exact for polynomials
through degree $5$, where three-point Simpson is exact only through
degree $3$, and the extra two degrees buy the two-decade accuracy
gain at lower cost. The lesson is the one the section states: for
smooth integrands, accuracy per evaluation favours Gauss heavily;
the engineer reserves the equispaced rules for tabulated data whose
abscissae are not free to choose.
\end{example}

\begin{example}
\textbf{Where the Gauss nodes come from.} The two-point Gauss-Legendre
rule on $[-1, 1]$ has four free parameters, the nodes $x_{1}, x_{2}$ and
the weights $w_{1}, w_{2}$, so it can be made exact for the four monomials
$1, x, x^{2}, x^{3}$, that is for every polynomial through degree $3$. The
exactness conditions are the moment equations $\sum_{i} w_{i} x_{i}^{k} =
\int_{-1}^{1} x^{k}\, \dd x$ for $k = 0, 1, 2, 3$, namely $w_{1} + w_{2} =
2$, $w_{1}x_{1} + w_{2}x_{2} = 0$, $w_{1}x_{1}^{2} + w_{2}x_{2}^{2} =
2/3$, and $w_{1}x_{1}^{3} + w_{2}x_{2}^{3} = 0$. Symmetry of the interval
suggests $x_{2} = -x_{1}$ and $w_{1} = w_{2}$; the first equation then gives
$w_{1} = w_{2} = 1$, the second and fourth are satisfied automatically, and
the third gives $2 x_{1}^{2} = 2/3$, so $x_{1} = 1/\sqrt{3}$. The two-point
rule is therefore $\int_{-1}^{1} f \approx f(-1/\sqrt{3}) + f(1/\sqrt{3})$,
exact through cubics with no weights to remember. The same nodes arise a
second way, which is the route that scales: $\pm 1/\sqrt{3}$ are the roots
of the degree-two Legendre polynomial $P_{2}(x) = \tfrac{1}{2}(3x^{2} - 1)$.
That the Gauss nodes are the roots of the orthogonal polynomial is the
general theorem. A rule whose nodes are the roots of the degree-$n$ member
of the family orthogonal under the integral's weight integrates the node
polynomial $\prod(x - x_{i})$ against any lower-degree polynomial to zero,
which is exactly the extra $n$ degrees of exactness that lift an $n$-node
rule from degree $n - 1$ (generic) to degree $2n - 1$ (Gauss). The
construction generalises by changing the orthogonal family to match the
weight: Laguerre for $e^{-x}$ on $[0, \infty)$, Hermite for $e^{-x^{2}}$ on
$\R$, Chebyshev for $(1 - x^{2})^{-1/2}$ on $[-1, 1]$.
\end{example}

\begin{mastery}
\textbf{Romberg integration and Richardson extrapolation.} The
composite trapezoidal rule has an error expansion in even powers of
$h$ (the Euler-Maclaurin formula): $T(h) = I + c_{1} h^{2} +
c_{2} h^{4} + c_{3} h^{6} + \cdots$ for a smooth integrand. Two
trapezoidal estimates at $h$ and $h/2$ can then be combined to
cancel the leading $h^{2}$ term:

\[
R(h) = \frac{4 T(h/2) - T(h)}{3} = I + O(h^{4}),
\]

which is exactly Simpson's rule. Repeating the cancellation on the
$h^{4}$, $h^{6}$, terms builds the \engterm{Romberg table}, each
column one Richardson-extrapolation step higher in order than the
last, reaching $O(h^{2m})$ accuracy from $m$ levels of trapezoidal
halving. Romberg is the cheapest route to high-order accuracy on a
smooth integrand when the abscissae must be equispaced (so Gauss is
unavailable), because it reuses every function evaluation across the
columns. The method fails silently on integrands without the smooth
even-power error expansion (those with endpoint singularities or
interior corners), where the extrapolation amplifies rather than
cancels the error; the adaptive Gauss-Kronrod routines of the next
subsection handle those integrands where Romberg cannot.
\end{mastery}

\subsection{The round-off floor of a refined grid}

Composite trapezoidal and Simpson rules have truncation error that
falls as $h^{2}$ and $h^{4}$ respectively. The round-off error of
the sum scales as $u \sqrt{n}$ (for $n$ subintervals, in the
statistical bound). Halving $h$ doubles $n$ and brings round-off up
by $\sqrt{2}$. The total error has a minimum at the $h$ where the
truncation error has fallen to the same order as the round-off
floor; refining the grid further is wasted compute that adds
round-off without lowering truncation.

The same pattern is sharper for the finite-difference derivative,
where round-off enters as $u/h$. Figure~\ref{fig:vol02:ch16:error-vs-h}
plots the total error of a centred-difference derivative on log-log
axes. The truncation slope is $+2$; the round-off slope is $-1$; the
total error is the upper envelope and has a minimum near
$h^{*} \approx u^{1/3} \approx 10^{-5.3}$ in binary64, with attainable
error near $10^{-10.6}$. The engineer who chases truncation error
past this point by halving $h$ further loses digits to round-off and
reports a worse derivative than at the optimum.

\begin{figure}[ht]
\centering
\begin{tikzpicture}[x=0.6cm, y=0.45cm, font=\footnotesize]

% Log-log: x = log10(h), y = log10(|error|)
% h from 10^0 down to 10^{-15}
% Truncation error decreases as h^p; round-off increases as 1/h for
% finite-difference derivative; total has a U shape with minimum near
% h ~ sqrt(u) ~ 10^{-8} for first-derivative central difference.

% Axes
\draw[->, thick] (-15.5,0) -- (1,0) node[right] {$\log_{10} h$};
\draw[->, thick] (-15.5,-16) -- (-15.5,1) node[above, align=center]
  {$\log_{10}\abs{\text{error}}$};

% Tick marks
\foreach \k in {0,-3,-6,-9,-12,-15} {
  \draw (\k,0) -- (\k,0.3);
  \node[anchor=south, font=\scriptsize] at (\k,0.3) {$10^{\k}$};
}
\foreach \d in {0,-4,-8,-12,-16} {
  \draw (-15.5,\d) -- (-15.8,\d);
  \node[anchor=east, font=\scriptsize] at (-15.8,\d) {$10^{\d}$};
}

% Truncation error: O(h^2) for central difference, slope +2
% y = 2*x - 2 (so at h=10^0 error = 10^{-2}; at h = 10^{-8} error = 10^{-18})
\draw[thick, engblue] (0,-2) -- (-8,-18);
\node[engblue, font=\scriptsize, anchor=west] at (-3,-7)
  {truncation $\sim h^{2}$};

% Round-off: ~ u / h, slope -1 (here y = -x - 16; at h=10^0 error = 10^{-16}, at h=10^{-15} error = 10^{-1})
\draw[thick, engred] (0,-16) -- (-15,-1);
\node[engred, font=\scriptsize, anchor=east] at (-12,-4)
  {round-off $\sim u/h$};

% Total error: pointwise max of the two
\draw[thick, black, dashed] plot coordinates {
  (0,-2) (-1,-4) (-2,-6) (-3,-8) (-4,-10) (-5,-12) (-5.5,-12.5)
  (-5.8,-12.7) (-6,-12.4) (-7,-11.8) (-8,-10.7) (-10,-8.5)
  (-12,-6.5) (-15,-1)
};
\node[font=\scriptsize, anchor=south] at (-6,-12.5)
  {optimal $h^{*} \approx u^{1/3}$};

% Mark optimum (h ~ 10^{-5.3}, error ~ 10^{-10.6}) for central diff
\fill[black] (-5.5,-12.5) circle (2pt);

\end{tikzpicture}
\caption[Truncation vs round-off in finite-difference differentiation]{The
total error of a finite-difference approximation to a derivative as a
function of step size $h$, on log-log axes. The truncation error of the
centred difference $f'(x) \approx (f(x + h) - f(x - h))/(2 h)$ scales
as $h^{2}$ and decreases with $h$. The round-off error of the
subtraction in the numerator scales as $u / h$ and increases as $h$
shrinks. The total error has a minimum near $h^{*} \approx u^{1/3}$,
with $u$ the unit round-off; for binary64, $h^{*} \approx 10^{-5.3}$,
producing an attainable error near $10^{-10.6}$. A naive engineer who
chases truncation error down past this point by shrinking $h$ further
loses digits to round-off and reports a worse derivative than at the
optimum. Source: authors' analysis; the $u^{1/3}$ optimum is standard,
see \cite[\S 5.7]{text:higham-numerical}.}
\label{fig:vol02:ch16:error-vs-h}
\end{figure}


The optimal step is algorithm-dependent: $h^{*} \sim u^{1/2}$ for
the forward difference, $u^{1/3}$ for the centred difference, $u^{1/4}$
for a Richardson-extrapolated centred difference; the attainable
relative accuracy is the algorithm's order divided by one plus that
order \cite[\S 5.7]{text:higham-numerical}. The reader who never
needs to write a finite-difference derivative by hand still inherits
the same trade-off when calling a black-box adaptive routine: the
tolerance the engineer requests must lie above the round-off floor
of the algorithm and the format, and a request below that floor
returns an answer the engineer cannot defend.

\begin{mastery}
\textbf{Automatic differentiation defeats the round-off floor.} The
finite-difference derivative is trapped between truncation and round-off:
the optimal step $h^{*} \sim u^{1/3}$ for the centred difference leaves an
unavoidable error near $u^{2/3} \approx 10^{-10.6}$ in binary64, far short
of the arithmetic's full precision. \engterm{Automatic differentiation}
(AD) computes derivatives to full working precision with no step size at
all, by propagating derivative information through the program's elementary
operations rather than perturbing the input. \engterm{Forward-mode} AD
carries each value as a \engterm{dual number} $a + b\varepsilon$ with
$\varepsilon^{2} = 0$, defining arithmetic so that $(a + b\varepsilon)(c +
d\varepsilon) = ac + (ad + bc)\varepsilon$ and, by the chain rule, every
elementary function $\phi$ acts as $\phi(a + b\varepsilon) = \phi(a) +
b\,\phi'(a)\varepsilon$. Evaluating $f$ at $x + 1\cdot\varepsilon$ then
returns $f(x) + f'(x)\varepsilon$, and the $\varepsilon$-component is the
exact derivative of the program, accurate to round-off and free of any
truncation. The cost of one forward-mode derivative is a small constant
times the cost of evaluating $f$. \engterm{Reverse-mode} AD (the
backpropagation of machine learning) computes the gradient of a scalar
output with respect to all inputs at a cost independent of the input
dimension, which is why training a network with millions of parameters is
feasible at all; Volume VII develops it for that use. The durable point
for this chapter is the contrast with finite differences: AD is exact up
to round-off because it differentiates the program symbolically at each
operation while evaluating numerically, where the finite difference
approximates the limit and pays the truncation-round-off trade of
Figure~\ref{fig:vol02:ch16:error-vs-h}. The complex-step derivative
$f'(x) \approx \Im[f(x + i h)]/h$ is a related trick that also escapes the
subtractive cancellation (there is no subtraction of nearly equal reals),
delivering full precision for analytic real functions at $h$ as small as
$10^{-200}$ \cite[\S 5.7]{text:higham-numerical}. The engineer who needs a
derivative inside an optimisation or a Newton solve reaches for AD or the
complex step before the finite difference, and accepts the finite
difference only when the function is a black box that admits neither.
\end{mastery}

\subsection{Adaptive quadrature}

Real engineering integrands are not uniformly smooth. An
\engterm{adaptive quadrature} routine subdivides the interval
where the local error estimate is large and accepts the
approximation where the local error estimate is small. The
typical pattern: estimate the integral on $[a, b]$ by two methods
of different order (e.g.\ Simpson and Simpson-on-halves);
estimate the error from the difference; if the error exceeds
tolerance, recurse on the two halves; otherwise return the
higher-order estimate.

The \texttt{scipy.integrate.quad}, MATLAB \texttt{integral}, and
GSL \texttt{gsl\_integration\_qag} routines are adaptive and
backed by Gauss-Kronrod pairs (a Gauss rule and a Kronrod
extension that share nodes, allowing efficient error estimation).
The engineer's working role is to feed a sensible interval, a
sensible tolerance, and to read the returned error estimate
seriously when the integrand has structure (a peak, a corner, a
log-singular endpoint).

\begin{figure}[ht]
\centering
\begin{tikzpicture}[x=1.0cm, y=1.6cm, font=\footnotesize]

% A peaked integrand on [0,1] with the recursive bisection pattern of
% adaptive Simpson drawn underneath: the interval is subdivided densely
% where the integrand is curved (near the peak) and left coarse where
% the integrand is flat. Tick depth encodes recursion depth.

% the curve: a Lorentzian-like peak centred at x=0.7
% f(x) = 1 / (1 + 400 (x-0.7)^2 ), scaled to fit
\draw[->, thick] (-0.2,0) -- (10.6,0) node[right] {$x$};
\draw[->, thick] (0,-0.05) -- (0,1.25) node[above] {$f(x)$};
\node[anchor=north, font=\scriptsize] at (0,-0.05) {$0$};
\node[anchor=north, font=\scriptsize] at (10,-0.05) {$1$};

% peak curve (x in screen units 0..10 maps to [0,1]; peak at screen 7)
\draw[thick, engblue] plot[domain=0:10, samples=120] (\x,
  {1/(1 + 4*(\x-7)*(\x-7))});

% Coarse subdivisions away from peak (depth 1): big intervals
\foreach \xb in {0,2.5} {
  \draw[engred, thick] (\xb,0) -- (\xb,-0.14);
}
% Medium subdivisions approaching peak (depth 2)
\foreach \xb in {5.0,6.0} {
  \draw[engred, thick] (\xb,0) -- (\xb,-0.20);
}
% Fine subdivisions at the peak (depth 3..4)
\foreach \xb in {6.5,6.75,7.0,7.25,7.5} {
  \draw[engred, thick] (\xb,0) -- (\xb,-0.26);
}
% Medium again on the far flank (depth 2)
\foreach \xb in {8.0,9.0} {
  \draw[engred, thick] (\xb,0) -- (\xb,-0.20);
}
% Endpoint
\draw[engred, thick] (10,0) -- (10,-0.14);

\node[engred, font=\scriptsize, anchor=north] at (1.25,-0.34)
  {coarse: flat region};
\node[engred, font=\scriptsize, anchor=north] at (7.0,-0.40)
  {fine: high curvature};

\end{tikzpicture}
\caption[Adaptive-quadrature subdivision following the integrand]{An
adaptive Simpson routine on a peaked integrand. The routine estimates the
local error on each panel by comparing a Simpson estimate to the sum of two
Simpson estimates on the halves, and recurses only where the discrepancy
exceeds the per-panel tolerance. The tick depth marks the recursion depth:
the flat flanks are accepted at coarse resolution after one or two
subdivisions, while the curved peak draws repeated bisection until the local
error falls below tolerance. The total work concentrates where the integrand
has structure, which is the property fixed-order rules lack. The worked trace
in the text carries one such subdivision numerically. Source: authors'
analysis.}
\label{fig:vol02:ch16:adaptive-quadrature}
\end{figure}


Figure~\ref{fig:vol02:ch16:adaptive-quadrature} shows the subdivision
pattern an adaptive routine produces on a peaked integrand: dense near the
peak, coarse on the flat flanks. The next example carries one level of that
recursion numerically so the error estimate and the accept-or-recurse
decision are visible as arithmetic rather than as a library call.

\begin{example}
\textbf{An adaptive Simpson trace.} Integrate $\int_{0}^{1} \frac{1}{1 +
16 x^{2}}\, \dd x$, whose exact value is $\tfrac{1}{4}\arctan 4 =
0.331161$, by adaptive Simpson with a tolerance $\tau = 10^{-3}$ on the
whole interval. The basic Simpson estimate on a panel $[a, b]$ with
midpoint $m = (a+b)/2$ is $S(a,b) = \frac{b-a}{6}(f(a) + 4 f(m) + f(b))$.
On $[0, 1]$: $f(0) = 1$, $f(0.5) = 1/5 = 0.2$, $f(1) = 1/17 = 0.0588$, so
$S(0, 1) = \frac{1}{6}(1 + 0.8 + 0.0588) = 0.30980$. Subdivide at $m =
0.5$ and apply Simpson on each half. Left half $[0, 0.5]$ with midpoint
$0.25$: $f(0.25) = 1/2 = 0.5$, $S(0, 0.5) = \frac{0.5}{6}(1 + 4(0.5) +
0.2) = \frac{0.5}{6}(3.2) = 0.26667$. Right half $[0.5, 1]$ with midpoint
$0.75$: $f(0.75) = 1/10 = 0.1$, $S(0.5, 1) = \frac{0.5}{6}(0.2 + 4(0.1) +
0.0588) = \frac{0.5}{6}(0.6588) = 0.05490$. The two-panel estimate is
$S_{2} = 0.26667 + 0.05490 = 0.32157$. The error estimate is the Richardson
quantity $\abs{S_{2} - S(0,1)}/15 = \abs{0.32157 - 0.30980}/15 = 0.01177/15
= 7.8 \times 10^{-4}$, which exceeds the per-panel tolerance the routine
assigns to $[0,1]$, so the routine \emph{recurses} on both halves. On the
left half the local error estimate after one further subdivision drops to
about $4 \times 10^{-5}$ and the panel is accepted; the steeper right flank
near the curvature of the integrand requires one more level before its
estimate falls below tolerance. The accepted total, $S_{2}$ plus the
Richardson correction $(S_{2} - S(0,1))/15 = 0.32157 + 0.00078 = 0.32235$
on the first level and refined further on recursion, converges to the exact
$0.33116$ as the recursion deepens where the integrand bends. The trace
exposes the two ideas the library hides: the error estimate is a comparison
of two orders on the same nodes (Simpson against Simpson-on-halves, the
$1/15$ from the $h^{4}$ error ratio), and the recursion is driven by that
estimate, not by a fixed grid. The reader who has carried this one level
reads \texttt{scipy.integrate.quad}'s returned error estimate as the same
quantity at machine scale.
\end{example}

\begin{mastery}
\textbf{Gauss-Kronrod node sharing.} The adaptive Simpson trace estimated
its error by halving the panel and comparing two rules of the same family,
which throws away the coarse rule's function values when it refines. The
production adaptive routines (QUADPACK's \texttt{qag}, behind
\texttt{scipy.integrate.quad} and MATLAB \texttt{integral}) instead pair a
Gauss rule with a \engterm{Kronrod extension} that reuses every Gauss node.
A degree-$n$ Gauss rule integrates polynomials through degree $2n - 1$ with
$n$ nodes; the Gauss-Kronrod extension adds $n + 1$ new nodes, interlacing
with the originals, to build a $(2n + 1)$-point rule exact through degree
$3n + 1$, and the $2n + 1$ points include all $n$ original Gauss points, so
the engineer pays only $n + 1$ \emph{new} evaluations to obtain a much
higher-order estimate alongside the original. The difference between the
Gauss and the Gauss-Kronrod values on the same panel is the local error
estimate, computed at $n + 1$ extra function calls rather than the doubling
an order-comparison would cost. The standard QUADPACK pairing is the
$7$-point Gauss with the $15$-point Kronrod, the (G7, K15) pair: fifteen
evaluations deliver a degree-$22$ estimate and a free error indicator. The
extra Kronrod nodes are not arbitrary; they are the roots of the Stieltjes
polynomial orthogonal to the Legendre polynomials under the Legendre weight,
the same orthogonality argument that placed the Gauss nodes, applied one
level up to place the points that extend the rule
\cite[ch.~3]{text:trefethen-bau1997}. The durable reading is that node
sharing is the quadrature analogue of the embedded Runge-Kutta pairs of the
next section: two rules of different order built on overlapping evaluations,
the cheap one a free error estimate for the accurate one. The same
design pattern, an accurate estimate paired with a cheap error indicator on
shared work, recurs in every adaptive numerical method the engineer will
meet.
\end{mastery}

\section{Numerical ODE/PDE solvers}\pathtag{standard}

Volume II Chapter 7 introduced forward Euler and recognised RK4.
This section develops the apparatus that distinguishes solvers
worth running from solvers that will silently lie.

\subsection{Forward and backward Euler}

The initial-value problem $\dot{y} = f(t, y)$, $y(t_{0}) = y_{0}$
on a time grid $t_{n} = t_{0} + n h$ produces, by \engterm{forward
Euler},

\[
y_{n+1} = y_{n} + h f(t_{n}, y_{n}),
\]

a one-step \engterm{explicit} method with local truncation error
$O(h^{2})$ and global error $O(h)$. \engterm{Backward Euler}
solves

\[
y_{n+1} = y_{n} + h f(t_{n+1}, y_{n+1})
\]

for $y_{n+1}$ at each step, an \engterm{implicit} method, also
first-order. The implicit equation is solved by Newton's method
or by a fixed-point iteration; for linear $f$, by a single linear
solve.

The two methods differ sharply in their stability. Applied to the
test equation $\dot{y} = \lambda y$ with $\lambda < 0$, forward
Euler is stable only when $\abs{1 + h\lambda} < 1$, that is when
$h < 2/\abs{\lambda}$; backward Euler is stable for all $h > 0$.
A method whose stability region includes the entire left half of
the complex plane is \engterm{A-stable}; backward Euler is
A-stable, forward Euler is not.

\subsection{Runge-Kutta methods}

A \engterm{Runge-Kutta} method takes one step by combining
several stage evaluations of $f$. The classical fourth-order
\engterm{RK4} is

\begin{align*}
k_{1} &= f(t_{n}, y_{n}), \\
k_{2} &= f(t_{n} + h/2, y_{n} + (h/2) k_{1}), \\
k_{3} &= f(t_{n} + h/2, y_{n} + (h/2) k_{2}), \\
k_{4} &= f(t_{n} + h, y_{n} + h k_{3}), \\
y_{n+1} &= y_{n} + \frac{h}{6}(k_{1} + 2 k_{2} + 2 k_{3} + k_{4}).
\end{align*}

The local truncation error is $O(h^{5})$, the global error
$O(h^{4})$. RK4 is the working baseline for nonstiff problems.
The cost per step is four function evaluations.

Adaptive Runge-Kutta methods (Cash-Karp, Dormand-Prince) embed two
RK methods of different order in shared stage evaluations, use the
difference as a local error estimate, and adjust $h$ to keep the
local error within tolerance. The Dormand-Prince DOPRI5(4) is the
default in MATLAB \texttt{ode45} and in \texttt{scipy.integrate.RK45}.

\begin{figure}[ht]
\centering
\begin{tikzpicture}[x=1.5cm, y=1.5cm, font=\footnotesize]

% Absolute-stability regions in the complex z = h*lambda plane for
% forward Euler (disk of radius 1 at -1) and RK4 (the familiar bulb
% reaching to about -2.78 on the negative real axis). Backward Euler is
% A-stable (entire left half-plane), shown as the shaded half-plane edge.

% Axes
\draw[->, thick] (-3.4,0) -- (1.4,0) node[right] {$\Re(z)$};
\draw[->, thick] (0,-3.2) -- (0,3.2) node[above] {$\Im(z)$};

\foreach \x in {-3,-2,-1,1} {
  \draw (\x,0) -- (\x,-0.08);
  \node[anchor=north, font=\scriptsize] at (\x,-0.08) {\x};
}
\foreach \y in {-3,-2,-1,1,2,3} {
  \draw (0,\y) -- (-0.08,\y);
  \node[anchor=east, font=\scriptsize] at (-0.08,\y) {\y};
}

% Forward Euler: |1 + z| <= 1, disk of radius 1 centred at (-1, 0)
\draw[thick, engblue, fill=engblue!12] (-1,0) circle (1);
\node[engblue, font=\scriptsize, anchor=north] at (-1,-0.15)
  {forward Euler};

% RK4 absolute-stability boundary (caricature of the classic bulb,
% crossing the real axis near -2.78, slightly fatter than the Euler disk)
\draw[thick, engred] plot[smooth cycle] coordinates {
  (0.35,0) (0.1,1.0) (-0.6,1.9) (-1.6,2.3) (-2.4,1.6)
  (-2.78,0) (-2.4,-1.6) (-1.6,-2.3) (-0.6,-1.9) (0.1,-1.0)
};
\node[engred, font=\scriptsize, anchor=west] at (-2.4,1.95)
  {RK4};

% Backward Euler: A-stable, whole left half-plane (mark the boundary)
\draw[thick, dashed, engamber!70!black] (0,-3.1) -- (0,3.1);
\node[engamber!70!black, font=\scriptsize, anchor=west, align=left]
  at (0.15,-2.6) {backward Euler:\\ entire left half-plane\\ (A-stable)};

\end{tikzpicture}
\caption[Absolute-stability regions of Euler and RK4]{Absolute-stability
regions in the complex plane $z = h\lambda$, where $\lambda$ is the
relevant Jacobian eigenvalue and $h$ the step. A method is stable for a
given $h\lambda$ when $z$ lies inside its region. Forward Euler is stable
only inside the unit disk centred at $-1$; on the negative real axis it
admits $h\lambda \in (-2, 0)$, hence $h < 2/\abs{\lambda}$. RK4 extends
the real-axis reach to about $h\lambda \approx -2.78$, a modest gain that
does not change the conclusion for stiff systems. Backward Euler is
$A$-stable: its region is the entire closed left half-plane, so any
$h > 0$ is stable for a decaying mode, which is why implicit methods are
the working choice when one eigenvalue is large and negative. Source:
authors' analysis; the regions are standard, see
\cite[ch.~10]{text:moler-numerical}.}
\label{fig:vol02:ch16:rk-stability}
\end{figure}


Figure~\ref{fig:vol02:ch16:rk-stability} draws the absolute-stability
regions of forward Euler, RK4, and backward Euler in the complex
$z = h\lambda$ plane. A one-step method applied to $\dot y = \lambda
y$ advances by $y_{n+1} = R(z) y_{n}$ for a \engterm{stability
function} $R(z)$; the method is absolutely stable where
$\abs{R(z)} \leq 1$. Forward Euler has $R(z) = 1 + z$ (the unit disk
at $-1$); RK4 has $R(z) = 1 + z + z^{2}/2 + z^{3}/6 + z^{4}/24$, the
degree-$4$ Taylor truncation of $e^{z}$, whose stability region
reaches $z \approx -2.78$ on the negative real axis; backward Euler
has $R(z) = 1/(1 - z)$, stable across the entire left half-plane.

\begin{example}
\textbf{Reading a step-size cap off the stability region.} Integrate
$\dot y = -50 y$ with RK4. The real-axis stability limit is
$h\lambda \gtrsim -2.78$, so $h \lesssim 2.78/50 \approx 0.0556$. A
step of $h = 0.05$ gives $z = -2.5$, $R(z) = 1 - 2.5 + 3.125 -
2.604 + 1.628 = 0.649$, $\abs{R} < 1$, stable. A step of $h = 0.07$
gives $z = -3.5$, $R(z) = 1 - 3.5 + 6.125 - 7.146 + 6.253 = 2.732$,
$\abs{R} > 1$, and the numerical solution grows without bound while
the true solution decays to zero, the unmistakable signature of a
stability violation. The engineer who sees a decaying physical
system produce an exponentially growing numerical solution reads the
step size against the stability region before reaching for any other
explanation; the cure is a smaller step for an explicit method, or
an A-stable implicit method that removes the cap entirely.
\end{example}

\subsection{Stiffness}

A system is \engterm{stiff} when the eigenvalues of its Jacobian
have widely separated magnitudes, with at least one eigenvalue of
large negative real part. The fast mode forces an explicit
solver to take steps far smaller than the timescale of interest,
not for accuracy but for stability; the engineer pays in compute
time for transient modes that have already decayed.

The diagnosis is empirical: an explicit solver that demands very
small steps despite the solution being smooth is the canonical
signal. The remedy is an implicit solver. Backward Euler, the
trapezoidal rule, the implicit Runge-Kutta family, and the
backward differentiation formulas (BDF1 through BDF6) are the
working choices. \texttt{scipy.integrate.BDF} and MATLAB
\texttt{ode15s} are the standard production routines for stiff
systems. The cost of an implicit step is a Newton-style linear
solve per step; the savings are the lifted step-size restriction.

\begin{figure}[ht]
\centering
\begin{tikzpicture}[x=1.1cm, y=2.2cm, font=\footnotesize]

% The stiff scalar test y' = -100 y, y(0)=1, true solution decays to 0
% almost instantly. Forward Euler at h=0.025 (just past the stability
% limit h=0.02) oscillates with growing amplitude; backward Euler at the
% same h tracks the smooth decay monotonically.

\draw[->, thick] (0,0) -- (6.4,0) node[right] {$t$};
\draw[->, thick] (0,-1.15) -- (0,1.25) node[above] {$y$};

\foreach \t/\lab in {0/0, 2/0.05, 4/0.10, 6/0.15} {
  \draw (\t,0) -- (\t,-0.05);
  \node[anchor=north, font=\scriptsize] at (\t,-0.05) {$\lab$};
}
\foreach \y in {-1,-0.5,0.5,1} {
  \draw (0,\y) -- (-0.08,\y);
  \node[anchor=east, font=\scriptsize] at (-0.08,\y) {$\y$};
}

% True solution e^{-100 t}: drops to ~0 by t=0.05 (screen x=2)
\draw[thick, black] plot[domain=0:6, samples=80] (\x, {exp(-1.0*\x*\x/1.2 - 0.5*\x)});
% (a fast monotone decay caricature reaching ~0 by x=2.5)
\node[font=\scriptsize, anchor=west] at (2.6,0.18) {true $e^{-100t}$};

% Backward Euler at h=0.025: monotone decay, slightly above true, stable
% nodes at screen x = 0,1,2,3,4,5 (h=0.025 -> 5 steps to t=0.15? use 6 nodes)
\draw[thick, engblue] plot coordinates {
  (0,1) (1,0.31) (2,0.097) (3,0.030) (4,0.0095) (5,0.0030) (6,0.0009)
};
\foreach \p in {(0,1),(1,0.31),(2,0.097),(3,0.030),(4,0.0095),(5,0.0030),(6,0.0009)} {
  \fill[engblue] \p circle (1.4pt);
}
\node[engblue, font=\scriptsize, anchor=west] at (3.2,0.40) {backward Euler ($h=0.025$)};

% Forward Euler at h=0.025: factor (1 + h*lambda) = 1 - 2.5 = -1.5,
% so y_n = (-1.5)^n: oscillates and grows. Plot the growing oscillation.
\draw[thick, engred] plot coordinates {
  (0,1) (1,-0.5) (2,0.25) (3,-0.375) (4,0.5625) (5,-0.844) (6,1.266)
};
\foreach \p in {(0,1),(1,-0.5),(2,0.25),(3,-0.375),(4,0.5625),(5,-0.844),(6,1.266)} {
  \fill[engred] \p circle (1.4pt);
}
\node[engred, font=\scriptsize, anchor=south west] at (4.1,-0.95)
  {forward Euler ($h=0.025$): $(-1.5)^{n}$};

\end{tikzpicture}
\caption[Explicit versus implicit integration of a stiff scalar mode]{The
stiff test equation $\dot y = -100\,y$, $y(0)=1$, integrated past the
explicit stability limit. Forward Euler advances by the factor
$1 + h\lambda = 1 - 100h$; at $h = 0.025$ this factor is $-1.5$, so the
numerical solution is $(-1.5)^{n}$, oscillating in sign and growing without
bound while the true solution decays to zero. Backward Euler advances by
$1/(1 - h\lambda) = 1/(1 + 100h)$, a factor in $(0,1)$ for every $h>0$, and
tracks the decay monotonically. The explicit method is not wrong in its
local truncation order; it is unstable because the step exceeds the bound
$h < 2/\abs{\lambda} = 0.02$ set by the fast eigenvalue, and stability, not
accuracy, governs the usable step. The worked comparison in the text carries
both schemes for several steps by hand. Source: authors' analysis; the
test equation and stability factors are standard, see
\cite[ch.~10]{text:moler-numerical}.}
\label{fig:vol02:ch16:stiff-explicit-implicit}
\end{figure}


\begin{example}
\textbf{Explicit versus implicit on a stiff mode, carried by hand.}
Integrate the scalar stiff equation $\dot y = -100\,y$, $y(0) = 1$, whose
exact solution $y(t) = e^{-100 t}$ has decayed to $e^{-5} = 0.0067$ by
$t = 0.05$ and to $e^{-10} = 4.5 \times 10^{-5}$ by $t = 0.1$. Take a step
$h = 0.025$, chosen just past the explicit stability limit $h <
2/\abs{\lambda} = 0.02$, and compare. \emph{Forward Euler} advances by the
factor $1 + h\lambda = 1 - 100(0.025) = -1.5$. The iterates are $y_{0} =
1$, $y_{1} = -1.5$, $y_{2} = 2.25$, $y_{3} = -3.375$, $y_{4} = 5.0625$:
the numerical solution oscillates in sign and grows in magnitude as
$(-1.5)^{n}$ while the true solution decays smoothly toward zero. The
explicit step is not inaccurate in the local-truncation sense; it is
unstable, because $\abs{1 + h\lambda} = 1.5 > 1$ places $h\lambda = -2.5$
outside the forward-Euler stability disk. \emph{Backward Euler} on the same
equation solves $y_{n+1} = y_{n} + h\lambda y_{n+1}$, giving $y_{n+1} =
y_{n}/(1 - h\lambda) = y_{n}/(1 + 2.5) = y_{n}/3.5$. The iterates are
$y_{0} = 1$, $y_{1} = 0.2857$, $y_{2} = 0.0816$, $y_{3} = 0.0233$, $y_{4}
= 0.0067$: a monotone decay that tracks the true solution, every step
stable because $\abs{1/(1 - h\lambda)} = 1/3.5 < 1$ for the decaying mode.
To resolve the same equation accurately with forward Euler the engineer
must drop to $h < 0.02$, say $h = 0.01$, requiring $10$ steps to reach
$t = 0.1$; backward Euler at $h = 0.025$ reaches the same point in $4$
steps, each step costing one (here trivial) linear solve. On the scalar
problem the saving is modest; on a system with a stiffness ratio of
$10^{6}$, the explicit method's stability cap forces a step a million times
smaller than accuracy requires, and the implicit step's per-step linear
solve is the price that buys back the million-fold step-count saving.
Figure~\ref{fig:vol02:ch16:stiff-explicit-implicit} plots both schemes; the
explicit oscillation and the implicit decay are the visual signature the
engineer learns to read off any stiff integration that has been pushed past
its explicit limit.
\end{example}

\begin{mastery}
\textbf{Order barriers.} Two theorems of Dahlquist bound what an
ODE solver can achieve, and the engineer who knows them stops
hunting for a method that cannot exist. The \engterm{first Dahlquist
barrier} states that a stable linear multistep method of $k$ steps
has order at most $k + 1$ for odd $k$ and $k + 2$ for even $k$: there
is no free lunch in which more past values buy unbounded order while
keeping stability. The \engterm{second Dahlquist barrier} is the
sharper constraint for stiff problems: no explicit method is
$A$-stable, and no $A$-stable linear multistep method has order
greater than $2$, with the trapezoidal rule the $A$-stable method of
smallest error constant at order $2$. The barrier explains the
architecture of every production stiff solver. The backward
differentiation formulas BDF1 and BDF2 are $A$-stable; BDF3 through
BDF6 are not $A$-stable but are $A(\alpha)$-stable (stable in a wedge
around the negative real axis) and are used because real stiff
systems have eigenvalues near the negative real axis where the wedge
suffices; BDF7 and higher are unstable and do not exist as usable
methods. When the engineer needs both high order and $A$-stability,
the implicit Runge-Kutta family (which the barriers do not constrain,
since it is not a linear multistep method) supplies methods such as
the Radau and Gauss collocation schemes, at the cost of a larger
nonlinear solve per step. The barriers are durable mathematics, not
tooling: they will hold in 2050 exactly as in 1963
\cite[ch.~10]{text:moler-numerical}.
\end{mastery}

\subsection{Multistep methods and predictor-corrector pairs}

The Runge-Kutta methods reach high order by taking several stage evaluations
of $f$ within one step, throwing away the work once the step is taken. A
\engterm{linear multistep method} reaches high order instead by reusing the
$f$-values from previous steps, at one new evaluation per step. The
\engterm{Adams-Bashforth} family is the explicit branch: the two-step
second-order method is

\[
y_{n+1} = y_{n} + \frac{h}{2}\big(3 f_{n} - f_{n-1}\big),
\]

which fits a line through the two most recent slopes $f_{n} = f(t_{n}, y_{n})$
and $f_{n-1}$ and integrates that line forward. The \engterm{Adams-Moulton}
family is the implicit branch: the one-step second-order member is the
trapezoidal rule $y_{n+1} = y_{n} + \tfrac{h}{2}(f_{n} + f_{n+1})$, and
because $f_{n+1}$ depends on the unknown $y_{n+1}$, the formula is solved
rather than evaluated. Pairing the two is the classic
\engterm{predictor-corrector} scheme: the explicit Adams-Bashforth step
\emph{predicts} $y_{n+1}^{P}$; the implicit Adams-Moulton formula, with
$f_{n+1}$ evaluated at the prediction, \emph{corrects} it; the difference
$y_{n+1}^{C} - y_{n+1}^{P}$ is a local error estimate that drives step-size
control. The cost is two $f$-evaluations per step (one for the corrector
slope, one carried forward), against RK4's four, which is why
variable-order Adams predictor-correctors are the working choice when $f$ is
expensive and the problem is nonstiff. The backward differentiation
formulas of the stiffness subsection are the implicit multistep family
tuned for the stiff case; the Adams pair is its nonstiff sibling, and a
production solver such as \texttt{LSODA} switches between the two as it
detects stiffness.

\begin{example}
\textbf{A two-step Adams-Bashforth start, by hand.} Integrate $\dot y = -y$,
$y(0) = 1$, with $h = 0.1$, whose exact solution is $e^{-t}$. A two-step
method needs two starting values; take $y_{0} = 1$ and bootstrap $y_{1}$
with one RK4 step, $y_{1} = 0.904837$ (the calculation exercise of this
chapter), so $f_{0} = -1$ and $f_{1} = -0.904837$. The Adams-Bashforth-2
step gives $y_{2} = y_{1} + \tfrac{h}{2}(3 f_{1} - f_{0}) = 0.904837 +
0.05(3(-0.904837) - (-1)) = 0.904837 + 0.05(-1.714512) = 0.818111$, against
the exact $e^{-0.2} = 0.818731$, an error of $6.2 \times 10^{-4}$, the
$O(h^{2})$ accuracy the two-step formula promises at one new evaluation per
step. Correcting with Adams-Moulton-2 (the trapezoidal rule) using the
predicted slope $f_{2}^{P} = -0.818111$ gives $y_{2}^{C} = y_{1} +
\tfrac{h}{2}(f_{1} + f_{2}^{P}) = 0.904837 + 0.05(-0.904837 - 0.818111) =
0.818690$, error $4.1 \times 10^{-5}$, a decade better than the predictor
alone, and the predictor-corrector difference $\abs{y_{2}^{C} - y_{2}^{P}}
= 5.8 \times 10^{-4}$ is the local error estimate the step controller reads.
The example shows the three working facts of multistep methods: they need a
starter (here one RK4 step) because they reference past values that do not
exist at $t_{0}$; they reach their order at one new $f$-evaluation per step;
and the predictor-corrector gap is a free error estimate.
\end{example}

\begin{mastery}
\textbf{Structure-preserving integrators.} The error analysis of this
chapter has measured a method by its order, the power of $h$ at which the
local error vanishes. For a long integration of a conservative system, order
is the wrong figure of merit. Consider a Hamiltonian system, a planet
orbiting a star or a frictionless pendulum, whose true trajectory conserves
total energy exactly and preserves phase-space volume (Liouville's theorem).
A high-order method such as RK4 conserves neither: applied to the harmonic
oscillator $\ddot q = -q$ over many periods, RK4's energy drifts
monotonically (downward, since its stability region pulls the amplitude in),
and after enough orbits the simulated planet spirals into the star though no
physical force draws it there. The defect is not poor order; RK4's local
error is tiny. The defect is that RK4 does not respect the geometric
structure the dynamics live on. A \engterm{symplectic integrator} preserves
the phase-space area form exactly, and as a consequence its energy error
does not drift: it oscillates within a bounded band for exponentially long
times. The simplest symplectic method is \engterm{symplectic Euler}, which
updates position and momentum in a staggered order ($p_{n+1} = p_{n} -
h\,\partial_{q}H(q_{n})$ then $q_{n+1} = q_{n} + h\,\partial_{p}H(p_{n+1})$),
and the workhorse is the second-order \engterm{Stormer-Verlet} (leapfrog)
method, ubiquitous in molecular dynamics and orbital mechanics precisely
because its energy stays bounded over billions of steps where a
higher-order non-symplectic method drifts away
\cite[ch.~1]{text:v2ch07-hairer-lubich-wanner-geometric}. The reason a
low-order symplectic method beats a high-order generic one on the long haul
is \engterm{backward error analysis} of the geometric kind: the symplectic
method is the exact flow of a nearby \emph{modified} Hamiltonian
$\tilde H = H + O(h^{2})$, and because it exactly conserves $\tilde H$, the
true $H$ stays within $O(h^{2})$ of constant forever rather than walking
away linearly in time. The durable engineering lesson outlives any solver
library: the figure of merit for a long simulation of a conserved quantity
is whether the method preserves the structure that protects the quantity,
not the order of its local truncation error. The same principle governs the
choice of a divergence-free discretisation for an incompressible flow and a
charge-conserving scheme for a plasma; structure preservation is the
property, order is the refinement.
\end{mastery}

\subsection{Finite differences and finite elements (preview)}

Partial differential equations on a spatial domain require a
\engterm{discretisation} of the domain. \engterm{Finite difference}
methods replace partial derivatives by difference quotients on a
grid, producing a coupled system of ODEs in time (the
\engterm{method of lines}) or a coupled algebraic system at each
time step. \engterm{Finite element} methods discretise the domain
into elements (triangles, tetrahedra, quadrilaterals), express the
unknown as a linear combination of basis functions over the
elements, and reduce the PDE to a sparse linear system by Galerkin
projection. \engterm{Finite volume} methods integrate the
governing PDE over control volumes, expressing conservation
locally; they are the working choice for fluid mechanics.

The numerical analysis of PDEs is the subject of Volume IV
Chapter 8 (computational continuum mechanics) and Volume IX. The
preview here names three numerical issues whose presence the
reader of this chapter should recognise:

\begin{itemize}[leftmargin=*]
  \item The \engterm{Courant-Friedrichs-Lewy} (CFL) condition,
    which bounds the time step in explicit hyperbolic-PDE solvers
    by the spatial grid divided by the wave speed. Violating the
    CFL produces an instability that grows exponentially in the
    grid spacing.
  \item \engterm{Numerical diffusion}, the spurious smoothing
    introduced by low-order upwind schemes, which can dominate
    the physical diffusion in advection-dominated problems.
  \item The \engterm{inf-sup} (Ladyzhenskaya-Babuska-Brezzi)
    condition for mixed finite-element formulations, whose
    violation produces spurious modes in incompressible-flow and
    elasticity solvers.
\end{itemize}

These are not ornaments. Each is a condition under which the
discretisation is well-posed; each has been the subject of
documented engineering failures in flow-solver and
structural-solver deployments.

\section{Reproducibility: deterministic floating-point, exact-arithmetic alternatives}\pathtag{standard}

A computational result an engineer reports should be reproducible.
The simplest definition: the same code, on the same input, on the
same machine, produces the same output. The strictest definition:
the same code, on the same input, on any conformant hardware,
produces the same output to the last bit. The engineering
practice usually sits between these, and the engineer should know
where their report sits.

\subsection{Sources of non-determinism in floating-point computation}

The same source code can produce different floating-point output
across runs for several distinct reasons.

\begin{itemize}[leftmargin=*]
  \item \engterm{Compiler reordering}. The compiler may evaluate
    $a + b + c$ as $(a + b) + c$ or as $a + (b + c)$. Floating-point
    addition is not associative; the two orders can differ by a
    rounding error per operation. The \texttt{-ffast-math} flag in
    GCC and Clang permits algebraic reordering aggressively, with
    speed gains and bit-reproducibility loss.
  \item \engterm{Fused multiply-add (FMA)}. Hardware FMA computes
    $a \cdot b + c$ with one rounding; a separate multiply-then-add
    rounds twice. The two differ at the last bit. Whether the
    compiler emits an FMA is a target-architecture and
    optimisation-level decision.
  \item \engterm{Parallel reduction}. A sum across threads may
    combine partial sums in a thread-count-dependent order. The
    rounding error of the final sum depends on the order. A
    parallel reduction is bit-reproducible only when the reduction
    is implemented to fix the order, at a cost in performance.
  \item \engterm{Math-library variation}. The transcendental
    functions $\sin$, $\cos$, $\exp$, $\log$ are not, in general,
    correctly rounded under IEEE 754. The implementation in
    glibc, in the Intel math library, in Apple's Accelerate, and
    in CUDA's libm are all conformant within the standard's
    tolerances, but differ at the last bit. An engineer who
    depends on the last bit of $\sin(0.123)$ depends on the
    library shipped with the platform.
  \item \engterm{Subnormal handling}. The flush-to-zero mode
    converts subnormals to zero on input or output of arithmetic
    operations, controlled by per-thread CPU flags that some
    libraries set without notification.
  \item \engterm{Random-seed and concurrency-order} effects in
    Monte Carlo and stochastic-optimisation code, which compound
    with the floating-point sources above.
\end{itemize}

\begin{mastery}
\textbf{Non-associativity quantified.} The root of every reordering
effect is that floating-point addition is not associative. A
concrete instance: let $a = 1$, $b = -1$, $c = 10^{-20}$ in
binary64. Then $(a + b) + c = 0 + 10^{-20} = 10^{-20}$, but
$a + (b + c) = a + \operatorname{fl}(-1 + 10^{-20})$, and since
$10^{-20}$ is far below the spacing of $1$, $\operatorname{fl}(-1 +
10^{-20}) = -1$ exactly, so $a + (b + c) = 0$. The two groupings
differ by the entire value of $c$. A parallel reduction that splits
the sum across threads chooses a grouping at run time, and the
grouping depends on the thread count, which is why the same code
gives different last bits on $2$ threads and on $8$. The cost of
fixing the grouping (a deterministic reduction tree, fixed
regardless of thread count) is a loss of the dynamic load balancing
that makes the parallel sum fast in the first place, typically a
$10$ to $30$ percent slowdown on a reduction-heavy kernel. The
engineering decision is whether the application needs bit-identical
output (regulatory audit, differential debugging, exact regression
tests) or only statistically reproducible output (most scientific
reporting). The first pays the determinism tax; the second documents
the run-to-run variance and moves on. Stating which choice the
report makes is the discipline; leaving it implicit is the defect.
\end{mastery}

\subsection{Practical reproducibility disciplines}

The engineer's defences are concrete.

\begin{itemize}[leftmargin=*]
  \item Pin the compiler version, optimisation level, and
    floating-point flags in a build manifest committed alongside
    the source. \texttt{-O2 -fno-fast-math} is the working
    default for reproducibility.
  \item Pin the random seed in any code that draws random numbers,
    and document the random-number generator (its name and its
    state size).
  \item Pin the library versions: NumPy, SciPy, BLAS (OpenBLAS
    versus MKL versus Apple Accelerate), LAPACK, and the math
    library. A \texttt{requirements.txt} or a Conda environment
    file or a Nix derivation is the artefact.
  \item For parallel code, set the thread count explicitly
    (\texttt{OMP\_NUM\_THREADS}, \texttt{MKL\_NUM\_THREADS}) and
    use deterministic-reduction routines where available (the
    Intel MKL Conditional Numerical Reproducibility mode, the
    \texttt{torch.use\_deterministic\_algorithms(True)} flag in
    PyTorch).
  \item Report the platform: operating system, CPU model, and
    library versions, in the project artefact.
\end{itemize}

The cost of bit-reproducibility is real. Production scientific code
that has been tuned for parallel performance often gives up
bit-reproducibility for speed. The discipline is then to report a
\emph{statistically} reproducible result: the answer to within a
stated tolerance under repeated runs, with the variance of the
output documented.

\subsection{Exact-arithmetic alternatives}

For some computations, floating-point is the wrong arithmetic.
\engterm{Exact rational arithmetic} represents numbers as ratios
of arbitrary-precision integers, with no rounding error. The Python
\texttt{fractions} module, the \texttt{Rational} type in many
computer-algebra systems, and SageMath provide working
implementations. The cost is performance: the integers grow as the
computation proceeds, and a long sequence of operations on
rationals can produce numerators and denominators of ruinous size.

\engterm{Arbitrary-precision floating-point} (Python's
\texttt{decimal} and \texttt{mpmath} modules, MPFR, the GMP MPF
type) lets the engineer set the precision per computation. A
condition number of $10^{30}$ is intractable in double precision
and routine in 100-digit MPFR.

\engterm{Symbolic computation} (SymPy, Mathematica, Maple)
operates on expressions rather than on numerical approximations,
yielding closed-form answers when they exist and exact rational
or algebraic outputs when they do not. The cost is the engineer's
time: setting up a symbolic computation is slower than a
floating-point one, and the symbolic tool's behaviour at the
edges of its capability is sometimes less predictable than the
floating-point library's.

\engterm{Interval arithmetic} represents each quantity as an
interval $[a, b]$ guaranteed to contain the true value, and
propagates intervals through the computation. The output is an
interval that bounds the true result, regardless of the
floating-point error of any operation. The intervals tend to grow
faster than the engineer would like (the \engterm{wrapping
effect}), but for verification of critical quantities, interval
arithmetic is the formal answer to the question \emph{can we
prove the result lies inside this band}. Sun, Boost, and INTLAB
provide working implementations.

The engineering choice among these alternatives is a working
decision. Most computations the engineer reports will be done in
double-precision floating-point with a documented build
environment. The alternatives are tools the engineer reaches for
when the problem demands them, and the report names the
arithmetic used.

\section{A worked computation with an error budget}\pathtag{core}

The chapter's apparatus, the three error sources, the conditioning, the
algorithm choice, the reproducibility discipline, earns its keep only when
it is carried through one realistic computation from the posed engineering
question to a verified number with a stated uncertainty. We run such a
computation here. The question is small enough to verify against an exact
solution and rich enough to exercise every error source the chapter names.

\subsection{The problem and its exact reference}

A uniform cantilever beam of length $L$, Young's modulus $E$, and
second moment of area $I$, fixed at $x = 0$ and free at $x = L$, carries
a uniformly distributed load $w$ per unit length. The static deflection
$y(x)$ satisfies the Euler-Bernoulli equation $E I\, y''''(x) = w$ with
the clamped-end conditions $y(0) = y'(0) = 0$ and the free-end conditions
$y''(L) = y'''(L) = 0$. The quantity of interest is the tip deflection
$y(L)$. This problem has the exact closed form

\[
y(x) = \frac{w}{24 E I}\left(x^{4} - 4 L x^{3} + 6 L^{2} x^{2}\right),
\qquad
y(L) = \frac{w L^{4}}{8 E I},
\]

so the engineer can check every numerical decision against a known answer,
which is the right way to build trust in a method before turning it loose
on a problem with no closed form. Take the working values $L = 2.000\,
\text{m}$, $E = 200\,\text{GPa}$ (structural steel), $I = 4.000 \times
10^{-6}\,\text{m}^{4}$, and $w = 5.000 \times 10^{3}\,\text{N}/\text{m}$.
The exact tip deflection is

\[
y(L) = \frac{(5000)(2.000)^{4}}{8 (200 \times 10^{9})(4.000 \times 10^{-6})}
= \frac{80000}{6.400 \times 10^{6}} = 1.250 \times 10^{-2}\,\text{m}
= 12.50\,\text{mm}.
\]

\begin{estimation}
\textbf{Question.} Before discretising, estimate the tip deflection from
dimensional reasoning and a single sanity number, and estimate which of the
three error sources (truncation, round-off, conditioning) will dominate the
numerical result on a grid of a few dozen intervals.

\smallskip
\textit{Estimate, before reading on.}

\smallskip
The deflection must scale as $w L^{4}/(E I)$ by dimensions; the only
question is the dimensionless prefactor, which for a tip-loaded or
distributed-load cantilever is an order-unity rational ($1/8$ for the
uniformly distributed case). With $w L^{4}/(E I) = (5000)(16)/(8 \times
10^{5}) = 0.1\,\text{m}$ and a prefactor near $1/8$, the tip deflection is
near $12\,\text{mm}$, a span-over-$160$ deflection that a structural
engineer reads as plausible for a working beam. For the error sources: a
second-order finite-difference scheme on $n$ intervals has truncation error
$O(h^{2}) = O(1/n^{2})$, near $4 \times 10^{-4}$ relative at $n = 50$; the
linear solve is tridiagonal with a condition number that grows as
$O(n^{4})$ for a fourth-order operator but stays near $10^{6}$ at $n = 50$,
so the round-off contribution $\kappa u$ is near $10^{-10}$; and the input
data ($E$ and $I$ from material and section tolerances) carry a few percent
each. The dominant source is therefore neither truncation nor round-off
but the input-data uncertainty, by two orders of magnitude. The estimate
sets the expectation the budget will confirm: refining the grid past the
data floor is wasted work.
\end{estimation}

\subsection{Discretisation and its truncation error}

Replace the fourth derivative by the standard five-point centred
difference on a uniform grid $x_{i} = i h$, $h = L/n$, $i = 0, \ldots, n$:

\[
y''''(x_{i}) \approx \frac{y_{i-2} - 4 y_{i-1} + 6 y_{i} - 4 y_{i+1}
+ y_{i+2}}{h^{4}},
\]

whose truncation error is $-\frac{h^{2}}{6} y^{(6)}(\xi) + O(h^{4})$, of
order $h^{2}$. The clamped conditions $y_{0} = 0$ and $y'(0) = 0$ (the
latter as $(y_{1} - y_{-1})/(2h) = 0$, introducing a ghost node $y_{-1} =
y_{1}$) and the free-end conditions on $y''$ and $y'''$ close the system at
the boundaries. The result is a banded linear system $\mat{A}\vect{y} =
\vect{f}$ with $\vect{f}_{i} = w/(E I)$ at the interior nodes, where
$\mat{A}$ is pentadiagonal of size $(n - 1) \times (n - 1)$ after the
boundary eliminations. The truncation error of the computed deflection
inherits the scheme's order: $\abs{y_{n}^{\text{computed}} - y(L)} =
O(h^{2})$ for this smooth (in fact quartic) solution, and because the exact
solution is a degree-four polynomial whose sixth derivative is zero, the
leading $h^{2}$ truncation term vanishes identically and the scheme is
exact up to round-off on this particular problem. The engineer notes this
as a feature of the verification case, not of the method: the method's
$O(h^{2})$ behaviour must be confirmed on a problem whose solution is not
a low-degree polynomial before the zero-truncation result here is trusted
as evidence of a correct implementation rather than of an accidental
exactness.

\subsection{The linear solve and its round-off}

The system is solved by banded LU with partial pivoting (the routine behind
\texttt{scipy.linalg.solve\_banded}), which is backward stable with
backward error of order $u$ times a modest growth factor. The forward error
of the solve is bounded by $\kappa(\mat{A}) \cdot u$ times that factor. The
condition number of the discrete fourth-derivative operator grows as
$\kappa(\mat{A}) = O(n^{4})$, the discrete analogue of the
fourth-order differential operator's unbounded spectrum; at $n = 50$,
$\kappa(\mat{A}) \approx 4 \times 10^{6}$, so the round-off contribution to
the tip deflection is bounded by $\kappa u \approx 4 \times 10^{6} \cdot
1.1 \times 10^{-16} \approx 4 \times 10^{-10}$, ten significant digits
intact. Doubling $n$ to $100$ raises $\kappa$ by $2^{4} = 16$ and the
round-off floor to $\approx 7 \times 10^{-9}$, still negligible against the
data. The $O(n^{4})$ growth is the reason a naive engineer who refines the
grid toward $n = 10^{4}$ in pursuit of truncation accuracy eventually meets
a round-off wall: at $n = 10^{4}$, $\kappa \approx 4 \times 10^{14}$ and
the solve retains only one or two digits. The defence, when high spatial
resolution is genuinely needed, is to reformulate (a second-order system
in $y$ and the moment $M = E I y''$ rather than the fourth-order system in
$y$ alone), which lowers the operator order and the conditioning growth
from $n^{4}$ to $n^{2}$.

\subsection{The input-data uncertainty}

The dominant term, as the estimation block predicted, is the uncertainty in
the inputs. Structural steel's Young's modulus is quoted as $200\,\text{GPa}$
with a typical batch-to-batch and temperature variation of about $\pm 3\,
\text{percent}$; a rolled or machined section's second moment of area
carries a tolerance from the dimensional tolerances on the cross-section,
typically $\pm 2$ to $\pm 5\,\text{percent}$ depending on the section class;
the distributed load $w$ in service is known to perhaps $\pm 5\,\text{percent}$.
Propagating these through $y(L) = w L^{4}/(8 E I)$ by the
relative-uncertainty rule of Volume I Chapter 4, with the relative
uncertainties combining in quadrature for independent sources,

\[
\frac{u(y_{L})}{y_{L}} = \sqrt{\left(\frac{u_{w}}{w}\right)^{2} +
\left(\frac{u_{E}}{E}\right)^{2} + \left(\frac{u_{I}}{I}\right)^{2}
+ \left(4\frac{u_{L}}{L}\right)^{2}}
\approx \sqrt{0.05^{2} + 0.03^{2} + 0.04^{2} + (4 \cdot 0.005)^{2}}
\approx 0.075,
\]

a relative uncertainty near $7.5\,\text{percent}$, dominated by the load
and the section tolerances and amplified in the length by the fourth power.
The tip deflection is therefore $12.5\,\text{mm} \pm 0.9\,\text{mm}$,
current as of the quoted material data, and the honest report carries that
band, not the ten-digit precision the linear solve preserves.

\subsection{The budget and the verified result}

\begin{figure}[ht]
\centering
\begin{tikzpicture}[x=1cm, y=1cm, font=\footnotesize]

% A horizontal stacked bar showing the error budget of the case-study
% computation (cantilever deflection by finite differences + linear
% solve). The contributions, on a log scale of magnitude, are:
% discretisation (truncation, O(h^2)), linear-solve round-off
% (kappa * u), input/data uncertainty (E and I tolerances), and the
% reported result tolerance. The point: the data uncertainty dominates,
% so refining the grid past the round-off floor buys nothing.

% Log axis: error magnitude from 10^{-16} (left) to 10^{-1} (right)
\draw[->, thick] (0,0) -- (11,0) node[right] {relative error magnitude};
\foreach \k/\x in {-16/0, -13/1.5, -10/3, -7/4.5, -4/6, -1/7.5} {
  \draw (\x,0) -- (\x,-0.12);
  \node[anchor=north, font=\scriptsize] at (\x,-0.12) {$10^{\k}$};
}

% bars: each from its left edge to a point at its magnitude. Use y-rows.
% round-off of linear solve: kappa ~ 10^4, u ~ 10^{-16} -> ~10^{-12} (x = 2)
\fill[engblue!30] (0,1.2) rectangle (2,1.7);
\draw (0,1.2) rectangle (2,1.7);
\node[anchor=west, font=\scriptsize] at (2.15,1.45)
  {linear-solve round-off $\kappa u \approx 10^{-12}$};

% discretisation: O(h^2) with h=L/50 -> ~10^{-3} (x = 6.5)
\fill[engred!30] (0,2.2) rectangle (6.5,2.7);
\draw (0,2.2) rectangle (6.5,2.7);
\node[anchor=west, font=\scriptsize] at (6.65,2.45)
  {discretisation $O(h^{2}) \approx 10^{-3}$};

% data uncertainty: E,I tolerances ~ 5% -> ~5x10^{-2} (x = 7.0)
\fill[engamber!40] (0,3.2) rectangle (7.0,3.7);
\draw (0,3.2) rectangle (7.0,3.7);
\node[anchor=west, font=\scriptsize] at (7.15,3.45)
  {input data $E,I$ tolerance $\approx 5\times 10^{-2}$};

% defensible reported tolerance (the max, dominated by data) ~ 5x10^{-2}
\draw[very thick, dashed, black] (7.0,0.9) -- (7.0,4.0)
  node[above, font=\scriptsize, anchor=south] {reported tolerance};

\end{tikzpicture}
\caption[Error budget of the case-study deflection computation]{The error
budget of the worked cantilever-deflection computation, with each
contribution placed on a logarithmic magnitude axis. The linear-solve
round-off, bounded by $\kappa(\mat{A})\,u \approx 10^{-12}$ for the
tridiagonal system, is the smallest term. The discretisation (truncation)
error of the second-order finite-difference scheme on $50$ intervals is near
$10^{-3}$. The input-data uncertainty, set by the manufacturing tolerances
on Young's modulus $E$ and the second moment of area $I$, is near five
percent and dominates the budget. The reported tolerance is the largest
term, so refining the grid below the discretisation level, or computing the
solve in higher precision, buys nothing the data can support. Naming which
term dominates is the discipline; the budget makes the dominant term
visible. Source: authors' analysis.}
\label{fig:vol02:ch16:error-budget}
\end{figure}


Figure~\ref{fig:vol02:ch16:error-budget} stacks the four contributions on
a logarithmic magnitude axis. The ranking is unambiguous: the input-data
uncertainty at $7.5 \times 10^{-2}$ dominates the discretisation error at
$\sim 10^{-3}$ (for a grid fine enough to resolve a non-polynomial
solution), which dominates the linear-solve round-off at $\sim 10^{-10}$.
The verified numerical result is

\[
y(L) = 12.5\,\text{mm} \pm 0.9\,\text{mm}\ (\text{relative } 7.5\,\text{percent}),
\]

with the numerical method contributing nothing the data can resolve. Three
disciplines close the computation. First, the method was verified against
the exact closed form before being trusted, and its $O(h^{2})$ order was
confirmed on a separate non-polynomial test problem so that the exactness
on the quartic was not mistaken for an implementation that happens to work.
Second, the dominant error source was named and quantified, so the
engineer knows that buying a finer grid or a higher precision is wasted
money against this data, and that tightening the section tolerance or the
load estimate is the only path to a tighter deflection. Third, the reported
number carries its uncertainty and its provenance, so a reader can audit
the claim. The same three steps scale to a finite-element model with a
million degrees of freedom, where the closed form is unavailable, the
verification runs against a manufactured solution or a mesh-refinement
study, and the error budget is the artefact that tells the engineer which
term to attack. The discipline does not change with the problem size; only
the cost of each step does.

\section{The fast Fourier transform and spectral differentiation}\pathtag{mastery}

Two of the most consequential algorithms in computational engineering, the
fast Fourier transform and the spectral-differentiation methods built on it,
share a single idea: a smooth periodic or analytic function is represented
exactly by a short list of coefficients in an orthogonal basis, and the
operations the engineer wants (differentiation, convolution, filtering) are
diagonal or near-diagonal in that basis. This section develops the working
arithmetic of both, at the depth a reader who calls \texttt{numpy.fft.fft}
or builds a Chebyshev solver should be able to defend.

\subsection{The discrete Fourier transform and its fast algorithm}

The \engterm{discrete Fourier transform} (DFT) of a length-$N$ sequence
$x_{0}, \ldots, x_{N-1}$ is the sequence

\[
X_{k} = \sum_{j=0}^{N-1} x_{j}\, \omega_{N}^{jk},
\qquad \omega_{N} = e^{-2\pi i/N},
\qquad k = 0, \ldots, N-1,
\]

a matrix-vector product $\vect{X} = \mat{F}_{N}\vect{x}$ with the symmetric
DFT matrix $(\mat{F}_{N})_{jk} = \omega_{N}^{jk}$. Computed as written, the
product costs $N^{2}$ complex multiply-adds. The \engterm{fast Fourier
transform} (FFT) computes the identical $\vect{X}$ in $O(N \log N)$
operations by exploiting the algebraic structure of $\omega_{N}$. The
radix-2 Cooley-Tukey factorisation splits the sum into its even-indexed and
odd-indexed terms,

\[
X_{k} = \underbrace{\sum_{j=0}^{N/2-1} x_{2j}\,\omega_{N/2}^{jk}}_{E_{k}}
+ \omega_{N}^{k}\underbrace{\sum_{j=0}^{N/2-1} x_{2j+1}\,\omega_{N/2}^{jk}}_{O_{k}},
\]

using $\omega_{N}^{2} = \omega_{N/2}$. Each of $E_{k}$ and $O_{k}$ is a DFT
of length $N/2$, and the periodicity $E_{k+N/2} = E_{k}$ lets the full
length-$N$ transform be assembled from the two half-transforms by the
\engterm{butterfly} combination $X_{k} = E_{k} + \omega_{N}^{k} O_{k}$ and
$X_{k+N/2} = E_{k} - \omega_{N}^{k} O_{k}$ for $k = 0, \ldots, N/2 - 1$.
Recursing the split until the transforms are length one gives the recurrence
$T(N) = 2 T(N/2) + O(N)$, whose solution is $T(N) = O(N \log N)$.

\begin{figure}[ht]
\centering
\begin{tikzpicture}[x=1cm, y=1cm, font=\footnotesize]

% Inputs on the left
\node[anchor=east] at (0,2) {$E_{k}$};
\node[anchor=east] at (0,0) {$O_{k}$};

% Nodes
\fill[engblue] (0,2) circle (2pt);
\fill[engblue] (0,0) circle (2pt);

% Twiddle multiply on the lower line
\node[draw, fill=white, inner sep=2pt] (tw) at (2,0) {$\times\,\omega_{N}^{k}$};

% Combination nodes on the right
\fill[engred] (4,2) circle (2pt);
\fill[engred] (4,0) circle (2pt);

% Edges: E_k straight to both outputs
\draw[->, thick] (0,2) -- (4,2);
\draw[->, thick] (0,2) -- (4,0);

% O_k through twiddle, then to both outputs (plus and minus)
\draw[->, thick] (0,0) -- (tw);
\draw[->, thick] (tw) -- (4,2);
\draw[->, thick] (tw) -- (4,0);

% Plus / minus annotations at the combination
\node[anchor=south west, font=\scriptsize] at (3.1,2.05) {$+$};
\node[anchor=north west, font=\scriptsize] at (3.1,-0.05) {$-$};

% Output labels
\node[anchor=west] at (4.15,2) {$X_{k} = E_{k} + \omega_{N}^{k} O_{k}$};
\node[anchor=west] at (4.15,0) {$X_{k+N/2} = E_{k} - \omega_{N}^{k} O_{k}$};

\end{tikzpicture}
\caption[The radix-2 FFT butterfly]{One radix-2 Cooley-Tukey butterfly. The
two inputs are the $k$-th values of the even-index and odd-index
half-transforms, $E_{k}$ and $O_{k}$. The odd value is scaled by the twiddle
factor $\omega_{N}^{k}$, and the two outputs $X_{k}$ and $X_{k+N/2}$ are formed
by one addition and one subtraction. A length-$N$ transform is $\log_{2} N$
such stages, each touching all $N$ values once, giving the $N \log_{2} N$
operation count of the fast Fourier transform. Source: authors' diagram of the
standard radix-2 decimation-in-time butterfly.}
\label{fig:vol02:ch16:fft-butterfly}
\end{figure}


Figure~\ref{fig:vol02:ch16:fft-butterfly} draws one radix-2 butterfly stage:
two inputs, one twiddle factor $\omega_{N}^{k}$, two outputs formed by a
sum and a difference. The full transform is $\log_{2} N$ such stages, each
touching all $N$ values once, which is the $N \log_{2} N$ operation count
read directly off the diagram.

\begin{example}
\textbf{What $O(N \log N)$ buys, in seconds.} A one-dimensional transform of
$N = 10^{6}$ points costs $N^{2} = 10^{12}$ operations by the direct DFT and
$N \log_{2} N \approx 2 \times 10^{7}$ operations by the FFT, a ratio near
$5 \times 10^{4}$. At a sustained $10^{10}$ operations per second, the direct
DFT takes about $100$ seconds and the FFT about two milliseconds. The gap
widens with $N$: at $N = 10^{9}$ the direct transform is a $10^{18}$-operation
computation that runs for years, while the FFT is a $3 \times 10^{10}$-operation
computation that runs in seconds. The FFT is the reason real-time spectral
analysis, digital filtering, large-scale convolution, and spectral PDE
solvers are feasible at all; the same physics computed through the direct DFT
would be intractable. Cooley and Tukey published the radix-2 algorithm in
1965, though Gauss had used the same factorisation by hand in 1805; the
modern rediscovery, arriving with digital computers that could exploit it,
is among the most cited results in computational science.
\end{example}

The DFT diagonalises circular convolution: the transform of a convolution is
the pointwise product of the transforms, so an $O(N^{2})$ convolution becomes
two forward FFTs, an $O(N)$ pointwise multiply, and one inverse FFT, total
$O(N \log N)$. The same diagonalisation makes differentiation of a periodic
function a multiplication: if $x_{j}$ samples a periodic $f$ on a uniform
grid, then the derivative samples are recovered by multiplying the
$k$-th Fourier coefficient by $ik$ (the wavenumber) and transforming back.
This is \engterm{Fourier spectral differentiation}, and for a smooth periodic
function its error decreases faster than any power of the grid spacing,
\engterm{spectral accuracy}.

\subsection{Chebyshev spectral differentiation}

Most engineering domains are not periodic. On a bounded non-periodic interval
the Fourier basis is replaced by the Chebyshev basis, and the uniform grid is
replaced by the Chebyshev points $x_{j} = \cos(j\pi/N)$, $j = 0, \ldots, N$,
the same cosine-spaced points that cured the Runge phenomenon. A function
sampled at the Chebyshev points has a unique interpolating polynomial, and
the derivative of that interpolant at the grid points is a linear map: the
sampled derivative vector is $\mat{D}_{N}\vect{u}$ for the
\engterm{Chebyshev differentiation matrix} $\mat{D}_{N}$, an $(N+1) \times
(N+1)$ dense matrix with entries known in closed form
\cite[ch.~6]{text:v2ch06-trefethen-spectral}. The corner entries are
$(\mat{D}_{N})_{00} = (2N^{2}+1)/6$ and $(\mat{D}_{N})_{NN} =
-(2N^{2}+1)/6$; the off-diagonal entries are $(\mat{D}_{N})_{ij} =
\frac{c_{i}}{c_{j}}\frac{(-1)^{i+j}}{x_{i}-x_{j}}$ with $c_{0} = c_{N} = 2$
and $c_{i} = 1$ otherwise; the interior diagonal entries are
$(\mat{D}_{N})_{jj} = -x_{j}/(2(1 - x_{j}^{2}))$.

\begin{example}
\textbf{The $N = 2$ Chebyshev differentiation matrix, and its accuracy.} For
$N = 2$ the Chebyshev points are $x_{0} = 1$, $x_{1} = 0$, $x_{2} = -1$. The
differentiation matrix is

\[
\mat{D}_{2} = \begin{psmallmatrix}
1.5 & -2 & 0.5 \\
0.5 & 0 & -0.5 \\
-0.5 & 2 & -1.5
\end{psmallmatrix}.
\]

Apply it to the sampled values of $f(x) = x^{2}$, which are $\vect{u} =
(1, 0, 1)^{\transpose}$: the product $\mat{D}_{2}\vect{u} = (1.5 + 0.5,
0.5 - 0.5, -0.5 - 1.5)^{\transpose} = (2, 0, -2)^{\transpose}$, which is
exactly $f'(x) = 2x$ sampled at $(1, 0, -1)$. The matrix differentiates
$x^{2}$ with zero error because the degree-two interpolant through three
points \emph{is} $x^{2}$, so its derivative is exact. This is the spectral
method's signature: for a polynomial of degree at most $N$ the differentiation
is exact, and for a smooth non-polynomial the error falls faster than any
power of $1/N$ as $N$ grows. A second-order finite difference would need
thousands of points to match what the Chebyshev method achieves with a
few dozen on a smooth problem.
\end{example}

\begin{figure}[ht]
\centering
\begin{tikzpicture}[x=0.55cm, y=0.40cm, font=\footnotesize]

% Semi-log: x = N (grid points), y = log10(max error)
% Spectral: straight line on semi-log (geometric decay) until floor.
% Finite difference O(h^2): on a semi-log-in-N axis this is a curve that
% flattens; we draw it as a slowly descending curve for contrast.

% Axes
\draw[->, thick] (0,-16) -- (34,-16) node[right] {$N$ (grid points)};
\draw[->, thick] (0,-16) -- (0,1) node[above, align=center]
  {$\log_{10}(\text{max error})$};

% x ticks
\foreach \k in {0,8,16,24,32} {
  \draw (\k,-16) -- (\k,-16.3);
  \node[anchor=north, font=\scriptsize] at (\k,-16.3) {$\k$};
}
% y ticks
\foreach \d in {0,-4,-8,-12,-16} {
  \draw (0,\d) -- (-0.4,\d);
  \node[anchor=east, font=\scriptsize] at (-0.4,\d) {$10^{\d}$};
}

% Spectral: geometric decay, straight line on semi-log, then floor at ~10^{-14}
\draw[thick, engblue] plot coordinates {
  (2,-0.5) (6,-3) (10,-5.5) (14,-8) (18,-10.5) (22,-13) (26,-14) (32,-14)
};
\node[engblue, font=\scriptsize, anchor=west] at (14,-6)
  {Chebyshev spectral};

% Finite difference O(h^2), h = 1/N: error ~ N^{-2}; on this axis a slow curve
\draw[thick, engred] plot coordinates {
  (2,-0.6) (4,-1.2) (8,-1.8) (16,-2.4) (24,-2.8) (32,-3.0)
};
\node[engred, font=\scriptsize, anchor=west] at (20,-2.4)
  {centred difference $O(h^{2})$};

% Round-off floor
\draw[dashed, enggray] (0,-14) -- (34,-14);
\node[enggray, font=\scriptsize, anchor=west] at (26,-15) {round-off floor};

\end{tikzpicture}
\caption[Spectral versus algebraic convergence]{Maximum error of the
derivative of a smooth test function against the number of grid points $N$, on
a semi-logarithmic axis. The Chebyshev spectral derivative falls along a
straight line, the signature of geometric (exponential) convergence, until it
reaches the round-off floor near $10^{-14}$; a few dozen points already exhaust
the available precision. The second-order centred finite difference falls only
algebraically and is decades behind at every $N$. Spectral accuracy holds for
smooth solutions on simple domains; a jump or corner in the function destroys
it. Source: authors' analysis, after the spectral-convergence demonstrations of
\cite[ch.~6]{text:v2ch06-trefethen-spectral}.}
\label{fig:vol02:ch16:spectral-convergence}
\end{figure}


Figure~\ref{fig:vol02:ch16:spectral-convergence} plots, on a semi-logarithmic
axis, the maximum error of the Chebyshev derivative of a smooth test function
against the number of grid points, beside the algebraic $O(h^{2})$ decay of a
centred finite difference. The spectral error falls along a straight line on
the semi-log plot, the signature of geometric (exponential) convergence,
until it reaches the round-off floor near $10^{-14}$; the finite-difference
error falls along a line of slope $-2$ on a log-log plot, decades slower. The
trade is that the spectral differentiation matrix is dense ($O(N^{2})$ work
per derivative, or $O(N \log N)$ via the FFT for the Chebyshev case) and
sensitive to non-smoothness: a single corner or jump in the function destroys
the spectral convergence and reintroduces oscillation, the Gibbs phenomenon.
The engineer reaches for spectral methods when the solution is smooth and the
domain is simple, and for finite elements when the geometry is complicated or
the solution has fronts and corners.

\begin{mastery}
\textbf{Interval arithmetic and verified computing.} Every error bound in
this chapter is an \emph{a priori} estimate: the engineer reasons about the
worst case before running, then trusts that the run stayed inside it. Interval
arithmetic delivers an \emph{a posteriori} guarantee instead, a machine-checked
enclosure of the true answer. Each quantity is carried as an interval
$[\underline{x}, \overline{x}]$ guaranteed to contain the true value, and
each operation rounds its result outward (the lower endpoint toward
$-\infty$, the upper toward $+\infty$, using the directed rounding modes IEEE
754 provides exactly for this purpose), so the output interval provably
contains the true result regardless of round-off. The \engterm{interval
Newton} method is the structure-revealing application: to enclose a root of
$f$ in $[a, b]$, form $N([a,b]) = m - f(m)/F'([a,b])$ where $m$ is the
midpoint and $F'([a,b])$ is an interval enclosure of $f'$ over the box; if
$N([a,b]) \subset [a,b]$ the method has \emph{proved} that a unique root
exists in the box and that every iterate stays enclosed, a theorem the
floating-point Newton of section 16.3 can never deliver. The cost is the
\engterm{wrapping effect}: naive interval arithmetic overestimates, because
it treats every occurrence of a variable as independent and so loses the
cancellation a real computation would enjoy; the interval for $x - x$ is
$[\underline{x} - \overline{x}, \overline{x} - \underline{x}]$, a band of
width $2(\overline{x} - \underline{x})$ rather than the true zero. Taming the
wrapping effect (centred forms, affine arithmetic, Taylor models) is the
craft of verified computing, and the payoff is a result the engineer can put
in front of a regulator as a proof rather than an estimate. Verified interval
computations have closed open mathematical problems (the Kepler conjecture on
sphere packing, the existence of the Lorenz attractor) and underwrite
safety-critical control where an unproven bound is unacceptable. INTLAB,
Boost.Interval, and Arb provide working implementations
\cite[ch.~26]{text:higham-numerical}.
\end{mastery}

\section{A second worked computation: a convection-diffusion boundary-value problem}\pathtag{core}

The cantilever computation of section 16.10 ran a fourth-order
boundary-value problem whose exact solution was a polynomial, so the
discretisation error vanished and the budget was dominated by input data.
A second case study earns its place by exercising the error sources the
first one suppressed: a problem whose discretisation can become
\emph{unstable} at the wrong grid resolution, whose remedy trades one error
for another, and whose verification requires a manufactured solution because
no clean closed form is available across the parameter range. The
one-dimensional steady convection-diffusion equation is the canonical
vehicle, and the lessons it teaches govern every advection-dominated flow,
heat-transfer, and transport solver the reader will later build.

\subsection{The problem and the Peclet number}

A scalar quantity $u(x)$, a temperature or a concentration, on $[0, 1]$ is
carried by a flow of speed $v$ and spread by diffusion of coefficient $D$,
in steady state with fixed end values:

\[
v\, u'(x) - D\, u''(x) = 0, \qquad u(0) = 0, \quad u(1) = 1.
\]

The exact solution is the exponential boundary-layer profile

\[
u(x) = \frac{e^{P x} - 1}{e^{P} - 1}, \qquad P = \frac{v}{D},
\]

where $P$ is the \engterm{global Peclet number}, the dimensionless ratio of
convective to diffusive transport. For small $P$ the profile is nearly
linear (diffusion dominates); for large $P$ it is flat across most of the
domain and rises steeply in a boundary layer of width $\sim 1/P$ near
$x = 1$ (convection dominates). The quantity of interest is the full profile,
and in particular the solver's behaviour as $P$ grows and the boundary layer
thins below the grid resolution.

\begin{estimation}
\textbf{Question.} Before discretising, estimate the grid spacing $h$ at
which a naive centred-difference scheme will begin to misbehave for
$P = 100$, and predict the qualitative form of the failure.

\smallskip
\textit{Estimate, before reading on.}

\smallskip
The boundary layer has width $\sim 1/P = 0.01$. A scheme that does not
resolve the layer, $h \gtrsim 1/P$, cannot represent the steep rise and will
either smear it or oscillate. The relevant local parameter is the
\engterm{cell Peclet number} $P_{h} = v h / D = P h$, the Peclet number
built on the grid spacing rather than the domain. Dimensional reasoning says
the scheme's character changes when $P_{h}$ crosses order unity: with $P =
100$, that is $h \sim 0.01$, or about $100$ grid intervals. Below that
resolution the centred scheme, which has no inherent bias toward the upwind
direction, is likely to produce spurious oscillations rather than smooth
smearing, because a centred approximation to a strongly one-directional
transport term is the discrete analogue of an unstable algebraic system. The
estimate sets the test: sweep $h$ across $P_{h} = 1$ and watch the centred
scheme break while an upwind scheme stays monotone.
\end{estimation}

\subsection{Two discretisations and the cell Peclet number}

On the uniform grid $x_{i} = i h$, $h = 1/n$, the \engterm{centred} scheme
approximates both derivatives to second order,

\[
v\,\frac{u_{i+1} - u_{i-1}}{2h} - D\,\frac{u_{i-1} - 2 u_{i} + u_{i+1}}{h^{2}}
= 0,
\]

which rearranges to the three-point stencil $(-1 - P_{h}/2)\,u_{i-1} +
2\,u_{i} + (-1 + P_{h}/2)\,u_{i+1} = 0$ with cell Peclet number $P_{h} =
vh/D$. When $P_{h} > 2$ the coefficient of $u_{i+1}$ turns positive, the
tridiagonal system loses diagonal dominance, and the discrete solution
develops node-to-node oscillations that have no counterpart in the smooth
exact profile. The \engterm{upwind} scheme replaces the centred convection
term with a one-sided difference biased into the flow,

\[
v\,\frac{u_{i} - u_{i-1}}{h} - D\,\frac{u_{i-1} - 2 u_{i} + u_{i+1}}{h^{2}}
= 0,
\]

whose stencil $(-1 - P_{h})\,u_{i-1} + (2 + P_{h})\,u_{i} - u_{i+1} = 0$
keeps all off-diagonal signs negative and the matrix diagonally dominant for
every $P_{h}$, so the upwind solution is monotone at any resolution. The
upwind scheme buys that robustness at a price the engineer must name. Its
truncation error is only first order, $O(h)$, and the leading error term is a
spurious diffusion: the upwind difference is equivalent to the centred
difference applied to the equation with diffusion coefficient $D + vh/2$
rather than $D$, an added \engterm{numerical diffusion} of magnitude $vh/2$.
The upwind scheme is stable because it has quietly thickened the boundary
layer, smearing the front the centred scheme oscillated around.

\begin{figure}[ht]
\centering
\begin{tikzpicture}[x=10cm, y=4cm, font=\footnotesize]

% u(x) on [0,1], boundary layer near x=1. P=100 exact: flat then steep rise.
% Axes
\draw[->, thick] (0,0) -- (1.08,0) node[right] {$x$};
\draw[->, thick] (0,0) -- (0,1.15) node[above] {$u$};

\foreach \k in {0,0.25,0.5,0.75,1} {
  \draw (\k,0) -- (\k,-0.02);
  \node[anchor=north, font=\scriptsize] at (\k,-0.02) {$\k$};
}
\foreach \d in {0,0.5,1} {
  \draw (0,\d) -- (-0.012,\d);
  \node[anchor=east, font=\scriptsize] at (-0.012,\d) {$\d$};
}

% Exact: flat near zero until ~0.9, then steep rise to 1 (boundary layer width 0.01)
\draw[thick, black] plot coordinates {
  (0,0) (0.5,0.007) (0.8,0.02) (0.9,0.05) (0.95,0.13) (0.98,0.35)
  (0.99,0.6) (1.0,1.0)
};
\node[black, anchor=south east, font=\scriptsize] at (0.92,0.62) {exact};

% Upwind: monotone but smeared, rises earlier and more gently
\draw[thick, engblue] plot coordinates {
  (0,0) (0.3,0.05) (0.5,0.12) (0.7,0.28) (0.8,0.42) (0.9,0.62)
  (0.95,0.75) (1.0,1.0)
};
\node[engblue, anchor=south east, font=\scriptsize] at (0.7,0.30) {upwind (smeared)};

% Centred: oscillates between nodes near the layer
\draw[thick, engred] plot coordinates {
  (0,0) (0.5,0.01) (0.7,-0.03) (0.75,0.08) (0.8,-0.06) (0.85,0.12)
  (0.9,-0.05) (0.95,0.4) (1.0,1.0)
};
\node[engred, anchor=south, font=\scriptsize] at (0.62,0.13) {centred (oscillatory)};

% zero line for reference (centred undershoots below it)
\draw[dotted, enggray] (0,0) -- (1,0);

\end{tikzpicture}
\caption[Cell-Peclet instability]{The steady convection-diffusion profile for
global Peclet number $P = 100$ on a grid too coarse to resolve the boundary
layer (cell Peclet number $P_{h} = 4 > 2$). The exact solution is flat across
most of the domain and rises steeply in a layer of width $\sim 1/P$ near
$x = 1$. The centred scheme loses diagonal dominance once $P_{h} > 2$ and
oscillates between nodes, overshooting and undershooting; the upwind scheme
stays monotone but smears the front over several cells through its
$vh/2$ numerical diffusion. Both failures are the cell-Peclet constraint in
two forms; the only clean cure is to resolve the layer. Source: authors'
analysis of the standard one-dimensional convection-diffusion test problem.}
\label{fig:vol02:ch16:peclet-instability}
\end{figure}


Figure~\ref{fig:vol02:ch16:peclet-instability} plots both discrete solutions
against the exact profile for $P = 100$ on a grid with $P_{h} = 4$ (too
coarse). The centred solution oscillates between nodes, overshooting and
undershooting the exact profile near the layer; the upwind solution is
monotone but visibly smeared, its front spread over several cells by the
$vh/2$ numerical diffusion. Neither is acceptable at this resolution, and the
two failures are the two faces of the cell-Peclet constraint: a scheme is
either oscillatory or over-diffusive when the layer is unresolved, and the
only clean cure is to resolve it.

\subsection{Verification by a manufactured solution}

The exact boundary-layer profile lets the engineer verify the solver, but
only on this particular right-hand side. To confirm the implementation's
order of accuracy independently of the boundary-layer regime, the disciplined
move is the \engterm{method of manufactured solutions}: choose a smooth
target $u_{\mathrm{m}}(x)$, substitute it into the differential operator to
compute the source term $s(x) = v\,u_{\mathrm{m}}' - D\,u_{\mathrm{m}}''$
that makes $u_{\mathrm{m}}$ the exact solution of the forced problem, then
run the solver with that source and Dirichlet data from $u_{\mathrm{m}}$ and
measure the error against $u_{\mathrm{m}}$. Take $u_{\mathrm{m}}(x) =
\sin(\pi x)$, giving $s(x) = v\pi\cos(\pi x) + D\pi^{2}\sin(\pi x)$, a smooth
forcing with no boundary layer. Running the centred scheme on this
manufactured problem at $P = 1$ (so the cell-Peclet constraint is not
active) and halving $h$ across $n = 10, 20, 40, 80$ should show the error
falling by a factor of four per halving, confirming the $O(h^{2})$ order; the
upwind scheme on the same manufactured problem falls by a factor of two,
confirming its $O(h)$ order. The manufactured solution separates the two
questions the single closed form conflated: \emph{is the code correct}
(does it converge at the design order on a smooth problem) and \emph{is the
discretisation adequate} (does it resolve the boundary layer of the actual
problem). An implementation that passes the manufactured-solution order test
but oscillates on the boundary-layer problem is correct code applied at
inadequate resolution, a diagnosis the engineer must distinguish from a
coding error.

\begin{figure}[ht]
\centering
\begin{tikzpicture}[x=1.4cm, y=0.42cm, font=\footnotesize]

% log-log: x = log10(n), y = log10(max error). n = 10,20,40,80 -> log10 ~ 1..1.9
% Centred slope -2; upwind slope -1.

% Axes
\draw[->, thick] (0.8,-12) -- (3.3,-12) node[right] {$\log_{10} n$};
\draw[->, thick] (0.8,-12) -- (0.8,1) node[above, align=center]
  {$\log_{10}(\text{max error})$};

% x ticks (n = 10, 20, 40, 80, 160)
\foreach \k/\lab in {1/10, 1.301/20, 1.602/40, 1.903/80, 2.204/160} {
  \draw (\k,-12) -- (\k,-12.3);
  \node[anchor=north, font=\scriptsize] at (\k,-12.3) {$\lab$};
}
% y ticks
\foreach \d in {0,-3,-6,-9,-12} {
  \draw (0.8,\d) -- (0.7,\d);
  \node[anchor=east, font=\scriptsize] at (0.7,\d) {$10^{\d}$};
}

% Centred scheme: slope -2. At n=10 (x=1) error 10^{-1.7}; at n=80 (x=1.903) error 10^{-3.5}
% line through (1,-1.7) and (2.204,-4.1) slope = (-4.1+1.7)/(2.204-1)= -2.0
\draw[thick, engblue] (1,-1.7) -- (2.204,-4.1);
\fill[engblue] (1,-1.7) circle (1.6pt);
\fill[engblue] (1.301,-2.3) circle (1.6pt);
\fill[engblue] (1.602,-2.9) circle (1.6pt);
\fill[engblue] (1.903,-3.5) circle (1.6pt);
\node[engblue, font=\scriptsize, anchor=west] at (1.7,-2.7)
  {centred, slope $-2$};

% Upwind scheme: slope -1. through (1,-1.0) and (2.204,-2.2)
\draw[thick, engred] (1,-1.0) -- (2.204,-2.2);
\fill[engred] (1,-1.0) circle (1.6pt);
\fill[engred] (1.301,-1.3) circle (1.6pt);
\fill[engred] (1.602,-1.6) circle (1.6pt);
\fill[engred] (1.903,-1.9) circle (1.6pt);
\node[engred, font=\scriptsize, anchor=west] at (1.7,-1.45)
  {upwind, slope $-1$};

\end{tikzpicture}
\caption[Work-precision diagram for the manufactured solution]{The
work-precision diagram for the convection-diffusion solver on the manufactured
solution $u_{\mathrm{m}}(x) = \sin(\pi x)$ at $P = 1$, on log-log axes. Each
grid halving (rightward by $\log_{10} 2$) drops the centred-scheme error by a
factor of four (slope $-2$) and the upwind-scheme error by a factor of two
(slope $-1$), the machine-checked confirmation of each scheme's order of
accuracy. A line that fails to reach its predicted slope signals an
implementation bug, a boundary-condition error, or an unresolved feature.
Reading the slope off this diagram is the single most useful verification an
engineer runs on a new discretisation. Source: authors' analysis using the
method of manufactured solutions.}
\label{fig:vol02:ch16:mms-convergence}
\end{figure}


Figure~\ref{fig:vol02:ch16:mms-convergence} is the \engterm{work-precision
diagram} for the manufactured problem: the maximum error against the grid
count on log-log axes, one line per scheme. The centred scheme's line has
slope $-2$, the upwind scheme's slope $-1$, and the slopes are the
machine-checked confirmation of the orders the truncation analysis predicted.
A line that fails to reach its predicted slope is the signature of an
implementation bug, a boundary-condition error, or an unresolved feature, and
reading the slope off this diagram is the single most useful verification an
engineer runs on a new discretisation.

\subsection{The budget and the resolved result}

The error budget for the convection-diffusion solve at the target $P = 100$
ranks differently from the cantilever's. Three contributions compete. The
discretisation error is $O(h^{2})$ for the centred scheme but only when
$P_{h} < 2$, which forces $h < 2/P = 0.02$, that is at least $50$ grid
intervals merely to keep the scheme stable and more to make the error small;
at $n = 400$ ($P_{h} = 0.25$) the centred error is near $10^{-4}$ relative.
The round-off contribution is the conditioning of the tridiagonal system
times $u$; the convection-diffusion operator's discrete condition number
grows as $O(n^{2})$ (a second-order operator, milder than the cantilever's
$n^{4}$), reaching $\sim 10^{5}$ at $n = 400$ and a round-off floor near
$10^{-11}$, negligible against the discretisation error. The modelling error
is the physical uncertainty in $v$ and $D$, which for a measured flow and a
tabulated diffusivity might be a few percent each, dominating once the grid
is fine enough. The ranking is the reverse of a naive expectation: the
binding constraint at coarse resolution is stability, the cell-Peclet cap,
rather than accuracy, and once the grid is fine enough to be stable and accurate,
the physical-parameter uncertainty takes over and dominates exactly as it did
for the cantilever.

\begin{figure}[ht]
\centering
\begin{tikzpicture}[x=1.4cm, y=0.42cm, font=\footnotesize]

% log-log: x = log10(n), y = log10(relative magnitude).
% Lines: discretisation error (slope -2, but only valid for n > 50 where P_h<2),
% round-off (slope +2, rising), physical-parameter floor (flat ~ 3e-2).

% Axes
\draw[->, thick] (1.2,-12) -- (4.2,-12) node[right] {$\log_{10} n$};
\draw[->, thick] (1.2,-12) -- (1.2,1) node[above, align=center]
  {$\log_{10}(\text{rel.\ magnitude})$};

% x ticks n = 20,50,100,400,2000
\foreach \k/\lab in {1.301/20, 1.699/50, 2/100, 2.602/400, 3.301/2000} {
  \draw (\k,-12) -- (\k,-12.3);
  \node[anchor=north, font=\scriptsize] at (\k,-12.3) {$\lab$};
}
% y ticks
\foreach \d in {0,-3,-6,-9,-12} {
  \draw (1.2,\d) -- (1.1,\d);
  \node[anchor=east, font=\scriptsize] at (1.1,\d) {$10^{\d}$};
}

% Stability threshold P_h=2 at n=50 (vertical line)
\draw[dashed, enggray] (1.699,-12) -- (1.699,0.5);
\node[enggray, font=\scriptsize, anchor=south, rotate=90] at (1.65,-8)
  {$P_{h}=2$ (stability)};
% Shade unstable region n<50 lightly (just a note)
\node[engred, font=\scriptsize, anchor=south] at (1.45,-1.5)
  {oscillatory};

% Discretisation error: slope -2, valid n>50. through (1.699,-1.5) and (3.301,-4.7)
\draw[thick, engblue] (1.699,-1.5) -- (3.301,-4.7);
\node[engblue, font=\scriptsize, anchor=west] at (2.1,-3.0)
  {discretisation $O(h^{2})$};

% Physical-parameter floor: flat ~ 3e-2 = 10^{-1.5}
\draw[thick, engorange] (1.699,-1.5) -- (4.0,-1.5);
\node[engorange, font=\scriptsize, anchor=south west] at (3.2,-1.5)
  {parameter floor $\sim 3\%$};

% Round-off: slope +2 (kappa ~ n^2), rising. through (1.699,-13) -> (3.301,-9.8)
\draw[thick, engred] (1.699,-13) -- (4.0,-8.4);
\node[engred, font=\scriptsize, anchor=west] at (3.0,-10.5)
  {linear-solve round-off};

% Crossing of discretisation with floor near n ~ 400
\fill[black] (2.602,-1.5) circle (1.8pt);
\node[font=\scriptsize, anchor=south] at (2.602,-1.3)
  {refinement wasted past here};

\end{tikzpicture}
\caption[Error budget for the convection-diffusion solve]{The four competing
magnitudes for the convection-diffusion boundary-value problem at global
Peclet number $P = 100$, against grid count $n$ on log-log axes. Left of the
dashed line ($P_{h} > 2$, fewer than about fifty intervals) the centred scheme
oscillates and no error claim is meaningful. Right of it the discretisation
error falls as $O(h^{2})$ until it drops below the physical-parameter floor
near $n = 400$, past which a finer grid is wasted money against the uncertainty
in $v$ and $D$. The linear-solve round-off rises as the discrete condition
number grows like $n^{2}$ but stays far below the binding terms. The binding
constraint at coarse resolution is stability, not accuracy. Source: authors'
analysis.}
\label{fig:vol02:ch16:conv-diff-budget}
\end{figure}


Figure~\ref{fig:vol02:ch16:conv-diff-budget} stacks the four magnitudes (the
cell-Peclet stability threshold, the discretisation error at the chosen grid,
the linear-solve round-off, the physical-parameter uncertainty) against grid
count, and the crossing points tell the engineer where each constraint binds.
Below $P_{h} = 2$ the solver oscillates and no error claim is meaningful;
above it the discretisation error falls as the grid refines until it drops
below the physical-parameter floor, past which refinement is wasted. The
disciplined deliverable for this problem states the cell-Peclet number of the
chosen grid (the evidence that the scheme is in its stable regime), the
verified order of accuracy from the manufactured-solution study, the
discretisation error at the working resolution, and the physical-parameter
uncertainty that sets the floor. The same four-part statement scales to a
three-dimensional finite-volume flow solver with a hundred million cells,
where the cell-Peclet constraint becomes a local mesh-adaptation criterion,
the manufactured solution becomes a code-verification suite, and the budget
is the artefact that decides whether the next computation should buy a finer
mesh or a better measurement of the inflow.

\begin{principle}[Stability before accuracy]
A discretisation has an order of accuracy and a stability constraint, and the
two are separate properties that fail in separate ways. The order governs how
fast the error falls as the grid refines; the stability constraint governs
whether the discrete problem is well-posed at the chosen resolution at all. A
scheme run outside its stability regime does not merely lose accuracy, it
produces a solution whose error bears no relation to the grid spacing, so the
order of accuracy is meaningless until the stability constraint is satisfied.
The engineer verifies stability first (the cell-Peclet cap for convection-diffusion,
the CFL condition for explicit time stepping, diagonal dominance for an
iterative solve), confirms the order second on a manufactured solution, and
quotes the discretisation error third, against the physical-parameter floor.
The convection-diffusion oscillation and the explicit-integrator blow-up of
section 16.6 are the same failure at different parameters: a discrete operator
pushed past the resolution its stability requires.
\end{principle}

\section{Failure: famous numerical failures}\pathtag{core}

Four named cases ground the chapter's machinery in deployed systems
whose numerical decisions destroyed the result. The Vancouver Stock
Exchange index drift is the gentlest, an accumulating one-sided
rounding bias with no casualties; the Patriot and Ariane cases are
lethal and destructive at the scale of human life and a launch
vehicle; the Pentium FDIV bug carries the lesson into hardware. The
cases share one structure: a numerical decision whose error was
small per operation, well defined, and trivially calculable, and an
operating context that multiplied or exposed the error past the
tolerance the deployment assumed.

\subsection{Vancouver Stock Exchange index drift (1982-1983)}

The Vancouver Stock Exchange introduced a composite index on
25 January 1982 at a nominal value of $1000.000$. Twenty-two months
later the recorded index closed near $525$, having drifted downward
at roughly $25$ points per month while the underlying market was
approximately flat \cite{paper:quinn-vse-1983}. No market event
explained the drift; the cause was an arithmetic rule.

\paragraph{Technical mechanism.}
The index was recomputed on each trade, of which there were roughly
$3000$ per trading day, and the running value was carried to three
decimal places. The update stage truncated the in-memory value to
three decimals rather than rounding it. Truncation of a value
uniformly distributed in its last retained place removes on average
half a unit in that place, a one-sided bias rather than the
zero-mean error of round-to-nearest. With roughly $3000$ updates per
day across about $250$ trading days, the index lost approximately
$3000 \times 250 \times 0.0005 = 375$ index points per year to the
cumulative truncation alone. Over twenty-two months the bias reached
about $574$ points: the recorded index closed at $524.811$ where the
value recomputed with round-half-to-even closed at $1098.892$, a
difference of $574.081$ \cite{paper:quinn-vse-1983}. The correction,
prompted by analyst attention to the persistent drift in a flat
market, recomputed the index retroactively with round-half-to-even
and republished the series \cite{web:cnn-vse-archive}.

\begin{figure}[ht]
\centering
\begin{tikzpicture}[x=0.45cm, y=0.011cm, font=\footnotesize]

% Vancouver Stock Exchange index: recorded value drifting down at
% roughly 25 points/month over 22 months from a nominal 1000, against
% the corrected (round-half-to-even) value that stayed near flat.

\draw[->, thick] (0,0) -- (24,0) node[right] {months since inception};
\draw[->, thick] (0,0) -- (0,1200) node[above, align=center]
  {index value};

\foreach \m in {0,6,12,18,22} {
  \draw (\m,0) -- (\m,-25);
  \node[anchor=north, font=\scriptsize] at (\m,-25) {\m};
}
\foreach \v in {500,1000} {
  \draw (0,\v) -- (-0.4,\v);
  \node[anchor=east, font=\scriptsize] at (-0.4,\v) {\v};
}

% Recorded index: linear downward drift ~25 pts/month from 1000 to ~525
\draw[thick, engred] (0,1000) -- (22,525);
\node[engred, anchor=north west, font=\scriptsize] at (14,640)
  {recorded (truncation)};

% Corrected index: roughly flat market, ends near 1099
\draw[thick, engblue] plot[smooth] coordinates {
  (0,1000) (6,1030) (12,1010) (18,1075) (22,1099)
};
\node[engblue, anchor=south west, font=\scriptsize] at (12,1015)
  {corrected (round-half-to-even)};

% The gap at month 22
\draw[<->, gray] (22,525) -- (22,1099);
\node[gray, anchor=west, font=\scriptsize, align=left] at (22.3,812)
  {$\approx 574$\\ points};

\end{tikzpicture}
\caption[Vancouver Stock Exchange index drift]{The Vancouver Stock
Exchange composite index from inception in January 1982 to the
correction in November 1983. The recorded index (red) drifted downward
at roughly $25$ points per month from its nominal $1000$, closing near
$525$ after twenty-two months, while the underlying market was roughly
flat. The cause was a truncate-to-three-decimals rule applied on each of
roughly $3000$ trades per day; each truncation removed on average half a
unit in the last decimal, a one-sided bias that accumulated to
approximately $574$ points. Recomputing retroactively with
round-half-to-even (blue) recovered a value near $1099$. The case is the
canonical illustration of accumulating one-sided rounding error, the
same mechanism, at a different scale, as the Patriot timing drift.
Source: \cite{paper:quinn-vse-1983}; magnitudes from the case registry.}
\label{fig:vol02:ch16:vse-drift}
\end{figure}


Figure~\ref{fig:vol02:ch16:vse-drift} plots the recorded and the
corrected index. The single-update error was far below the index's
display precision; the cumulative error over millions of updates was
visible at index-level scale. The mechanism is the same one section
16.2 names as accumulated round-off, and the same one the
round-to-nearest-ties-to-even rule of section 16.1 was designed to
defeat: a one-sided rounding rule applied a million times turns a
sub-display-precision per-step error into a structural bias.

\paragraph{Lesson for numerical computation.}
The case is the cleanest demonstration in the chapter of why the
choice of rounding mode is an engineering decision, not an
implementation detail. Truncation (round-toward-zero) is biased;
round-to-nearest-ties-to-even is unbiased, which is exactly why IEEE
754 makes it the default. An engineer who replaces the default with
truncation, for speed or for a misplaced sense that truncation is
the conservative choice, has introduced a one-sided bias that
accumulates with the number of operations. The defence is
institutional as much as numerical: a parallel computation, an
independent audit, or a comparison against a broader-market index
would have caught the drift within months, and none was in place.
The same accumulating-bias mechanism reappears, lethally, in the
Patriot case below, where the accumulated quantity is a clock offset
rather than an index value.

\subsection{Patriot Dhahran (1991)}

A US Army Patriot missile battery at Dhahran, Saudi Arabia, failed
to intercept an incoming Iraqi Scud missile on 25 February 1991.
Twenty-eight US service members died and approximately one
hundred were wounded. The General Accounting Office investigation
identified the cause as a software-clock arithmetic error in the
Patriot's range-gate prediction \cite{acc:gao-patriot-1992}.

\paragraph{Technical mechanism.}
The Patriot's tracking software represented elapsed time as the
number of tenths of a second since the system was started, stored
as an integer counter and converted to seconds as a real number
when needed. The conversion used a 24-bit fixed-point
representation of the constant $1/10$. The exact decimal $0.1$ has
no exact binary representation; in 24-bit fixed point it appears as

\[
0.1_{10} = 0.0001100110011001100110011_{2} + \delta,
\qquad \delta \approx 9.5 \times 10^{-8}.
\]

The truncation error per conversion was about $9.5 \times 10^{-8}$
seconds, accumulating linearly with the integer counter. After
approximately 100 hours of continuous operation, the counter had
reached about $3.6 \times 10^{6}$, and the cumulative conversion
error had reached about $0.34$ seconds. The Patriot's range-gate
prediction used the converted time to forecast the Scud's position
at the next radar look. At the Scud's speed of about $1{,}676
\, \text{m}/\text{s}$, a time offset of $0.34$ seconds
corresponded to a predicted-position error of about $570$ metres,
larger than the range gate. The system reported "no track" and did
not engage \cite{acc:gao-patriot-1992}. The mechanism is
quantitative throughout: the per-tick error, the counter rate, the
cumulative duration, and the target velocity multiply to a single
predicted-position error that the engineer can recompute. The
reconstruction is in
\texttt{code/patriot\_drift.py} and the resulting series is in
\texttt{data/patriot-drift.csv};
Figure~\ref{fig:vol02:ch16:patriot-drift} plots both the timing and
position drift against operating hours.

\begin{figure}[p]
\centering
\begin{tikzpicture}[x=0.85cm, y=0.55cm, font=\footnotesize]

% Panel A: cumulative timing error vs operating hours
\begin{scope}[shift={(0,0)}]
  \node[anchor=south, font=\bfseries] at (5,7.5)
    {Cumulative timing error};
  \draw[->, thick] (0,0) -- (11,0) node[right] {operating hours};
  \draw[->, thick] (0,0) -- (0,7) node[above] {cumulative error (s)};
  % x ticks at 0, 20, 40, 60, 80, 100
  \foreach \h in {0,20,40,60,80,100} {
    \pgfmathsetmacro{\xp}{\h/10}
    \draw (\xp,0) -- (\xp,-0.15);
    \node[anchor=north, font=\scriptsize] at (\xp,-0.15) {\h};
  }
  % y ticks at 0, 0.1, 0.2, 0.3, 0.4
  \foreach \e in {0,1,2,3} {
    \pgfmathsetmacro{\yp}{\e*2}
    \draw (0,\yp) -- (-0.15,\yp);
    \pgfmathsetmacro{\elabel}{\e/10}
    \node[anchor=east, font=\scriptsize] at (-0.15,\yp) {\elabel};
  }
  % Linear accumulation: error = 9.5e-8 * 10 * 3600 * hours = 3.42e-3 * hours
  % at h = 100, error = 0.342 s -> y = 0.342 * 20 = 6.84 (with y-scale s -> 20 units/s)
  \draw[thick, engred] (0,0) -- (10,6.84);
  % Annotate range-gate threshold (e.g. 0.10 s)
  \draw[thick, dashed, gray] (0,2) -- (10,2);
  \node[gray, anchor=west, font=\scriptsize] at (0.2,2.3)
    {range-gate threshold $\approx 0.1$ s};
  % Mark Dhahran failure point: 100 h, 0.34 s
  \fill[engred] (10,6.84) circle (2.5pt);
  \node[engred, anchor=south east, font=\scriptsize, align=right] at (10,6.84)
    {Dhahran:\\$100$ h, $0.34$ s};
\end{scope}

% Panel B: Predicted-position error vs operating hours (Scud at 1676 m/s)
\begin{scope}[shift={(0,-10)}]
  \node[anchor=south, font=\bfseries] at (5,7.5)
    {Predicted-position error at Scud speed};
  \draw[->, thick] (0,0) -- (11,0) node[right] {operating hours};
  \draw[->, thick] (0,0) -- (0,7) node[above] {position error (m)};
  \foreach \h in {0,20,40,60,80,100} {
    \pgfmathsetmacro{\xp}{\h/10}
    \draw (\xp,0) -- (\xp,-0.15);
    \node[anchor=north, font=\scriptsize] at (\xp,-0.15) {\h};
  }
  \foreach \e in {0,200,400,600} {
    \pgfmathsetmacro{\yp}{\e/100}
    \draw (0,\yp) -- (-0.15,\yp);
    \node[anchor=east, font=\scriptsize] at (-0.15,\yp) {\e};
  }
  % Position error = 1676 * 3.42e-3 * hours = 5.73 * hours metres
  % at h = 100, position = 573 m -> y = 5.73 (with 1 unit = 100 m)
  \draw[thick, engred] (0,0) -- (10,5.73);
  % range gate width ~ 137 m (typical for Patriot at intercept geometry; illustrative)
  \draw[thick, dashed, gray] (0,1.37) -- (10,1.37);
  \node[gray, anchor=west, font=\scriptsize] at (0.2,1.6)
    {range-gate half-width $\approx 137$ m (illustrative)};
  \fill[engred] (10,5.73) circle (2.5pt);
  \node[engred, anchor=south east, font=\scriptsize, align=right] at (10,5.73)
    {Dhahran:\\$570$ m};
\end{scope}

\end{tikzpicture}
\caption[Patriot Dhahran cumulative timing and position error]{Patriot
Dhahran cumulative error in two panels. Top: cumulative timing error
from the 24-bit truncation of $0.1$, accumulating linearly at
approximately $9.5 \times 10^{-8} \,\text{s}$ per tenth-of-a-second
tick, reaching approximately $0.34 \,\text{s}$ at the 100-hour mark.
Bottom: the resulting predicted-position error for a Scud-class target
at approximately $1{,}676 \,\text{m/s}$, reaching approximately
$570 \,\text{m}$ at 100 hours, larger than the Patriot's range-gate
half-width. The illustration shows the canonical pattern of
accumulated-truncation failure: a per-operation error that is far
below any single-measurement tolerance becomes mission-critical when
multiplied by an operational duration the original software design
did not bound. The plot is a reconstruction of the timing model in
the GAO report \cite{acc:gao-patriot-1992}; the range-gate half-width
in the lower panel is illustrative.}
\label{fig:vol02:ch16:patriot-drift}
\end{figure}


\paragraph{Organisational context.}
The GAO investigation identified contributing institutional factors
beyond the technical error \cite{acc:gao-patriot-1992}. Patriot
batteries were not designed to operate continuously for one hundred
hours; the operational concept assumed periodic restarts that would
reset the time counter. Field operators in the Gulf War deployment
had been instructed to keep the systems running continuously to
maintain readiness. The accumulated-error effect had been observed
in field tests and reported through the chain, but the patch and
the operational-restart guidance moved through different
distribution channels and arrived at field units inconsistently.
A software patch correcting the conversion was developed and
distributed before the Dhahran incident. The patch arrived at the
Dhahran battery on 26 February 1991, the day after the strike.
The Patriot organisation was not blind to the issue; the field unit
that suffered the strike had received neither the patch nor the
operational-restart guidance that would have substituted for it.
The organisational pattern recurs widely in long-deployment
mission-critical software: a corrected version exists, a workaround
exists, neither one has reached the unit in the field, and the
deployment continues against the original software whose
precondition no longer holds.

\paragraph{Lesson for numerical computation.}
The case is canonical for this chapter because it illustrates a
single principle of section 16.2 in deployed code at fatal scale.
The exact decimal $0.1$ is not representable in binary
floating-point at any finite precision; every conversion incurs a
small representation error; the error accumulates with the
counter. A 24-bit truncation of $0.1$ is mathematically the same
operation as a 53-bit truncation in IEEE 754 binary64; the binary
representation of $0.1$ is non-terminating in any binary base. The
Patriot deployed software whose precondition (operate for short
intervals between restarts) was no longer met by the operational
context (continuous Gulf War standby), and whose accumulated
arithmetic error reached lethal scale before the patch did. Three
specific reading lessons follow. First, accumulated round-off (or
its fixed-point analogue) is the dominant error source in any
long-running computation that performs the same conversion or the
same arithmetic step at a fixed rate; the bound the engineer must
state is the operating-duration bound. Second, the choice of binary
representation for a decimal constant is a per-conversion decision
that the engineer must defend in the precondition; the conversion
$\operatorname{fl}(0.1)$ has a defined error in every binary
floating-point format, and the system designer who treats $0.1$ as
representable has shipped a defect. Third, the chain from the
arithmetic error to the engineering failure is mechanical and
calculable; the engineer who could not have computed the
predicted-position error from the operating hours and the target
velocity is not equipped to deploy the algorithm. The reader who
has run the \texttt{patriot\_drift.py} reconstruction at the
chapter's accompanying code, and who has recovered the
$0.34 \,\text{s}$ timing offset and the $570 \,\text{m}$
position offset as a function of operating hours, has internalised
the case at the depth this chapter is built around.

\subsection{Ariane 5 Flight 501 (1996)}

Ariane 5 Flight 501, the maiden flight of the European Space
Agency's new launcher, broke up about 39 seconds after main
engine ignition on 4 June 1996. The Lions Inquiry Board
identified the cause as an unhandled software exception in the
inertial reference system, triggered when a 64-bit floating-point
representation of an internal horizontal-bias variable was
converted to a 16-bit signed integer and the result exceeded the
integer's range $[-32768, 32767]$ \cite{acc:ariane501-ifr}.

\paragraph{Technical mechanism.}
The conversion ran inside an alignment function reused unchanged
from Ariane 4. For Ariane 4, the trajectory kept the variable
within range; for Ariane 5, the trajectory developed horizontal
velocity faster, the variable did not stay within range, and the
conversion overflowed. The variable in question was an internal
horizontal-bias quantity (named BH in the report) that the Ariane
4 design analysis had certified would never exceed
$\pm 32{,}767$ over the active duration of the alignment function;
that analysis was not re-run for the Ariane 5 flight envelope
\cite{acc:ariane501-ifr}. The exception was unhandled. The primary
inertial reference system shut down. The backup, running the same
software on the same input, shut down for the same reason. The
on-board computer, deprived of valid attitude data, interpreted
diagnostic output from the inertial reference systems as flight
data, commanded extreme nozzle deflections, and the launcher broke
up under aerodynamic loads; the range-safety self-destruct system
then activated. The conversion is the simplest possible numerical
procedure: read a double-precision floating-point value, cast to
a 16-bit signed integer. The error is the simplest possible:
the input lay outside the destination type's range and the
operation was not defined.

\paragraph{Organisational context.}
The Lions Inquiry Board found the failure to be a verification-domain
problem: software qualified for one launcher was reused on another
whose flight envelope had not been re-analysed against the reused
code. The board observed that the Ariane 5 qualification had treated
the reused inertial-reference software as already qualified by
Ariane 4 service history, rather than as new software whose flight
envelope must be re-verified. Seven internal variables in the
alignment function had been analysed for the Ariane 4 envelope;
four had been protected against overflow on grounds that their
ranges could not, in Ariane 4, exceed the destination's range; BH
was one of the four. The protection decision was per-variable and
per-flight-envelope; the rebadge to Ariane 5 did not re-run the
range analysis. The board recommended explicit reuse-with-modification
verification, including range analysis of all internal variables
against the new flight envelope, defensive programming for arithmetic
operations on flight-critical values, and software-architecture
review for reused components in qualification \cite{acc:ariane501-ifr}.
The institutional pattern is recurrent in aerospace: software
inherited from a previous vehicle is treated as a sunk cost rather
than as a component whose envelope must be re-derived. The Ariane
501 inquiry is the canonical case in which the institutional
shortcut produced a loss-of-vehicle failure.

\paragraph{Lesson for numerical computation.}
Volume II Chapter 9 already used this case as the canonical
refusal-of-an-unsafe-numerical-operation example. The reading is the
same here, sharper. The conversion is the simplest possible numerical
procedure, and the precondition is the simplest possible precondition
(the input lies in $[-32768, 32767]$). The procedure ran in a context
whose envelope had not been verified against the precondition. The
discipline that should have refused the operation, and that section
16.2's forward-error apparatus formalises, is the same discipline
section 9.8 articulated: name the precondition, verify the input
meets it, refuse the operation when verification fails. The Patriot
and Ariane cases together teach the two scales of the same lesson.
Patriot is the accumulated-error scale: the per-operation error is
microscopic and acceptable, but the operational duration multiplies
it past tolerance. Ariane is the range-envelope scale: the
per-operation error is well-defined when the input is in range and
catastrophic when it is not, and the operational envelope places
the input out of range. Both deployments lacked the same artefact:
a run-time check whose job was to detect the precondition violation
and to surface it before the algorithm consumed the bad input. The
arithmetic error in each case was undisputed and trivially
calculable; the failure was the absence of the discipline that
section 16.2 names and that the principle below formalises.

\subsection{Pentium FDIV bug (1994), short form}

The Pentium FDIV bug \cite{paper:coe-tang-pratt-1995,paper:edelman-1997}
sits at a third scale, complementing the Patriot and Ariane cases.
The Intel Pentium processor's floating-point division unit
implemented the Sweeney-Robertson-Tocher (SRT) radix-4 algorithm
with a lookup table of pre-computed quotient digits. The table
shipped on the Pentium die was missing five of its 1067 entries; the
empty cells were read as zero by the division hardware. The
probability of triggering an affected entry on a uniformly random
operand pair is approximately one in nine billion; the canonical
test case $4{,}195{,}835.0 / 3{,}145{,}727.0$ returned approximately
$1.33382044$ on an affected Pentium against the correct
$1.33373906$, a relative error near $6 \times 10^{-5}$. The defect
was discovered by Thomas Nicely of Lynchburg College in June 1994
while computing reciprocal sums of twin primes; the public response
through Internet discussion led Intel to announce an unconditional
processor-replacement programme on 20 December 1994, with an
eventual write-off of approximately 475 million US dollars in
1994 dollars.

The Pentium case extends the chapter's reading along two axes.
First, the defect was in hardware, not in software; the lookup
table was burned into silicon, and the remedy required physical
replacement of the processor. The discipline of representing,
checking, and refusing a numerical operation must extend to the
hardware that performs the operation; the consumer cannot patch
a fused circuit. Second, the operand-region the defect affected was
not covered by Intel's pre-release verification; the defect
survived ordinary statistical testing because random operand pairs
hit the affected cells at a rate too low to register against a
sampling-based test budget. The case is the canonical lesson for
formal verification of arithmetic units: random testing cannot
sample the operand space densely enough to find rare
correctness failures, and a unit whose function is to compute a
mathematical operation must be checked against the mathematics, not
against a sample of inputs. Intel's substantial post-1994 investment
in formal verification of arithmetic units, including the use of
theorem provers for the floating-point division and square-root
algorithms in subsequent processor generations, traces directly to
the FDIV episode \cite{paper:edelman-1997}.

\subsection{The principle}

\begin{principle}[Numerical refusal]
A numerical operation has three sources of error: the truncation
of the algorithm, the round-off of the arithmetic, and the
conditioning of the problem. The engineer who deploys the
operation states a precondition (a representation envelope, a
conditioning bound, a smoothness requirement, an operational
duration) that bounds each source of error within the engineering
tolerance. The deployment includes a mechanism to verify the
precondition at run time when the operating envelope can change,
and to refuse the operation, fail safely, or escalate when
verification fails. A reported numerical result without a stated
precondition and a verified bound on each error source is, for
engineering purposes, not a result. The discipline is independent
of the algorithm, of the format, of the language, and of the
domain; the conversion that destroyed Ariane 5 Flight 501 and
the accumulated arithmetic error that the Patriot at Dhahran did
not survive are the same failure mode at different scales.
\end{principle}

\section{Project}\pathtag{core}

\begin{project}[Reproduce a published numerical result]
The reader chooses one published numerical result from a journal
paper, a canonical numerical-textbook exercise, or a documented
case study, reproduces it on a modern floating-point system,
accounts for any drift between the published value and the
reader's, and identifies the dominant numerical error source.
The deliverable is the reproduction artefact plus a 1500-word
reflection.

\textbf{Track}. Simulation with analysis. \textbf{Hazard class}:
none. \textbf{Tools}. Paper, pencil, and a programming
environment with linear-algebra and quadrature primitives:
Python with NumPy and SciPy, MATLAB, Julia, Octave, R, or
equivalent. \textbf{Time}. Approximately one to two hours to
choose and read the source, three to five hours to implement and
run the reproduction, two to three hours to analyse drift, three
to four hours to write the reflection. Total approximately nine
to fourteen hours.

\textbf{Source choice}. The reader chooses one of the following
categories of source.

\begin{itemize}[leftmargin=*]
  \item A chapter exercise from Trefethen and Bau,
    \emph{Numerical Linear Algebra}
    \cite{text:trefethen-bau1997}: the conditioning of a Hilbert
    matrix as a function of order, the orthogonality loss of
    classical versus modified Gram-Schmidt on the same matrix,
    the convergence rate of the QR algorithm on a chosen $4
    \times 4$ matrix.
  \item A chapter exercise from Higham, \emph{Accuracy and
    Stability of Numerical Algorithms}
    \cite{text:higham-numerical}: backward error of LU on a
    pathological matrix, accuracy loss in catastrophic
    cancellation, the variance computation comparison
    (textbook two-pass vs naive one-pass).
  \item A chapter exercise from Moler, \emph{Numerical Computing
    with MATLAB} \cite{text:moler-numerical}: the simulation of a
    chosen ODE with stiffness, a numerical-integration example
    with an integrand singularity, a documented root-finding
    example with a known failure mode.
  \item A specific numerical figure from a published engineering
    paper of the reader's choice. The reader records the citation,
    the relevant equations, and the numerical method described.
\end{itemize}

\textbf{Reproduction protocol}. The reader produces:

\begin{enumerate}[leftmargin=*]
  \item A written statement of the original problem, its
    governing equations, the numerical method described, and the
    published result (a number, a table, a curve).
  \item An implementation of the method in the reader's
    programming environment. The implementation may use library
    primitives (LU, QR, SVD, ODE solvers, quadrature) but the
    reader documents which primitives and which library versions.
  \item The reader's reproduction of the result, presented in
    the same form as the original (table, curve, or scalar).
  \item A drift table: original value, reader's value, absolute
    difference, relative difference. The drift is reported to at
    least three significant figures.
\end{enumerate}

\textbf{Drift accounting}. For each row of the drift table where
the relative drift exceeds $10^{-10}$, the reader identifies the
dominant error source from among:

\begin{itemize}[leftmargin=*]
  \item Different floating-point format (e.g.\ original in single
    precision, reader's in double).
  \item Different algorithm (e.g.\ original used Gauss-Seidel,
    reader's library uses GMRES).
  \item Different ordering or parallelisation of operations.
  \item Different math library (transcendental-function
    differences).
  \item Different convergence tolerance for an iterative method.
  \item Reader's input data differs from the published data
    (e.g.\ a constant rounded differently).
\end{itemize}

\textbf{Reflection (1500 words)}. The reader writes a 1500-word
narrative addressing:

\begin{itemize}[leftmargin=*]
  \item Which numerical decisions in the original work were the
    most consequential for reproducibility, and which were
    incidental.
  \item Whether the published value was reported with sufficient
    precision and provenance to make the reproduction possible
    without ambiguity.
  \item What the reproduction cost in terms of the reader's
    time, software-environment setup, and code-debugging effort.
  \item What an honest reproducibility statement for the original
    paper would have included that was missing.
  \item Whether the reader's reproduction would itself be
    reproducible by a third party from the artefact alone.
\end{itemize}

\textbf{Deliverable}. The source statement, the implementation
code, the reproduced result, the drift table, the
floating-point-environment description (compiler, BLAS, math
library, OS, hardware), and the 1500-word reflection. The
artefact lives in the reader's repository at \repopath{vol02-ch16/project/}.

\textbf{Phases}. The project decomposes into five phases.

\begin{enumerate}[leftmargin=*]
  \item \emph{Source selection and reading} (1-2 hours). Read the
    chosen source closely; isolate the numerical result; list every
    algorithmic and parametric decision the source records and
    every one it does not. The list of \emph{unrecorded} decisions
    is the project's working hypothesis about where the drift will
    come from.
  \item \emph{Implementation} (3-5 hours). Translate the source's
    method into runnable code in the reader's environment. Use
    library primitives where the source used library primitives;
    write from scratch where the source did. Pin compiler and
    library versions.
  \item \emph{Run and tabulate} (1-2 hours). Execute the
    reproduction, log the output, and assemble the drift table.
  \item \emph{Drift diagnosis} (2-3 hours). For each row whose
    drift exceeds $10^{-10}$, attribute the dominant error source
    from the list under \emph{drift accounting} above. The
    attribution is the reader's working hypothesis; the reflection
    treats it as such.
  \item \emph{Reflection writing} (3-4 hours). Draft the 1500-word
    narrative against the five bullet questions; rewrite for
    voice.
\end{enumerate}

\textbf{Success criteria}. The project is complete when:

\begin{itemize}[leftmargin=*]
  \item Every value in the drift table is traceable to a specific
    code path in the reader's implementation;
  \item every drift greater than $10^{-10}$ has a named attribution
    to one of the six listed sources;
  \item the floating-point-environment description is sufficient
    that a third party could reproduce the reader's reproduction on
    different hardware;
  \item the reflection addresses each of the five bullets with at
    least one paragraph of analytical content.
\end{itemize}

\textbf{Common failure modes}. Three recurring failure modes the
reader should anticipate.

\begin{itemize}[leftmargin=*]
  \item The reader cannot reproduce the source value to any precision
    and concludes the source is wrong. The likely cause: a
    parametric difference (different tolerance, different
    convergence criterion, different starting iterate) that the
    source did not record. The remedy: reread the source for
    candidate parameters and tune.
  \item The reader reproduces the source value to twenty digits and
    concludes the project is trivial. The likely cause: the chosen
    source is in a regime where the algorithm is bit-stable across
    platforms (e.g.\ a small symmetric positive definite linear
    system in double precision). The remedy: choose a more
    aggressive source or extend the analysis to a parameter sweep
    that exposes the regime change.
  \item The reflection becomes a description of what was done rather
    than an analysis of what was learned. The remedy: each of the
    five bullets begins with an analytical claim and supports it
    with the project's evidence; description supports the claim,
    not the other way around.
\end{itemize}
\end{project}

\section{Exercises}

The exercises test calculation, derivation, estimation,
simulation, design, diagnosis, and judgment. The editor should
attach the solution and rubric file before release.

\subsection*{Calculation}\pathtag{core}

\begin{exercise}
Compute the unit round-off $u$ and the machine epsilon
$\epsilon_{m}$ for IEEE 754 binary64 and binary32. State the
spacing between $1.0$ and the next representable double-precision
number, and the spacing between $2^{20}$ and the next
representable double-precision number above it.
\end{exercise}

\paragraph{Solution sketch.}
For binary64, $u = 2^{-53} \approx 1.110 \times 10^{-16}$ and
$\epsilon_{m} = 2^{-52} \approx 2.220 \times 10^{-16}$. For binary32,
$u = 2^{-24} \approx 5.960 \times 10^{-8}$ and
$\epsilon_{m} = 2^{-23} \approx 1.192 \times 10^{-7}$. The spacing
between consecutive representables at a normal number $x = 2^{e}$ is
$\epsilon_{m} \cdot 2^{e}$. At $x = 1.0 = 2^{0}$ the spacing is
$\epsilon_{m} \approx 2.220 \times 10^{-16}$. At $x = 2^{20}$ the
spacing is $\epsilon_{m} \cdot 2^{20} \approx 2.328 \times 10^{-10}$,
larger by a factor of $2^{20} \approx 10^{6}$. The lesson is that
absolute spacing scales with the value; relative spacing is
constant.

\begin{exercise}
The number $0.1$ is stored as a binary64 floating-point value.
Compute the absolute representation error $\operatorname{fl}(0.1) - 0.1$ to
three significant figures, and verify that it has the magnitude
predicted by the unit round-off.
\end{exercise}

\paragraph{Solution sketch.}
The exact value of the closest binary64 representable to $0.1$ is
the rational $7205759403792794 / 2^{56}
= 0.1000000000000000055511151231257827021\ldots$. The absolute
representation error is $\operatorname{fl}(0.1) - 0.1
\approx +5.55 \times 10^{-18}$. The unit-round-off prediction is
$\abs{\operatorname{fl}(0.1) - 0.1} \leq 0.1 \cdot u
\approx 1.11 \times 10^{-17}$. The observed error is about half the
bound, consistent with the bound being worst-case under
round-to-nearest. Verification with NumPy (the
\texttt{code/floating\_point\_demo.py} script reconstructs the
exact rational via \texttt{float.as\_integer\_ratio}) confirms the
sign and magnitude.

\begin{exercise}
For the function $f(x) = e^{x} - 1$, evaluate $f$ at $x = 10^{-8}$
in double precision (a) by the direct formula and (b) by the
Taylor approximation $x + x^{2}/2$. Report both values and the
relative error of (a). Identify the source of the error.
\end{exercise}

\paragraph{Solution sketch.}
At $x = 10^{-8}$ the exact value is $e^{x} - 1
\approx 1.0000000050000000 \times 10^{-8}$. The Taylor approximation
(b) returns $x + x^{2}/2 = 1.00000005 \times 10^{-8}$, correct to
all visible digits. The direct evaluation (a) computes
$e^{x} \approx 1.0000000100000001$ rounded to binary64, then
subtracts $1$ to obtain approximately $1.0 \times 10^{-8}$; the
result has only one or two correct digits because the subtraction
cancels the leading 16 identical digits of $e^{x}$ and $1$. The
relative error of (a) is roughly $5 \times 10^{-9}$. The source is
catastrophic cancellation: the well-conditioned function $f$ is
being computed by a numerically unstable algorithm. The C standard
library function \texttt{expm1(x)} returns the stable Taylor form
and is the working choice; the platform-equivalent in MATLAB is
\texttt{expm1}, in NumPy \texttt{numpy.expm1}.

\begin{exercise}
Apply five iterations of bisection to find the root of $f(x) =
x^{3} - x - 2$ in the bracket $[1, 2]$. Tabulate $a_{k}$, $b_{k}$,
$c_{k} = (a_{k} + b_{k})/2$, and $f(c_{k})$ at each step. State
the bracket length after the fifth iteration.
\end{exercise}

\paragraph{Solution sketch.}
$f(1) = -2$, $f(2) = +4$. Step 1: $c_{1} = 1.5$,
$f(c_{1}) = -0.125$; replace $a$, new bracket $[1.5, 2]$.
Step 2: $c_{2} = 1.75$, $f(c_{2}) = +1.609$; replace $b$, new
bracket $[1.5, 1.75]$.
Step 3: $c_{3} = 1.625$, $f(c_{3}) = +0.666$; new bracket
$[1.5, 1.625]$.
Step 4: $c_{4} = 1.5625$, $f(c_{4}) = +0.252$; new bracket
$[1.5, 1.5625]$.
Step 5: $c_{5} = 1.53125$, $f(c_{5}) = +0.0593$; new bracket
$[1.5, 1.53125]$. Bracket length after five steps is
$0.5 / 2^{5} = 1/64 = 0.015625$. The root is approximately
$x^{*} \approx 1.5214$. The error of $c_{5}$ is approximately
$0.010$, consistent with the bound (half the bracket).

\begin{exercise}
Apply five iterations of Newton's method to the same problem, $f
(x) = x^{3} - x - 2$, starting from $x_{0} = 1.5$. Tabulate
$x_{k}$ and $f(x_{k})$, and verify the quadratic convergence by
comparing $\abs{e_{k+1}}$ and $\abs{e_{k}}^{2}$ at each step.
\end{exercise}

\paragraph{Solution sketch.}
$f(x) = x^{3} - x - 2$, $f'(x) = 3x^{2} - 1$. The root, to many
digits, is $x^{*} \approx 1.521379706804568$. Newton:
$x_{0} = 1.5$, $f = -0.125$, $f' = 5.75$, $x_{1} = 1.521739\ldots$,
$e_{1} \approx 3.6 \times 10^{-4}$.
$x_{2} = 1.521379722\ldots$, $e_{2} \approx 1.5 \times 10^{-8}$.
$x_{3} = 1.521379706804568\ldots$, $e_{3} \approx 3 \times 10^{-16}$
(at the binary64 ceiling).
Verification: $e_{1}^{2} \approx 1.3 \times 10^{-7}$ vs
$\abs{e_{2}} \approx 1.5 \times 10^{-8}$; $e_{2}^{2}
\approx 2.3 \times 10^{-16}$ vs $\abs{e_{3}} \approx 3 \times 10^{-16}$.
The ratio $\abs{e_{k+1}} / e_{k}^{2}$ converges to
$C = f''(x^{*}) / (2 f'(x^{*})) = 6 x^{*} / (2 (3 x^{*2} - 1))
\approx 0.6$, matching the quadratic-convergence theory. The
contrast with the bisection exercise is stark: five Newton steps
recover machine precision; bisection needs roughly 50.

\begin{exercise}
Compute the LU factorisation with partial pivoting of $\mat{A} =
\begin{psmallmatrix} 0.001 & 1 \\ 1 & 1 \end{psmallmatrix}$, both
with and without pivoting, and solve $\mat{A}\vect{x} = (1, 2)
^{\transpose}$ in three-decimal-digit floating-point. Compare the
two solutions to the exact solution.
\end{exercise}

\paragraph{Solution sketch.}
The exact solution is $x_{1} = 1/0.999 \approx 1.001$, $x_{2} =
1 - 0.001 x_{1} \approx 0.999$. Without pivoting (pivot $0.001$):
the multiplier is $1/0.001 = 1000$; eliminating gives $U_{22} =
1 - 1000 \cdot 1 = -999$, and the right-hand side updates to
$2 - 1000 \cdot 1 = -998$. In three-digit arithmetic these round
to $-999$ and $-998$, giving $x_{2} \approx 0.999$, then
$x_{1} = (1 - 1 \cdot 0.999) / 0.001 = 0.001 / 0.001 = 1.00$. The
sensitivity is now extreme: in three-digit arithmetic
$1 - 0.999 = 0.001$ is the difference of two nearly equal numbers,
so a single bit error in $x_{2}$ swings $x_{1}$ by $1000$. With
partial pivoting, swap rows so the larger element ($1$) is the
pivot. Multiplier is $0.001/1 = 0.001$, eliminating gives
$U_{22} = 1 - 0.001 \cdot 1 = 0.999$ and right-hand side updates
to $1 - 0.001 \cdot 2 = 0.998$, giving $x_{2} = 0.998/0.999 \approx
0.999$ and $x_{1} = 2 - x_{2} \approx 1.00$. Both solutions are
correct to three digits, but the no-pivot solution is at the mercy
of round-off in $x_{2}$; partial pivoting bounds the multiplier by
$1$ in magnitude and stabilises the back-substitution.

\begin{exercise}
Apply the composite trapezoidal rule with $n = 4, 8, 16, 32$ to
$\int_{0}^{1} e^{x}\, \dd x$. Tabulate the approximation and the
error against the exact value $e - 1$, and verify that the error
behaves as $O(h^{2})$.
\end{exercise}

\paragraph{Solution sketch.}
The exact value is $e - 1 = 1.7182818284\ldots$.
For $n = 4$ ($h = 0.25$): $T_{4} \approx 1.7272219046$, error
$\approx +8.94 \times 10^{-3}$. For $n = 8$ ($h = 0.125$):
$T_{8} \approx 1.7205185\ldots$, error $\approx +2.24 \times
10^{-3}$. For $n = 16$: $T_{16}$, error $\approx +5.59 \times 10^{-4}$.
For $n = 32$: error $\approx +1.40 \times 10^{-4}$. Each halving of
$h$ divides the error by approximately 4, confirming the $O(h^{2})$
prediction. The truncation-error formula
$-(b - a) h^{2} f''(\xi) / 12$ with $f''(x) = e^{x}$ on $[0, 1]$
gives a maximum predicted error of $(1)(0.25)^{2}(e)/12
\approx 1.4 \times 10^{-2}$ for $n = 4$, consistent with the
observed $+0.009$ (the sign is positive because the trapezoidal
rule over-estimates for convex integrands).

\begin{exercise}
Apply Simpson's rule with $n = 4$ to the same integral
$\int_{0}^{1} e^{x}\, \dd x$ and compare the error to the
trapezoidal-rule error from the previous exercise. State the
ratio of the two errors.
\end{exercise}

\paragraph{Solution sketch.}
With $n = 4$, $h = 0.25$: Simpson gives
$S_{4} = (h/3) [f(0) + 4 f(0.25) + 2 f(0.5) + 4 f(0.75) + f(1)]
\approx 1.7183188\ldots$, error $\approx +3.70 \times 10^{-5}$.
The trapezoidal-rule error from the previous exercise was
$+8.94 \times 10^{-3}$. Ratio: trapezoidal/Simpson
$\approx 240$. The truncation-error formulas predict the ratio
$(h^{2}/12 \cdot f''_{\max}) / (h^{4}/180 \cdot f^{(4)}_{\max})
= 15 / h^{2}$ (with $f^{(k)}_{\max} \approx e$ for both); at
$h = 0.25$ this is $15 / 0.0625 = 240$, exactly matching. The
exercise drills the practical reading: doubling the order of the
rule (from 2 to 4) buys two extra decades of accuracy at the
same node count.

\begin{exercise}
Take one step of forward Euler and one step of RK4 for $\dot{y} =
-y$, $y(0) = 1$, with $h = 0.1$. Compute both numerical values of
$y(0.1)$ and compare to the exact value $e^{-0.1}$.
\end{exercise}

\paragraph{Solution sketch.}
Exact value: $e^{-0.1} = 0.90483742\ldots$. Forward Euler:
$y_{1} = y_{0} + h f(t_{0}, y_{0}) = 1 + 0.1 \cdot (-1) = 0.9$.
Error: $-4.84 \times 10^{-3}$.
RK4: $k_{1} = -1$; $k_{2} = -(1 + 0.05 \cdot (-1)) = -0.95$;
$k_{3} = -(1 + 0.05 \cdot (-0.95)) = -0.9525$;
$k_{4} = -(1 + 0.1 \cdot (-0.9525)) = -0.90475$;
$y_{1} = 1 + (0.1/6)(-1 + 2(-0.95) + 2(-0.9525) - 0.90475)
= 0.90483750$. Error: $+8.3 \times 10^{-8}$. RK4 is roughly five
decades more accurate than forward Euler at the same step,
consistent with the $O(h^{5})$ local-truncation prediction vs
forward Euler's $O(h^{2})$ local.

\subsection*{Derivation}\pathtag{standard}

\begin{exercise}
Derive the rounding-error bound $\operatorname{fl}(a + b) = (a + b)(1 + \delta)$
with $\abs{\delta} \leq u$ from the IEEE 754 round-to-nearest
specification, assuming $a + b$ is in the normal range. State
where the assumption fails (overflow, underflow, subnormals).
\end{exercise}

\paragraph{Solution sketch.}
Let $s = a + b$ be the exact mathematical sum. Suppose $s$ lies in
the normal range with $\abs{s} \in [2^{e}, 2^{e+1})$ for some
exponent $e$. The representables near $s$ are spaced by
$\epsilon_{m} \cdot 2^{e} = 2 u \cdot 2^{e}$, so the nearest
representable $\operatorname{fl}(s)$ differs from $s$ by at most
half a spacing: $\abs{\operatorname{fl}(s) - s} \leq u \cdot 2^{e}
\leq u \abs{s}$. Writing $\operatorname{fl}(s) = s (1 + \delta)$
with $\delta = (\operatorname{fl}(s) - s)/s$ gives
$\abs{\delta} \leq u$. Failure cases: overflow when $\abs{s} >
\Omega_{\max}$ (the maximum representable), in which case
$\operatorname{fl}(s) = \pm \infty$ and the relative bound is
meaningless; underflow when $\abs{s} < 2 u \cdot 2^{e_{\min}}$,
where $e_{\min}$ is the minimum normal exponent, in which case
$\operatorname{fl}(s)$ is subnormal and the absolute error is
bounded by $u \cdot 2^{e_{\min}}$ but the relative error can be
arbitrarily large.

\begin{exercise}
Derive the truncation error $-\frac{(b-a)h^{2}}{12} f''(\xi)$ of
the composite trapezoidal rule by Taylor expansion of the
integrand around the midpoint of each subinterval. State the
smoothness assumption on $f$.
\end{exercise}

\paragraph{Solution sketch.}
On a single subinterval $[x_{k}, x_{k+1}]$ of width $h$, expand
$f$ around the midpoint $m = (x_{k} + x_{k+1})/2$:
$f(x) = f(m) + f'(m)(x - m) + \frac{1}{2} f''(m)(x - m)^{2}
+ O((x-m)^{3})$. Integrate over the subinterval: the odd term
vanishes by symmetry; the integral is
$h f(m) + \frac{h^{3}}{24} f''(m) + O(h^{5})$. The trapezoidal
estimate on the same subinterval is $(h/2)(f(x_{k}) + f(x_{k+1}))
= h f(m) + \frac{h^{3}}{8} f''(m) + O(h^{5})$ (Taylor-expanding
the endpoints around $m$). The per-subinterval error is
$T_{k} - I_{k} = h^{3} f''(m)/8 - h^{3} f''(m)/24
= h^{3} f''(m)/12$. Sum over $n = (b - a)/h$ subintervals and
apply the mean-value theorem for sums: the total error is
$(b - a) h^{2} f''(\xi)/12$ for some $\xi \in (a, b)$. The
trapezoidal estimate is larger than the integral when $f'' > 0$
(convex $f$), so the error is conventionally written with a
negative sign: $-(b-a) h^{2} f''(\xi)/12$. Smoothness assumption:
$f \in C^{2}([a, b])$ (twice continuously differentiable).

\begin{exercise}
Derive the quadratic-convergence rate of Newton's method at a
simple root $x^{*}$ from the Taylor expansion $f(x_{k}) = f'(x^{*})
e_{k} + \frac{1}{2} f''(x^{*}) e_{k}^{2} + O(e_{k}^{3})$ and the
update rule. State the constant $C$ in the bound $\abs{e_{k+1}}
\leq C \abs{e_{k}}^{2}$.
\end{exercise}

\paragraph{Solution sketch.}
Write $e_{k} = x_{k} - x^{*}$. By Taylor expansion at $x^{*}$:
$f(x_{k}) = f'(x^{*}) e_{k} + \frac{1}{2} f''(x^{*}) e_{k}^{2}
+ O(e_{k}^{3})$ and $f'(x_{k}) = f'(x^{*}) + f''(x^{*}) e_{k}
+ O(e_{k}^{2})$. The Newton step
$x_{k+1} - x^{*} = x_{k} - x^{*} - f(x_{k})/f'(x_{k})$. Substituting:
$e_{k+1} = e_{k} - \frac{f'(x^{*}) e_{k} + \frac{1}{2} f''(x^{*})
e_{k}^{2} + O(e_{k}^{3})}{f'(x^{*}) + f''(x^{*}) e_{k} + O(e_{k}^{2})}$.
Expand the denominator as a geometric series:
$1/(f'(x^{*}) (1 + (f''(x^{*})/f'(x^{*})) e_{k}))
= (1/f'(x^{*})) (1 - (f''/f') e_{k} + O(e_{k}^{2}))$. After cancellation
of the linear term and collecting $e_{k}^{2}$,
$e_{k+1} = \frac{f''(x^{*})}{2 f'(x^{*})} e_{k}^{2} + O(e_{k}^{3})$.
The constant is $C = \abs{f''(x^{*}) / (2 f'(x^{*}))}$, defined at a
simple root since $f'(x^{*}) \neq 0$.

\begin{exercise}
Show that the Householder reflector $\mat{H} = \mat{I} - 2
\vect{v}\vect{v}^{\transpose}/(\vect{v}^{\transpose}\vect{v})$ is
orthogonal and is its own inverse. Verify that for any nonzero
vector $\vect{x}$, the choice $\vect{v} = \vect{x} + \sgn(x_{1})
\norm{\vect{x}}_{2} \vect{e}_{1}$ produces an $\mat{H}\vect{x}$
that is a multiple of $\vect{e}_{1}$.
\end{exercise}

\paragraph{Solution sketch.}
Let $\beta = 2/(\vect{v}^{\transpose}\vect{v})$ so
$\mat{H} = \mat{I} - \beta \vect{v}\vect{v}^{\transpose}$.
Orthogonality: $\mat{H}^{\transpose}\mat{H} = (\mat{I} - \beta
\vect{v}\vect{v}^{\transpose})^{2} = \mat{I} - 2\beta
\vect{v}\vect{v}^{\transpose} + \beta^{2}
\vect{v}(\vect{v}^{\transpose}\vect{v})\vect{v}^{\transpose}
= \mat{I} - 2\beta\vect{v}\vect{v}^{\transpose} + 2\beta
\vect{v}\vect{v}^{\transpose} = \mat{I}$, using
$\beta(\vect{v}^{\transpose}\vect{v}) = 2$. So $\mat{H}^{\transpose}
= \mat{H}^{-1} = \mat{H}$ (since $\mat{H}$ is symmetric and orthogonal
it is its own inverse). For the second part: with $\sigma = \sgn(x_{1})
\norm{\vect{x}}_{2}$ and $\vect{v} = \vect{x} + \sigma \vect{e}_{1}$:
$\vect{v}^{\transpose}\vect{v} = \norm{\vect{x}}_{2}^{2} + 2\sigma x_{1}
+ \sigma^{2} = 2\sigma(x_{1} + \sigma)$ (using
$\sigma^{2} = \norm{\vect{x}}_{2}^{2}$); $\vect{v}^{\transpose}\vect{x}
= \norm{\vect{x}}_{2}^{2} + \sigma x_{1} = \sigma(x_{1} + \sigma)$.
Then $\mat{H}\vect{x} = \vect{x} - 2(\vect{v}^{\transpose}\vect{x})
\vect{v}/(\vect{v}^{\transpose}\vect{v})
= \vect{x} - \vect{v} = -\sigma \vect{e}_{1}$, a multiple of
$\vect{e}_{1}$. The sign choice $\sgn(x_{1})$ avoids cancellation in
forming $\vect{v}$.

\begin{exercise}
Derive the stability region of forward Euler applied to the
scalar test equation $\dot{y} = \lambda y$ with $\lambda \in \C$.
State the region in the complex plane and verify that the method
is not A-stable. Repeat for backward Euler and verify that it is.
\end{exercise}

\paragraph{Solution sketch.}
Forward Euler on $\dot y = \lambda y$: $y_{n+1} = y_{n} + h \lambda
y_{n} = (1 + h\lambda) y_{n}$. The iteration is stable when
$\abs{1 + h\lambda} \leq 1$, i.e.\ when $h\lambda$ lies in the closed
disk of radius 1 centred at $-1$. The region in the complex
$z = h\lambda$ plane is $\{z : \abs{z + 1} \leq 1\}$. For
$\lambda < 0$ real, this gives $h \leq 2/\abs{\lambda}$. The
region does not contain the entire left half-plane (e.g.\
$z = -2 + 2i$ has $\abs{z + 1} = \sqrt{1 + 4} > 1$), so forward
Euler is not A-stable.
Backward Euler: $y_{n+1} = y_{n} + h\lambda y_{n+1}$, giving
$y_{n+1} = y_{n}/(1 - h\lambda)$. Stable when $\abs{1/(1 - h\lambda)}
\leq 1$, i.e.\ $\abs{1 - h\lambda} \geq 1$. The region is the
complement of the open disk of radius 1 centred at $+1$. This
region contains the entire closed left half-plane (for any
$z$ with $\Re z \leq 0$, $\abs{1 - z} \geq 1$), so backward Euler
is A-stable \cite{text:v2ch16-dahlquist-1956}.

\subsection*{Estimation}\pathtag{standard}

\begin{exercise}
Estimate the maximum number of double-precision floating-point
additions an engineer can perform in a sequential sum before the
worst-case relative round-off error exceeds $10^{-10}$. Compare
to the typical (statistical) bound $\sqrt{n} u$.
\end{exercise}

\paragraph{Solution sketch.}
Worst-case bound: $n u < 10^{-10}$ gives $n < 10^{-10}/u
= 10^{-10}/(1.11 \times 10^{-16}) \approx 9 \times 10^{5}$. So
naive sequential summation of about a million double-precision
numbers risks the worst-case bound. Statistical bound:
$\sqrt{n} u < 10^{-10}$ gives $n < 10^{-20}/u^{2} \approx
8 \times 10^{11}$. The typical attainable count is therefore
nearly a trillion before random-sign rounding accumulation
approaches the $10^{-10}$ bound. The worst case is pessimistic
for benign inputs; pairwise summation
($O(\log n) u$) or Kahan summation
($O(1) u$) defends against the worst case at modest cost.

\begin{exercise}
Estimate the wall-clock time to compute the LU factorisation of a
dense $n \times n$ matrix on a laptop running at $10$ Gflops, for
$n = 10^{3}, 10^{4}, 10^{5}$. State at which $n$ the operation
becomes a project rather than a routine call.
\end{exercise}

\paragraph{Solution sketch.}
LU costs approximately $\frac{2}{3} n^{3}$ flops. At $10^{10}$
flops/s: $n = 10^{3}$ gives $6.7 \times 10^{8}$ flops, time
$\approx 0.07 \,\text{s}$ (routine). $n = 10^{4}$ gives
$6.7 \times 10^{11}$ flops, time $\approx 67 \,\text{s}$
(coffee break). $n = 10^{5}$ gives $6.7 \times 10^{14}$ flops, time
$\approx 19 \,\text{h}$ (project, and the $\approx 80 \,\text{GB}$
of double-precision storage exceeds typical laptop RAM, so the
estimate is academic; iterative or out-of-core methods are
required). The threshold between routine and project on a single
laptop sits near $n = 10^{4}$ for direct dense LU, current as of
2025; a high-core-count workstation extends this by perhaps a
decade.

\begin{exercise}
Estimate the loss of precision (in decimal digits) in solving a
linear system whose condition number is $10^{12}$ in double
precision. Compare to the loss in $80$-bit extended precision and
in $128$-bit quad precision, using the appropriate unit round-offs.
\end{exercise}

\paragraph{Solution sketch.}
The rule of thumb: relative error in the solution is bounded by
$\kappa(\mat{A}) \cdot u$ for a backward-stable algorithm. The
number of correct decimal digits in the solution is approximately
$\log_{10}(1/u) - \log_{10} \kappa$. For binary64 ($u
\approx 1.1 \times 10^{-16}$, $\log_{10}(1/u) \approx 16$): with
$\kappa = 10^{12}$, correct digits $\approx 16 - 12 = 4$. For
80-bit extended ($u \approx 5.4 \times 10^{-20}$, $\log_{10}(1/u)
\approx 19$): correct digits $\approx 19 - 12 = 7$. For binary128
quad ($u \approx 9.6 \times 10^{-35}$, $\log_{10}(1/u) \approx 34$):
correct digits $\approx 34 - 12 = 22$. Format choice buys a few
digits per format step; the conditioning of the problem is fixed
and the algorithm cannot recover digits the problem has lost. The
engineer who reports more digits than the rule-of-thumb count
over-claims.

\begin{exercise}
Estimate the time step a forward-Euler integrator must take to be
stable for a stiff ODE $\dot{y} = -100\,y$, and compare to the
time step required for $\dot{y} = -y$. State the ratio. Estimate
how much wall-clock time is wasted on the fast mode if the slow
mode dominates the solution of interest.
\end{exercise}

\paragraph{Solution sketch.}
Forward-Euler stability on $\dot y = \lambda y$ requires
$h < 2/\abs{\lambda}$. For $\lambda = -100$: $h < 0.02$. For
$\lambda = -1$: $h < 2$. Ratio: $100$. If the slow mode dominates
the solution (say on $t \in [0, 10]$), the natural step for
accuracy is $\sim 0.1$ to resolve the $\lambda = -1$ dynamics,
but the integrator is forced to $h < 0.02$ for stability of the
$\lambda = -100$ mode. The factor of 5 in step count is wasted
compute on the fast mode whose dynamics have already decayed.
An implicit solver (backward Euler or trapezoidal) is A-stable
and can take $h \sim 0.1$, recovering the wasted factor.
At realistic stiffness ratios (e.g.\ chemical kinetics with
$\lambda \sim -10^{6}$ and slow scale $-1$), the wasted factor
reaches $10^{6}$, which is the difference between hours and
seconds.

\begin{exercise}
Estimate the number of correctly rounded decimal digits in the
result of a transcendental computation $\sin(x) - x + x^{3}/6$ at
$x = 0.001$ in double precision, accounting for catastrophic
cancellation. State a numerically stable alternative.
\end{exercise}

\paragraph{Solution sketch.}
The true value is $\sin(x) - x + x^{3}/6 = x^{5}/120 - x^{7}/5040
+ \ldots \approx 8.33 \times 10^{-18}$ at $x = 10^{-3}$. In double
precision, $\sin(x)$ rounds to a value differing from the exact in
the 16th digit; $x$ and $x^{3}/6 \approx 1.67 \times 10^{-10}$ are
represented to full precision. The subtraction $\sin(x) - x$
returns approximately $-x^{3}/6 + O(10^{-19})$; adding $x^{3}/6$
cancels the leading term and leaves a residual on the order of
$10^{-19}$, which is below the round-off floor of $\sin(x)$ at
$\sim 10^{-16}$. The naive answer has approximately
$\log_{10}(8.3 \times 10^{-18}) - \log_{10}(10^{-16}) \approx -2$
correct digits; the answer is dominated by rounding noise. The
stable alternative uses the Taylor identity directly: compute
$x^{5}/120$ (and optionally subtract $x^{7}/5040$ for $\sim 19$
digits). The
\texttt{code/catastrophic\_cancellation.py} script tabulates the
difference across $x = 10^{0}$ to $10^{-12}$ and reproduces the
prediction.

\subsection*{Simulation}\pathtag{standard}

\begin{exercise}
Implement bisection, Newton, and the secant method for $f(x) =
\cos(x) - x$ starting from a sensible initial guess. Plot the
error $\abs{x_{k} - x^{*}}$ on a log scale against the iteration
number. Verify the linear, quadratic, and superlinear rates from
the slopes.
\end{exercise}

\paragraph{Solution sketch.}
The root is $x^{*} = 0.7390851332151607$. Use the bracket
$[0, 1]$ for bisection and the secant method; start Newton at
$x_{0} = 0.5$. Reference implementation in
\texttt{code/convergence\_orders.py}. Expected behaviour:
bisection reduces the error by a factor of 2 per step (linear,
slope $-\log_{10} 2$ on the semilog plot); Newton roughly
squares the error each step (quadratic, slope steepens to
$-2$ on log-log of $e_{k+1}$ vs $e_{k}$); secant has order
$\varphi \approx 1.618$. All three saturate at the binary64 floor
near $10^{-16}$. The chapter's
Figure~\ref{fig:vol02:ch16:convergence-orders} sketches the same
comparison.

\begin{exercise}
Implement the composite trapezoidal rule, the composite Simpson's
rule, and a Gauss-Legendre rule (using the library nodes and
weights) for $\int_{0}^{\pi/2} \sin(x)\, \dd x$. Plot the error
on a log scale against the number of nodes for each method, on
the same axes.
\end{exercise}

\paragraph{Solution sketch.}
The exact value is $\int_{0}^{\pi/2} \sin x\, dx = 1$. Use
NumPy or SciPy: trapezoidal via \texttt{numpy.trapezoid},
Simpson via \texttt{scipy.integrate.simpson}, Gauss-Legendre via
\texttt{numpy.polynomial.legendre.leggauss}. For $n = 4, 8, 16, 32$
nodes the trapezoidal error decays as $O(h^{2})$ (slope $-2$ on
log-log of error vs $n$); Simpson decays as $O(h^{4})$ (slope $-4$);
Gauss-Legendre decays super-algebraically (exponential decay for
this smooth analytic integrand). The Gauss-Legendre rule reaches
the binary64 floor in fewer than 10 nodes. The plot confirms the
theoretical orders by the slopes.

\begin{exercise}
Solve the system $\mat{A}\vect{x} = \vect{b}$ for the Hilbert
matrix $\mat{A}_{ij} = 1/(i + j - 1)$ of orders $n = 5, 8, 10,
12, 14$, with $\vect{b}$ chosen so that the exact solution is
$\vect{x} = (1, 1, \ldots, 1)^{\transpose}$. Solve by LU with
partial pivoting and by SVD with a tolerance cutoff. Tabulate the
relative error against the exact solution for both methods.
\end{exercise}

\paragraph{Solution sketch.}
Reference implementation in \texttt{code/hilbert\_solve.py}.
Expected relative-error progression (LU with partial pivoting,
binary64): $n = 5$ gives $\sim 10^{-12}$; $n = 8$, $\sim 10^{-8}$;
$n = 10$, $\sim 10^{-4}$; $n = 12$, $\sim 10^{-1}$; $n = 14$, of
order unity (the solution has no correct digits). The condition
number grows roughly as $\kappa \approx (\sqrt{2} + 1)^{4n}/\sqrt{n}$,
so $\log_{10}\kappa$ grows by about $1.5$ per unit of $n$;
expected loss-of-digits is $\log_{10}\kappa - 16$. The truncated
SVD with cutoff $\sigma_{1} u n$ produces a regularised solution
whose error stays bounded (at the cost of a biased estimator).
The data table in \texttt{data/hilbert-conditioning.csv} lists the
condition numbers for $n = 3$ to $15$. The exercise drills the
chapter's central point: conditioning is the dominant constraint;
algorithm choice within the well-conditioned regime is secondary,
algorithm choice in the ill-conditioned regime requires
regularisation.

\begin{exercise}
Implement forward Euler, backward Euler, and RK4 for the
nonstiff problem $\dot{y} = -y \sin(t)$, $y(0) = 1$ on $t \in
[0, 10]$. For each method, plot the global error against the
step size on a log-log scale and verify the predicted order of
each method.
\end{exercise}

\paragraph{Solution sketch.}
The exact solution is $y(t) = \exp(\cos t - 1)$. For step sizes
$h = 2^{-k}$ ($k = 0, 1, \ldots, 10$) compute the numerical
solution at $t = 10$ and the absolute error against $y(10)
= \exp(\cos 10 - 1) \approx 0.176$. On log-log axes the slopes
are: forward Euler $\sim 1$ (first order), backward Euler $\sim 1$
(first order), RK4 $\sim 4$ (fourth order). RK4 hits the binary64
floor near $h \approx 10^{-3}$; refining further loses to
round-off as the per-step error budget shrinks below the
arithmetic floor. The plot validates the global-error theory and
exposes the round-off floor in the high-precision regime, mirroring
Figure~\ref{fig:vol02:ch16:error-vs-h}.

\begin{exercise}
Take a stiff problem $\dot{y}_{1} = -1000(y_{1} - \cos t)$,
$\dot{y}_{2} = -y_{2} + y_{1}$, with $y_{1}(0) = y_{2}(0) = 0$,
and integrate it with both an explicit RK4 (taking the largest
$h$ that remains stable) and an implicit BDF or backward-Euler
solver. Compare the wall-clock time, the number of steps, and the
final-value accuracy.
\end{exercise}

\paragraph{Solution sketch.}
Reference implementation in \texttt{code/stiffness\_demo.py}.
RK4 stability requires $h \abs{\lambda_{\max}} \lesssim 2.78$;
with $\lambda \approx -1000$, $h \lesssim 2.5 \times 10^{-3}$.
On $t \in [0, 10]$ this is $\sim 4000$ steps, several milliseconds
of wall-clock time on a current laptop. BDF (\texttt{scipy.
integrate.solve\_ivp} with \texttt{method="BDF"}) takes on the
order of 100 steps and similar wall-clock time but with the
overhead per step of one Newton-style linear solve. The
final-value agreement between the two integrators is at the
tolerance the implicit solver requests ($10^{-6}$ relative by
default). As stiffness ratio increases (replace $-1000$ with
$-10^{6}$), the wall-clock ratio swings dramatically in BDF's
favour; that is the working motivation for the implicit family.

\subsection*{Design}\pathtag{standard}

\begin{exercise}
Design a numerical pipeline for the following problem: given $m$
noisy data points $(x_{i}, y_{i})$ with $m = 10^{4}$, fit a
polynomial of degree $p = 5$ to the data and produce the fitted
coefficient vector with a stated precision. Specify the
floating-point format, the algorithm (normal equations, QR, SVD),
the conditioning check, and the precision the deliverable will
report.
\end{exercise}

\paragraph{Discussion.}
Defensible pipeline. Format: binary64 throughout; binary32 risks
losing digits when the Vandermonde columns span several orders of
magnitude in $x$. Algorithm: form the design matrix $\mat{V} \in
\R^{m \times (p+1)}$ with columns $1, x_{i}, x_{i}^{2}, \ldots,
x_{i}^{5}$, then solve the least-squares problem by Householder QR
(via \texttt{numpy.linalg.lstsq} or \texttt{scipy.linalg.lstsq}).
Avoid the normal equations $\mat{V}^{\transpose}\mat{V}
\vect{\beta} = \mat{V}^{\transpose}\vect{y}$: they square the
condition number ($\kappa(\mat{V}^{\transpose}\mat{V}) =
\kappa(\mat{V})^{2}$). For a degree-5 Vandermonde, $\kappa(\mat{V})$
can reach $10^{5}$ even for benign $x$-grids; squaring it pushes
the normal-equations form near the binary64 limit. Conditioning
check: compute $\kappa(\mat{V})$ via SVD or via
\texttt{numpy.linalg.cond}; if $\kappa > 10^{6}$, consider
centring and scaling $x$ to $[-1, 1]$ or switching to a Chebyshev
basis (Vandermonde-in-Chebyshev), which is dramatically better
conditioned. Reported precision: $\log_{10}(1/u) - \log_{10}\kappa
- 1$ digits, where the $-1$ accounts for the data noise's
contribution to the backward error. For a well-conditioned
Vandermonde, this is $\sim 10$ correct digits; report two or three
fewer than that to leave margin.

\begin{exercise}
Design a reproducibility manifest for a numerical paper of the
reader's choice (or a hypothetical paper described in three
sentences). The manifest specifies compiler version, optimisation
flags, BLAS variant, math-library version, thread count, random
seed, hardware, and operating system. State which of these the
reader's chosen target paper would have left undocumented.
\end{exercise}

\paragraph{Discussion.}
A defensible manifest names each of the following: language and
version (e.g.\ Python 3.12.2 or Julia 1.10.3); compiler / interpreter
version; optimisation flags (\texttt{-O2 -fno-fast-math} for C/C++/
Fortran builds, no automatic vectorisation flags that reorder
floating-point sums); BLAS / LAPACK variant and version (OpenBLAS
0.3.27, MKL 2024.0, Apple Accelerate as bundled with macOS 14.5);
math library (glibc 2.39, Apple Accelerate, Intel libm); thread
count (\texttt{OMP\_NUM\_THREADS=1} for bit-reproducibility);
random seed (with the RNG identified: PCG64, Mersenne Twister,
Philox); hardware (CPU model and microcode revision, GPU model
and driver version where relevant); operating system and kernel
version; and the date of the run. A typical published numerical
paper, current as of 2025, names the language and the library used
and perhaps the operating system; the rest is implicit. The
reader's audit will typically find the BLAS variant, the thread
count, the optimisation flags, and the random seed (in stochastic
work) undocumented. A reproduction artefact that fixes all eight
items by means of a Conda environment or a Nix derivation or a
container is the working standard.

\begin{exercise}
Design a safeguarded root-finder that combines bisection and
Newton's method. Specify the initial bracket, the Newton-step
proposal, the acceptance test (the Newton step lies in the
current bracket), the fallback (bisect when the Newton step is
rejected), and the convergence tolerance. State the worst-case
and the best-case iteration counts.
\end{exercise}

\paragraph{Discussion.}
Reference implementation in \texttt{code/newton\_safeguarded.py}.
The algorithm maintains a bracket $[a, b]$ with $f(a) f(b) < 0$
throughout. At each iteration: compute $f(x_{k})$ and $f'(x_{k})$;
form the Newton proposal $x_{N} = x_{k} - f(x_{k})/f'(x_{k})$; if
$x_{N} \in (a, b)$ and (an optional further sufficient-decrease
check) is satisfied, accept $x_{N}$ as $x_{k+1}$; otherwise fall
back to bisection, $x_{k+1} = (a + b)/2$. Update the bracket
based on the sign of $f(x_{k+1})$. Convergence test:
$\abs{f(x_{k+1})} < \tau_{f}$ or $(b - a) < \tau_{x}$, with
$\tau_{f}, \tau_{x}$ above the round-off floor of the problem
($\sim u \abs{f(x)} \kappa_{f}$, $\sim u \abs{x^{*}}$
respectively). Worst case: pure bisection, $\sim \log_{2}((b-a)/
\tau_{x}) \approx 50$ for a unit bracket and double precision.
Best case: Newton everywhere, $\sim \log_{2}(\log_{2}(1/\tau))
\approx 6$ at quadratic convergence to machine precision. Brent's
method and \texttt{scipy.optimize.brentq} implement a tested
variant of this design.

\subsection*{Diagnosis}\pathtag{standard}

\begin{exercise}
A colleague reports that solving $\mat{A}\vect{x} = \vect{b}$ for
their $1000 \times 1000$ matrix produces results that differ
between MATLAB and NumPy at the fifth decimal digit. List the
candidate explanations and the diagnostic procedure the reader
would run to identify the dominant cause.
\end{exercise}

\paragraph{Solution sketch.}
Candidate explanations. (1) The matrix is ill-conditioned and
$\kappa(\mat{A}) \cdot u \sim 10^{-5}$, so a five-digit cross-
platform disagreement is at the algorithm's stability limit; both
results are correct in the backward sense. (2) The two
implementations call different BLAS variants (MKL vs OpenBLAS vs
Apple Accelerate), and the BLAS reduction order differs, producing
last-bit differences that the conditioning amplifies. (3) The
parallel thread count differs and the parallel reduction order
differs. (4) The two implementations use different algorithms
(LU vs QR vs SVD), and the matrix is rank-deficient enough that
the choice matters. (5) Different default tolerances or different
treatment of small singular values in a regularised solve.
Diagnostic. Compute $\kappa(\mat{A})$ and the residual
$\norm{\mat{A}\vect{x} - \vect{b}}/\norm{\vect{b}}$ on each
platform. If $\kappa \approx 10^{10}$ and both residuals are
small ($\sim 10^{-12}$), the disagreement is expected; both
$\vect{x}$ vectors are backward-stable solutions and the
difference is the conditioning fan-out of last-bit BLAS
reductions. If the residuals differ substantially, one of the
two solves is broken; inspect the algorithm choice and the
tolerance settings. Pin the BLAS and the thread count
(\texttt{OMP\_NUM\_THREADS=1}, \texttt{MKL\_CBWR=COMPATIBLE}) for
the reproducibility audit.

\begin{exercise}
A simulation of a chemical reaction network using
\texttt{scipy.integrate.RK45} reports that the integrator
demands time steps of $10^{-9}$ seconds despite the reaction
proceeding on a $1$-second time scale. Diagnose the issue and
propose a remedy.
\end{exercise}

\paragraph{Solution sketch.}
Stiffness. The reaction network contains at least one fast mode
with $\abs{\lambda_{\max}} \sim 1/(2 \times 10^{-9})
\approx 5 \times 10^{8} \,\text{s}^{-1}$. The slow chemistry of
interest evolves on a $1 \,\text{s}$ scale, but the explicit
\texttt{RK45} solver is forced to a step that resolves the fast
mode for \emph{stability}, not for accuracy. The diagnostic
confirmation: compute the Jacobian $\partial \dot{\vect{y}}/
\partial \vect{y}$ at a representative state and inspect its
eigenvalues; if the spectral radius is several orders of magnitude
larger than $1/T$ where $T$ is the integration horizon, the
system is stiff. Remedy: switch to an implicit, A-stable
integrator. \texttt{scipy.integrate.solve\_ivp(method="BDF")} or
\texttt{method="LSODA"} (which auto-detects stiffness and
switches between Adams and BDF) is the working choice. The
implicit solver pays a Newton-style linear solve per step but
takes steps three to nine decades larger.

\begin{exercise}
A least-squares fit of a degree-10 polynomial to evenly spaced
data on $[0, 100]$ produces coefficients that vary by a factor of
$10^{8}$ between two runs of nominally the same code. Diagnose
the conditioning issue and prescribe two changes to the
formulation that bring the conditioning under control.
\end{exercise}

\paragraph{Solution sketch.}
The Vandermonde matrix on $[0, 100]$ for a degree-10 polynomial
has columns $1, x, x^{2}, \ldots, x^{10}$ ranging in magnitude
from $1$ to $10^{20}$. The columns are far from orthogonal and
the condition number is enormous (typically $\sim 10^{18}$),
which exceeds binary64's representable range. Tiny perturbations
in the data swing the coefficient vector by orders of magnitude;
two runs of the same code that differ only in reduction order or
BLAS variant produce the observed $10^{8}$ swing.
Remedy 1: rescale and re-centre. Map $x$ to $\tilde x =
(x - 50)/50 \in [-1, 1]$ and fit in $\tilde x$. The Vandermonde
columns now span order unity, and $\kappa$ drops by several
orders of magnitude.
Remedy 2: switch to an orthogonal basis. Use the Chebyshev
polynomials of the first kind on $[-1, 1]$ as the basis
(\texttt{numpy.polynomial.Chebyshev.fit}). The Vandermonde-
in-Chebyshev matrix is nearly orthogonal; $\kappa$ for degree 10
is typically $\sim 10$, twenty orders of magnitude better than
the monomial Vandermonde on $[0, 100]$.

\subsection*{Failure analysis}\pathtag{core}

\begin{exercise}
Reread the Patriot Dhahran case study in section 16.8.1. State
which of the three error sources of section 16.2 (truncation,
round-off, conditioning) is the dominant source of the failure,
and which two are not the primary mechanism. State the
verification discipline that, if deployed, would have caught the
error before the deployment.
\end{exercise}

\paragraph{Solution sketch.}
Dominant source: \emph{round-off} (in the sense of representation
error of the constant $0.1$ in finite binary precision), accumulated
over $\sim 3.6 \times 10^{6}$ ticks. The error per conversion is
about $9.5 \times 10^{-8}$, the same per-tick representation residual
in every tick; summing produces a linearly accumulating bias rather
than a random-sign cancellation, so the worst-case bound $nu$ is
tight. Not primary: \emph{truncation error}, since the algorithm
(integer counter + scalar multiplication) has no discretisation
parameter; \emph{conditioning}, since the predicted-position
function is well-conditioned in the time argument
($\kappa \approx 1$ for the linear map). Verification discipline:
state and check the operating-duration precondition. The system's
design assumed periodic restarts; the design analysis should have
included a maximum-continuous-operation duration with a built-in
check (a watchdog or a logged warning) at, say, $50 \%$ of the
duration at which the accumulated drift exceeds the range-gate
tolerance. A run-time accumulator-bound check, with the bound
derived from the per-conversion error and the range-gate width,
would have surfaced the issue before the strike. The
\texttt{code/patriot\_drift.py} script and
\texttt{data/patriot-drift.csv} reconstruction lets the reader
compute the bound directly.

\begin{exercise}
Reread the Ariane 5 Flight 501 case study in section 16.8.2. The
case is sometimes described as a "floating-point error." Defend
or refute that description in three sentences, citing the
specific numerical operation that failed and the precondition
that was violated. Compare to the Patriot case in one further
sentence.
\end{exercise}

\paragraph{Solution sketch.}
The description is misleading. The failed operation was a
type conversion from binary64 floating-point to a 16-bit signed
integer; the precondition for the conversion is that the input
lie in $[-32{,}768, 32{,}767]$, and the violation was that the
horizontal-bias value at the Ariane 5 trajectory exceeded the
upper bound. The arithmetic was correctly rounded throughout; the
failure was the unhandled overflow on a type-conversion operation
whose precondition the new flight envelope no longer satisfied,
not a floating-point rounding error. Compared to the Patriot case
(accumulated representation error over operating duration), Ariane
501 is a single-operation range-envelope failure rather than a
multi-operation accumulated-error failure; both share the deeper
pattern of a precondition that the operating context silently
violated.

\subsection*{Judgment}\pathtag{standard}

\begin{exercise}
A working paper claims a least-squares fit to a $10^{6}$-point
dataset using the normal equations on a Vandermonde matrix of
degree 12, and reports residuals of $10^{-7}$ and coefficients to
six significant figures. State, in three sentences, the question
the reader would put to the paper's authors before accepting the
result. Reference at least one quantity from this chapter.
\end{exercise}

\paragraph{Discussion.}
Question one: what is $\kappa(\mat{V}^{\transpose}\mat{V})$? A
degree-12 Vandermonde on an unrescaled $x$-range can give
$\kappa(\mat{V}) \sim 10^{12}$ to $10^{15}$; the normal equations
square this, putting $\kappa(\mat{V}^{\transpose}\mat{V})$ near or
past binary64's representable range, and the six-figure
coefficients are then below the round-off floor of the algorithm.
Question two: was the data rescaled to $[-1, 1]$ and / or fitted
in a Chebyshev basis? If not, the reported residual of $10^{-7}$
on $10^{6}$ points may reflect a fit dominated by the
high-leverage endpoints rather than a globally accurate model.
Question three: what BLAS / LAPACK call was used, and was the
algorithm normal-equations or backward-stable QR? If normal
equations, the authors are reporting more digits than the
algorithm warrants; switching to \texttt{numpy.linalg.lstsq} or
to Chebyshev fitting is the recommended remedy. (The chapter's
\emph{rule of thumb}: correct digits $\approx \log_{10}(1/u)
- \log_{10} \kappa$.)

\begin{exercise}
An embedded-systems engineer proposes to replace IEEE 754 binary64
arithmetic with binary16 throughout a flight-control inner loop,
citing memory and bandwidth savings. State, in three to five
sentences, what the engineer would have to verify before the
substitution is defensible, and what the engineer would have to
add to the deployment to make the substitution safe in the field.
\end{exercise}

\paragraph{Discussion.}
Verifications before substitution. (1) Range: binary16 represents
roughly $6 \times 10^{-5}$ to $6.5 \times 10^{4}$; the engineer
demonstrates that every intermediate value of the flight-control
inner loop lies in this range across the entire operating
envelope, including all failure scenarios. (2) Precision: binary16
has 11 significant binary digits ($\sim 3$ to $4$ decimal); the
engineer demonstrates that the loop's stability margin and
control accuracy are not eroded by the four-decade loss of
precision relative to binary64. (3) Algorithm conditioning: any
matrix inverse or linear solve in the loop has $\kappa$ small
enough that $\kappa \cdot u_{\text{binary16}} \approx \kappa \cdot
5 \times 10^{-4}$ remains within the control tolerance. (4) Sensor
noise floor: the input quantisation must dominate the round-off
floor for the substitution to be net-positive in any sense beyond
memory savings.
Field-safety additions. Run-time range and overflow checks on
every flight-critical computation, with a defined fall-back
(redundant binary64 path, safe-mode controller, or pilot
notification) when a value approaches the binary16 range
boundary; explicit verification of denormal handling and of the
math-library behaviour on the target hardware; a regression test
suite that exercises the loop across the operating envelope at
the substituted precision; and a formal hazard analysis that
treats every precondition violation as a credible failure mode.

\begin{exercise}
A client requests a "reproducible" Monte Carlo simulation of a
financial portfolio. State, in three sentences, the three
distinct senses of reproducibility that the request might mean,
and the operational discipline corresponding to each. State which
of the three is the working default in the engineer's
deliverable.
\end{exercise}

\paragraph{Discussion.}
Three senses. (1) \emph{Bit-reproducible}: the same code on the
same hardware produces identical output bit-for-bit; achieved by
pinning the random seed, the RNG, the BLAS, the thread count, and
the floating-point flags. (2) \emph{Statistically reproducible}:
re-runs produce summary statistics (mean, variance, value-at-risk
quantile) that agree within a stated tolerance; achieved by
running enough Monte Carlo samples that the Monte Carlo standard
error is below the tolerance, and reporting the tolerance.
(3) \emph{Cross-platform reproducible}: any conformant
implementation, on any hardware, with the same algorithm
specification and the same sampled inputs, produces identical
output; achieved by using bit-reproducible RNGs, bit-reproducible
reductions, and a fully specified algorithm, and is the strictest
and most expensive sense. The working default in a financial
deliverable is (1) for an internal audit trail and (2) for the
reported risk numbers; (3) is reserved for regulator-facing
artefacts where independent re-derivation is required.
